% L'exode rural — Allemand 1ère A — Cours signé OnBuch+
% Programme officiel MINESEC. Compiler avec : tectonic ALA23-l-exode-rural.tex
\newcommand{\DOCMATIERE}{Allemand}
\newcommand{\DOCNIVEAU}{1ère A}
\newcommand{\DOCMODULE}{Chapitre 5 : Citoyenneté}
\newcommand{\DOCLECON}{ALA23}
\newcommand{\DOCTITRE}{L'exode rural}
% OnBuch+ — préambule « cours signés » ENRICHI (compilé avec tectonic / XeLaTeX)
% Système visuel « Pop » d'OnBuch+ : fond crème (jamais blanc), bordures encre
% épaisses, ombres « dures » décalées, titres Archivo Black en capitales.
% Version MATHÉMATIQUES : graphes (pgfplots/tikz), boîtes pédagogiques
% (définition, propriété, méthode, attention, le savais-tu, exercice résolu,
% exercice, TP, bilan), page de garde et signature OnBuch+ ancrée en bas.
%
% Macros attendues dans main.tex AVANT \input{preamble} :
%   \DOCMATIERE  \DOCNIVEAU  \DOCMODULE  \DOCLECON (numéro)  \DOCTITRE
\documentclass[11pt]{article}

\usepackage{fontspec}
\usepackage[ngerman,french]{babel} % français = langue principale ; l'allemand via \all{..} / textebox
\usepackage[a4paper,margin=2cm,top=3.4cm,bottom=2.8cm,headheight=1.4cm,footskip=1.3cm]{geometry}
\usepackage{amsmath,amssymb,amsfonts}
\usepackage{mathtools} % \xmapsto, \coloneqq... souvent utilisés par les modèles
\usepackage{cancel} % simplifications barrées (bilans de forces)
% \abs{z} (module, valeur absolue) et \norm{u} (norme), utilisés par les modèles
\providecommand{\abs}{}\let\abs\relax\DeclarePairedDelimiter{\abs}{\lvert}{\rvert}
\providecommand{\norm}{}\let\norm\relax\DeclarePairedDelimiter{\norm}{\lVert}{\rVert}
\usepackage[table]{xcolor} % [table] : fournit aussi \rowcolors (alternance de lignes)
\usepackage{pagecolor}
\usepackage{tikz}
\usetikzlibrary{arrows.meta,positioning,calc,shapes.geometric,shapes.misc,shapes.arrows,decorations.pathmorphing,decorations.pathreplacing,decorations.markings,patterns,shadows,mindmap,trees,fit,backgrounds,matrix,intersections,angles,quotes}
\usepackage{pgfplots}
\usepackage[version=4]{mhchem} % équations nucléaires \ce{^{235}_{92}U}
\usepackage[european,siunitx]{circuitikz} % schémas électriques
\pgfplotsset{compat=1.18}
\usepgfplotslibrary{fillbetween,groupplots} % aires entre courbes (calcul intégral)
\usepackage{enumitem}
\usepackage{fancyhdr}
% [most] charge toutes les bibliothèques tcolorbox (listings, minted, raster,
% documentation...) et gonfle inutilement la mémoire TeX sur un document long
% (~300 boîtes) : on ne charge que ce qui est réellement utilisé.
\usepackage[skins,breakable]{tcolorbox}
\usepackage{graphicx}
\usepackage{colortbl}
\usepackage{booktabs}
\usepackage{tabularx}
\usepackage{diagbox} % case d'en-tête diagonale des tableaux à double entrée
% Colonnes extensibles centrée / gauche / droite (C, L, R), souvent utilisées par les modèles
\newcolumntype{C}{>{\centering\arraybackslash}X}
\newcolumntype{L}{>{\raggedright\arraybackslash}X}
\newcolumntype{R}{>{\raggedleft\arraybackslash}X}
\usepackage{array}
\usepackage{multicol}
\usepackage{multirow}
\usepackage{siunitx}
% parse-numbers=false : accepte aussi \SI{2,5 \times 10^{-3}}{...}
\sisetup{locale=FR,output-decimal-marker={,},group-minimum-digits=4,parse-numbers=false}
\DeclareSIUnit\ohm{\ensuremath{\symup{\Omega}}} % U+2126 absent de la police
\DeclareSIUnit\division{div} % réglages d'oscilloscope (ms/div, V/div)
\DeclareSIUnit\tr{tr} % tours (vitesse de rotation, ex. \SI{1500}{\tr\per\minute})
\DeclareSIUnit\molar{M} % concentration molaire (ex. \SI{3}{\molar}, \SI{1}{\micro\molar})
\DeclareSIUnit\an{an} % année (le modèle confond parfois avec \year, un compteur TeX) — ex. \SI{5}{\kilo\meter\per\an}
\DeclareSIUnit\FCFA{FCFA} % franc CFA (ex. \SI{500}{\FCFA\per\kilo\gram})
\DeclareSIUnit\degreeBrix{\textdegree Brix} % teneur en sucre d'un sirop/jus (réfractométrie)
\DeclareSIUnit\month{mois} % le modèle utilise parfois \month, un compteur TeX, comme unité de durée
\usepackage{qrcode}
% \degree : commande fréquemment utilisée par les modèles pour un angle ou une
% température en dehors de \SI{}{} (non fournie par les paquets chargés).
\providecommand{\degree}{^{\circ}}
% \qed (fin de démonstration), utilisé par les modèles sans amsthm
\providecommand{\qed}{\hfill$\square$}
% \barre : double barre des valeurs interdites dans un tableau de variations (tkz-tab)
\providecommand{\barre}{\ensuremath{\|}}
% environnement proof (amsthm non chargé) utilisé par les modèles
\ifdefined\proof\else\newenvironment{proof}[1][Démonstration]{\par\noindent\textit{#1.}\ }{\qed\par}\fi
\usepackage{lastpage}
\usepackage{hyperref}
\hypersetup{hidelinks}

% ============================================================
% Palette « Pop » OnBuch+ — mêmes hex que la classe P (plus_ui.dart)
% ============================================================
\definecolor{poporange}{HTML}{F59321}
\definecolor{popdark}{HTML}{E07A0C}
\definecolor{popcream}{HTML}{FFF6E8}
\definecolor{popink}{HTML}{1C1714}
\definecolor{poppurple}{HTML}{7A5AE0}
\definecolor{popgreen}{HTML}{2E9E6A}
\definecolor{poppink}{HTML}{E0457B}
\definecolor{popblue}{HTML}{2F7DE1}
\definecolor{popgold}{HTML}{C98A12}
\definecolor{popmuted}{HTML}{8A7F76}
\definecolor{popline}{HTML}{EADFD2}
\definecolor{popgray}{HTML}{8A7F76}
\colorlet{popgrey}{popgray}
% Alias tolérants (le modèle nomme parfois les couleurs différemment)
\colorlet{popwhite}{white}
\colorlet{popblack}{popink}
\colorlet{popred}{poppink}
\colorlet{popyellow}{popgold}
% Teintes claires (fonds de tableaux/encadrés) — le modèle nomme parfois
% les couleurs avec un suffixe L (ex. popblueL) sans les avoir définies.
\colorlet{poporangeL}{poporange!20}
\colorlet{popdarkL}{popdark!20}
\colorlet{poppurpleL}{poppurple!20}
\colorlet{popgreenL}{popgreen!20}
\colorlet{poppinkL}{poppink!20}
\colorlet{popblueL}{popblue!20}
\colorlet{popgoldL}{popgold!20}
\colorlet{popredL}{popred!20}
\colorlet{popgrayL}{popgray!20}
% Teintes claires pour fonds de tableaux / zones
\colorlet{popblueL}{popblue!10!white}
\colorlet{poporangeL}{poporange!14!white}
\colorlet{popgreenL}{popgreen!12!white}
\colorlet{poppurpleL}{poppurple!12!white}
\colorlet{poppinkL}{poppink!12!white}
\colorlet{popgoldL}{popgold!14!white}

\pagecolor{popcream}

% ============================================================
% Typographie
% ============================================================
\defaultfontfeatures{Ligatures=TeX}
\setmainfont{PlusJakartaSans-Medium.ttf}[Path=../../../fonts/,BoldFont=PlusJakartaSans-Bold.ttf]
% Symboles, API et emojis absents de la police principale : repli sur DejaVu Sans (caractères actifs)
\newfontface\symfont{DejaVuSans.ttf}[Path=../../../fonts/]
\makeatletter
\newcommand\symactive[1]{\catcode#1=13 \begingroup\lccode`\~=#1 \lowercase{\endgroup\def~}{{\symfont\char#1}}}
\newcommand\symblank[1]{\catcode#1=13 \begingroup\lccode`\~=#1 \lowercase{\endgroup\def~}{}}
\newcommand\symhyph[1]{\catcode#1=13 \begingroup\lccode`\~=#1 \lowercase{\endgroup\def~}{\mbox{-}}}
\newcount\symc
\newcommand\symrange[2]{\symc=#1 \@whilenum\symc<#2 \do{\expandafter\symactive\expandafter{\the\symc}\advance\symc by 1 }}
\newcommand\symblankrange[2]{\symc=#1 \@whilenum\symc<#2 \do{\expandafter\symblank\expandafter{\the\symc}\advance\symc by 1 }}
\makeatother
\symrange{"0250}{"02FF}\symactive{"2713}\symactive{"2714}\symactive{"2717}\symactive{"2718}\symactive{"2610}\symactive{"2611}\symactive{"2612}
\symhyph{"2011}\symhyph{"2E1A}\symblankrange{"1F300}{"1FAFF}\symblank{"FE0F}
% Maths en sans-serif (Fira Math) avec chiffres et lettres droites de Plus
% Jakarta, pour que les formules se fondent dans le texte.
\usepackage[math-style=ISO,bold-style=ISO]{unicode-math}
\setmathfont{FiraMath-Regular.otf}[Path=../../../fonts/]
\setmathfont{PlusJakartaSans-Medium.ttf}[Path=../../../fonts/,range={up/num,up/{Latin,latin},\mathup}]
\usepackage{icomma} % virgule décimale sans espace en mode math
% Caractères mathématiques Unicode absents des polices de texte (Plus Jakarta,
% Archivo Black) : rendus par leur équivalent mathématique, en texte comme en
% formule — sinon ils s'affichent en carré vide (ex. « Arithmétique dans ℤ »).
\usepackage{newunicodechar}
\newunicodechar{ℝ}{\ensuremath{\mathbb{R}}}
\newunicodechar{ℤ}{\ensuremath{\mathbb{Z}}}
\newunicodechar{ℂ}{\ensuremath{\mathbb{C}}}
\newunicodechar{ℕ}{\ensuremath{\mathbb{N}}}
\newunicodechar{ℚ}{\ensuremath{\mathbb{Q}}}
\newunicodechar{𝔻}{\ensuremath{\mathbb{D}}}
\newunicodechar{α}{\ensuremath{\alpha}}
\newunicodechar{β}{\ensuremath{\beta}}
\newunicodechar{γ}{\ensuremath{\gamma}}
\newunicodechar{δ}{\ensuremath{\delta}}
\newunicodechar{ε}{\ensuremath{\varepsilon}}
\newunicodechar{θ}{\ensuremath{\theta}}
\newunicodechar{λ}{\ensuremath{\lambda}}
\newunicodechar{μ}{\ensuremath{\mu}}
\newunicodechar{σ}{\ensuremath{\sigma}}
\newunicodechar{τ}{\ensuremath{\tau}}
\newunicodechar{φ}{\ensuremath{\varphi}}
\newunicodechar{ψ}{\ensuremath{\psi}}
\newunicodechar{ω}{\ensuremath{\omega}}
\newunicodechar{Σ}{\ensuremath{\Sigma}}
% majuscules produites par \MakeUppercase dans les titres (α-aminés -> Α)
\newunicodechar{Α}{\ensuremath{\alpha}}
\newunicodechar{Β}{\ensuremath{\beta}}
\newunicodechar{Γ}{\ensuremath{\Gamma}}
\newunicodechar{Δ}{\ensuremath{\Delta}}
\newunicodechar{Ε}{\ensuremath{\varepsilon}}
\newunicodechar{Θ}{\ensuremath{\Theta}}
\newunicodechar{Λ}{\ensuremath{\Lambda}}
\newunicodechar{Μ}{\ensuremath{\mu}}
\newunicodechar{Τ}{\ensuremath{\tau}}
\newunicodechar{Φ}{\ensuremath{\Phi}}
\newunicodechar{Ψ}{\ensuremath{\Psi}}
\newunicodechar{Ω}{\ensuremath{\Omega}}
\newunicodechar{↦}{\ensuremath{\mapsto}}
\newunicodechar{∈}{\ensuremath{\in}}
\newunicodechar{∉}{\ensuremath{\notin}}
\newunicodechar{∧}{\ensuremath{\wedge}}
\newunicodechar{⇔}{\ensuremath{\Leftrightarrow}}
\newunicodechar{⟺}{\ensuremath{\Longleftrightarrow}}
\newunicodechar{⇒}{\ensuremath{\Rightarrow}}
\newunicodechar{⟹}{\ensuremath{\Longrightarrow}}
\newunicodechar{∘}{\ensuremath{\circ}}
\newunicodechar{∪}{\ensuremath{\cup}}
\newunicodechar{∩}{\ensuremath{\cap}}
\newunicodechar{≡}{\ensuremath{\equiv}}
\newunicodechar{⊂}{\ensuremath{\subset}}
\newunicodechar{⊥}{\ensuremath{\perp}}
\newunicodechar{∃}{\ensuremath{\exists}}
\newunicodechar{∀}{\ensuremath{\forall}}
\newunicodechar{∛}{\ensuremath{\sqrt[3]{\,}}}
\newunicodechar{∜}{\ensuremath{\sqrt[4]{\,}}}
\newunicodechar{⃗}{\ensuremath{{}^{\to}}}
\newunicodechar{ⁿ}{\ifmmode^{n}\else\textsuperscript{\ensuremath{n}}\fi}
\newunicodechar{ˣ}{\ifmmode^{x}\else\textsuperscript{\ensuremath{x}}\fi}
\newunicodechar{ᵇ}{\ifmmode^{b}\else\textsuperscript{\ensuremath{b}}\fi}
\newunicodechar{ʳ}{\ifmmode^{r}\else\textsuperscript{\ensuremath{r}}\fi}
\newunicodechar{ᵘ}{\ifmmode^{u}\else\textsuperscript{\ensuremath{u}}\fi}
\newunicodechar{ʸ}{\ifmmode^{y}\else\textsuperscript{\ensuremath{y}}\fi}
\newunicodechar{ᵃ}{\ifmmode^{a}\else\textsuperscript{\ensuremath{a}}\fi}
\newunicodechar{ᵉ}{\ifmmode^{e}\else\textsuperscript{\ensuremath{e}}\fi}
\newunicodechar{ᵏ}{\ifmmode^{k}\else\textsuperscript{\ensuremath{k}}\fi}
\newunicodechar{ᵖ}{\ifmmode^{p}\else\textsuperscript{\ensuremath{p}}\fi}
\newunicodechar{ᴺ}{\ifmmode^{N}\else\textsuperscript{\ensuremath{N}}\fi}
\newunicodechar{ᶜ}{\ifmmode^{c}\else\textsuperscript{\ensuremath{c}}\fi}
\newunicodechar{ʰ}{\ifmmode^{h}\else\textsuperscript{\ensuremath{h}}\fi}
\newunicodechar{ᵠ}{\ifmmode^{q}\else\textsuperscript{\ensuremath{q}}\fi}
\newunicodechar{ₙ}{\ifmmode_{n}\else\textsubscript{\ensuremath{n}}\fi}
\newunicodechar{ₐ}{\ifmmode_{a}\else\textsubscript{\ensuremath{a}}\fi}
\newunicodechar{ᵢ}{\ifmmode_{i}\else\textsubscript{\ensuremath{i}}\fi}
\newunicodechar{ᵣ}{\ifmmode_{r}\else\textsubscript{\ensuremath{r}}\fi}
\newunicodechar{ₖ}{\ifmmode_{k}\else\textsubscript{\ensuremath{k}}\fi}
\newunicodechar{ᵦ}{\ifmmode_{\beta}\else\textsubscript{\ensuremath{\beta}}\fi}
\newunicodechar{ₓ}{\ifmmode_{x}\else\textsubscript{\ensuremath{x}}\fi}
\newfontfamily\popdisplay{ArchivoBlack-Regular.ttf}[Path=../../../fonts/]
\newfontfamily\poplogo{SpaceGrotesk-Bold.ttf}[Path=../../../fonts/]
\newfontfamily\popsemi{PlusJakartaSans-SemiBold.ttf}[Path=../../../fonts/]
\newfontfamily\popxbold{PlusJakartaSans-ExtraBold.ttf}[Path=../../../fonts/]
\newfontfamily\popxboldls{PlusJakartaSans-ExtraBold.ttf}[Path=../../../fonts/,LetterSpace=3.0]

\setlist[enumerate]{leftmargin=1.7em,itemsep=5pt,topsep=4pt}
\setlist[itemize]{leftmargin=1.7em,itemsep=4pt,topsep=4pt,label={\color{poporange}\textbullet}}
\setlength{\parindent}{0pt}
\setlength{\parskip}{5pt}
\renewcommand{\arraystretch}{1.25}

% pgfplots : style Pop par défaut
\pgfplotsset{
  every axis/.append style={axis line style={popink,line width=1pt},tick style={popink},
    grid style={popline,line width=0.5pt},label style={font=\small},tick label style={font=\footnotesize}},
  popcurve/.style={poporange,line width=1.8pt,no markers,smooth},
}
% Même style disponible dans un \draw TikZ simple (hors pgfplots)
\tikzset{popcurve/.style={poporange,line width=1.8pt}}
\tikzset{popgrid/.style={popline,very thin}}
\colorlet{popcurve}{poporange} % le nom du style est parfois utilisé comme couleur

% ============================================================
% Pill / squiggle
% ============================================================
\newcommand{\poppill}[3]{%
  \tikz[baseline=-0.55ex]\node[draw=popink,line width=1.2pt,rounded corners=2.6pt,%
    fill=#2,inner xsep=6.5pt,inner ysep=3pt]{%
    \popxboldls\fontsize{7}{7}\selectfont\color{#3}\MakeUppercase{#1}};%
}
\newcommand{\popsquiggle}[1]{%
  \tikz[baseline=-0.4ex]{
    \draw[#1,line width=2.6pt,line cap=round]
      plot [smooth] coordinates {(0,0.09) (0.34,0.01) (0.68,0.21) (1.02,0.01) (1.36,0.10)};
  }%
}
% Mot-clé surligné (marqueur orange sous le texte)
\newcommand{\cle}[1]{{\popxbold\color{popdark}#1}}

% ============================================================
% En-tête / pied
% ============================================================
\pagestyle{fancy}
\fancyhf{}
\renewcommand{\headrulewidth}{0pt}
\renewcommand{\footrulewidth}{0pt}
\lhead{\raisebox{-0.15ex}{\poplogo\fontsize{13}{13}\selectfont\color{popink}OnBuch\color{poporange}+}}
\rhead{\poppill{\DOCMATIERE\ \textbullet\ \DOCNIVEAU\ \textbullet\ Leçon \DOCLECON}{white}{popink}}
\lfoot{\popxbold\fontsize{7.6}{7.6}\selectfont\color{popmuted}\MakeUppercase{Cours signé OnBuch+}\ \textbullet\ {\popsemi\DOCTITRE}}
\rfoot{\tikz[baseline=-0.6ex]\node[rounded corners=8.5pt,fill=popink,minimum height=17pt,minimum width=17pt,inner xsep=5pt,inner sep=0]{\popxbold\fontsize{8}{8}\selectfont\color{popcream}\thepage\,/\,\pageref*{LastPage}};}

% ============================================================
% Titres
% ============================================================
\newcommand{\chaptitre}[2]{%
  \begin{tcolorbox}[enhanced,colback=poporange,colframe=popink,boxrule=2.6pt,arc=11pt,
      drop shadow=popink,left=15pt,right=15pt,top=12pt,bottom=11pt,before skip=3pt,after skip=18pt]
    \poppill{#2}{popink}{popcream}\par\vspace{9pt}
    {\raggedright\hyphenpenalty=10000\exhyphenpenalty=10000\popdisplay\fontsize{22}{24}\selectfont\color{popcream}\MakeUppercase{#1}\par}
  \end{tcolorbox}
}
% Section numérotée : pastille orange avec le numéro + titre Archivo
\newcounter{popsec}
\newcommand{\coursec}[1]{%
  \stepcounter{popsec}\setcounter{popsub}{0}%
  \par\vspace{16pt}\noindent
  \begin{minipage}{\linewidth}
  \tikz[baseline=-0.7ex]\node[draw=popink,line width=1.6pt,fill=poporange,rounded corners=4pt,minimum size=22pt,inner sep=0]{\popdisplay\fontsize{12}{12}\selectfont\color{popcream}\thepopsec};%
  \hspace{8pt}{\popdisplay\fontsize{15}{17}\selectfont\color{popink}\MakeUppercase{#1}}\par
  \vspace{4pt}\noindent{\color{poporange}\rule{36pt}{3.6pt}}
  \end{minipage}\par\nopagebreak\vspace{8pt}
}
\newcounter{popsub}
\newcommand{\courssub}[1]{%
  \stepcounter{popsub}%
  \par\vspace{9pt}\noindent{\popxbold\fontsize{11.5}{13}\selectfont\color{popink}{\color{poporange}\thepopsec.\thepopsub}\hspace{6pt}#1}\par\nopagebreak\vspace{4pt}
}

% ============================================================
% Boîtes pédagogiques — même recette : fond blanc, bordure épaisse colorée,
% ombre dure encre, badge plein. Titre optionnel [#1].
% ============================================================
\tcbset{popbase/.style={breakable,enhanced,colback=white,boxrule=2.3pt,arc=9pt,
  shadow={3pt}{-3pt}{0pt}{popink},left=14pt,right=14pt,top=11pt,bottom=11pt,before skip=11pt,after skip=15pt,
  fontupper=\popsemi\fontsize{10.6}{14.5}\selectfont\color{popink},
  fontlower=\popsemi\fontsize{10.6}{14.5}\selectfont\color{popink}}}
% Coche dessinée (le glyphe ✓ n'existe pas dans les polices du document)
\newcommand{\popcheck}[1]{\tikz[baseline=-0.5ex]\draw[#1,line width=1.6pt,line cap=round,line join=round](-0.13,0.0)--(-0.04,-0.09)--(0.13,0.1);}
\providecommand{\coche}{\popcheck{popgreen}}
\providecommand{\resultat}[1]{\par\smallskip\noindent\fcolorbox{popgreen}{popgreenL}{\parbox{\dimexpr\linewidth-2\fboxsep-2\fboxrule\relax}{\textbf{Résultat.} #1}}\par\smallskip}
\newcommand{\popboxhead}[3]{\poppill{#1}{#2}{white}\ifx\relax#3\relax\else\hspace{6pt}{\popxbold\fontsize{10.6}{12}\selectfont\color{popink}#3}\fi\par\vspace{7pt}}

\newtcolorbox{aretenir}[1][]{popbase,colframe=popgreen,before upper={\popboxhead{À retenir}{popgreen}{#1}}}
% Texte en allemand (document de lecture, dialogue, extrait) : cadre
% d'étude. Un titre optionnel (source, auteur) s'affiche dans l'en-tête.
\newtcolorbox{textebox}[1][]{popbase,colframe=popblue,before upper={\popboxhead{Text}{popblue}{#1}\begin{otherlanguage}{ngerman}},after upper={\end{otherlanguage}}}
% Marques ✓ / ✗ (forme correcte / fautive) souvent utilisées par les modèles
\providecommand{\cross}{\textcolor{poppink}{\ensuremath{\times}}}
\providecommand{\xmark}{\cross}\providecommand{\crossmark}{\cross}\providecommand{\cmark}{\ensuremath{\checkmark}}
% Ligne de réponse / justification à compléter (inventée par les modèles)
\providecommand{\justification}[1]{\par\noindent\textit{Justification~:} \dotfill\par}
% \textipa (package tipa non chargé) : la police ne couvre pas l'API -> simple passage
\providecommand{\textipa}[1]{#1}
% Soulignement (ulem non chargé) : \uline, \ul
\providecommand{\uline}[1]{\underline{#1}}\providecommand{\ul}[1]{\underline{#1}}
% \glossaire{\item ...} inventée par les modèles : liste de glossaire
\providecommand{\glossaire}[1]{\begin{itemize}#1\end{itemize}}
% \alert (beamer) inventée par les modèles : texte mis en évidence
\providecommand{\alert}[1]{\textbf{\textcolor{poppink}{#1}}}
% variantes ASCII d'ordinaux écrites en macro
\providecommand{\iere}{\textsuperscript{re}}\providecommand{\ieme}{\textsuperscript{e}}\providecommand{\ere}{\textsuperscript{re}}\providecommand{\eme}{\textsuperscript{e}}\providecommand{\er}{\textsuperscript{er}}\providecommand{\ie}{\textsuperscript{e}}
% Ordinaux français écrits en macro par les modèles : 1\ière, 3\ième
\newcommand{\ière}{\textsuperscript{re}}\newcommand{\ième}{\textsuperscript{e}}
% \exemple (inventée par les modèles) : étiquette « Exemple : »
\providecommand{\exemple}{\textbf{Exemple~:}\ }
% \square utilisable aussi hors mode math (cases à cocher)
\AtBeginDocument{\let\squareMath\square\renewcommand{\square}{\ensuremath{\squareMath}}}
% Mot ou phrase en allemand à l'intérieur d'un paragraphe français
\newcommand{\all}[1]{\ifmmode\text{\textit{#1}}\else\begingroup\selectlanguage{ngerman}\itshape #1\endgroup\fi}
\let\esp\all % alias toléré
\newtcolorbox{exemplebox}[1][]{popbase,colframe=poporange,before upper={\popboxhead{Exemple}{poporange}{#1}}}
\newtcolorbox{definition}[1][]{popbase,colframe=popblue,before upper={\popboxhead{Définition}{popblue}{#1}}}
\newtcolorbox{propriete}[1][]{popbase,colframe=poppurple,before upper={\popboxhead{Propriété}{poppurple}{#1}}}
\newtcolorbox{methode}[1][]{popbase,colframe=popgold,before upper={\popboxhead{Méthode}{popgold}{#1}}}
\newtcolorbox{attention}[1][]{popbase,colframe=poppink,before upper={\popboxhead{Attention}{poppink}{#1}}}
\newtcolorbox{experience}[1][]{popbase,colframe=popblue,colback=popblueL,before upper={\popboxhead{Expérience}{popblue}{#1}}}
% « Le savais-tu ? » : carte violette pleine (contexte, culture, Cameroun)
\newtcolorbox{savaistu}[1][]{popbase,colback=poppurple,colframe=popink,
  fontupper=\popsemi\fontsize{10.4}{14.2}\selectfont\color{white},
  before upper={\poppill{Le savais-tu\,?}{white}{poppurple}\ifx\relax#1\relax\else\hspace{6pt}{\popxbold\color{white}#1}\fi\par\vspace{7pt}}}
% Exercice résolu : énoncé en haut, \tcblower puis la solution sur fond crème
\newcounter{popexo}\newcounter{popexr}
\newtcolorbox{exoresolu}[1][]{popbase,colframe=poporange,colbacklower=popcream,
  segmentation style={popink,line width=1.2pt,dashed},
  before upper={\stepcounter{popexr}\popboxhead{Exercice résolu \thepopexr}{poporange}{#1}},
  before lower={\poppill{Solution}{popink}{popcream}\par\vspace{6pt}}}
% Exercice (non résolu en place) — la correction est dans la partie Corrigés
\newtcolorbox{exercice}[1][]{popbase,colframe=popink,
  before upper={\stepcounter{popexo}\popboxhead{Exercice \thepopexo}{popink}{#1}}}
% Corrigé
\newtcolorbox{corrige}[1][]{popbase,colframe=popgreen,colback=popgreenL,
  before upper={\popboxhead{Corrigé}{popgreen}{#1}}}
% Objectifs / prérequis / bilan
\newtcolorbox{objectifs}{popbase,colframe=popink,colback=poporangeL,
  before upper={\popboxhead{Objectifs de la leçon}{popink}{}}}
\newtcolorbox{prerequis}{popbase,colframe=popink,colback=popblueL,
  before upper={\popboxhead{Prérequis}{popblue}{}}}
\newtcolorbox{bilan}{popbase,colframe=popink,colback=white,boxrule=2.8pt,
  before upper={\popboxhead{Fiche bilan}{poporange}{L'essentiel à connaître pour l'examen}}}
% Figure encadrée avec légende
\newtcolorbox{popfigure}[1][]{enhanced,breakable=false,colback=white,colframe=popline,boxrule=1.4pt,arc=8pt,
  left=10pt,right=10pt,top=10pt,bottom=8pt,before skip=10pt,after skip=12pt,halign=center,
  lower separated=false,halign lower=center,
  fontlower=\popsemi\fontsize{9}{11.5}\selectfont\color{popmuted},
  #1}
\newcommand{\legende}[1]{\par\vspace{6pt}{\popsemi\fontsize{9}{11.5}\selectfont\color{popmuted}#1}}

% ============================================================
% Signature OnBuch+
% ============================================================
\newcommand{\poplogoinline}[1]{{\poplogo\fontsize{#1}{#1}\selectfont\color{popink}OnBuch\color{poporange}+}}
\newcommand{\signatureonbuch}{%
  \begin{tcolorbox}[enhanced,colback=popink,colframe=popink,boxrule=2.2pt,arc=10pt,
      drop shadow=poporange,left=14pt,right=14pt,top=10pt,bottom=10pt,before skip=0pt,after skip=0pt]
    \tikz[baseline=-0.7ex]\node[circle,fill=popgreen,draw=popcream,line width=1.3pt,minimum size=22pt,inner sep=0]{\popcheck{white}};%
    \hspace{9pt}\begin{minipage}[c]{0.62\linewidth}
      {\popxbold\fontsize{12}{13}\selectfont\color{popcream}Signé\ \textbullet\ {\poplogo OnBuch\color{poporange}+}}\par\vspace{2pt}
      {\raggedright\popsemi\fontsize{8}{10.5}\selectfont\color{popline}\MakeUppercase{Équipe pédagogique OnBuch+}\\\MakeUppercase{Conforme au programme officiel MINESEC}\par}
    \end{minipage}\hfill
    {\poplogo\fontsize{22}{22}\selectfont\color{popcream}OnBuch\color{poporange}+}
  \end{tcolorbox}%
}

% ============================================================
% Page de garde
% #1 titre · #2 matière · #3 niveau · #4 module · #5 n° leçon · #6 durée ·
% #7 description / objectifs (texte ou itemize)
% ============================================================
\newcommand{\pagegarde}[7]{%
  \thispagestyle{empty}
  \newgeometry{a4paper,margin=1.7cm,top=1.6cm,bottom=1.4cm}%
  \begin{tikzpicture}[remember picture,overlay]
    \fill[poporange!18] ([xshift=-3.2cm,yshift=-3.6cm]current page.north east) circle (4.6cm);
    \fill[poppurple!14] ([xshift=2.4cm,yshift=7.6cm]current page.south west) circle (3.8cm);
  \end{tikzpicture}%
  {\poplogo\fontsize{30}{30}\selectfont\color{popink}OnBuch\color{poporange}+}\hfill
  \raisebox{0.6ex}{\poppill{Cours signé \textbullet\ Premium}{popink}{popcream}}\par
  \vspace{1pt}\popsquiggle{poppurple}\par
  \vspace{8pt}
  \begin{tcolorbox}[enhanced,colback=poporange,colframe=popink,boxrule=2.8pt,arc=13pt,
      drop shadow=popink,left=18pt,right=18pt,top=12pt,bottom=13pt,after skip=10pt]
    \poppill{#2 \textbullet\ #3}{popink}{popcream}\hspace{5pt}\poppill{#4}{white}{popink}\par\vspace{8pt}
    {\popdisplay\fontsize{14}{14}\selectfont\color{popink}LEÇON #5}\par\vspace{4pt}
    {\raggedright\hyphenpenalty=10000\exhyphenpenalty=10000\popdisplay\fontsize{25}{27}\selectfont\color{popcream}\MakeUppercase{#1}\par}
    \vspace{8pt}
    {\popxbold\fontsize{9.5}{9.5}\selectfont\color{popink}\MakeUppercase{Durée conseillée : #6}}
  \end{tcolorbox}
  \begin{tcolorbox}[enhanced,colback=white,colframe=popink,boxrule=2.2pt,arc=10pt,
      drop shadow=popink,left=16pt,right=16pt,top=10pt,bottom=10pt,after skip=8pt]
    \poppill{Dans cette leçon}{poppurple}{white}\par\vspace{5pt}
    \popsemi\fontsize{9.3}{12.6}\selectfont\color{popink}#7
  \end{tcolorbox}
  \noindent\begin{minipage}[t]{0.487\linewidth}
    \begin{tcolorbox}[enhanced,colback=popgreen,colframe=popink,boxrule=2.2pt,arc=10pt,
        drop shadow=popink,left=11pt,right=11pt,top=10pt,bottom=10pt,height=3.1cm,valign=center]
      \tcbox[colback=white,colframe=popink,boxrule=1.3pt,arc=5pt,boxsep=0pt,nobeforeafter,
        left=3pt,right=3pt,top=3pt,bottom=3pt,valign=center]{\qrcode[height=1.9cm]{https://play.google.com/store/apps/details?id=com.luvvix.onbuch&hl=fr}}%
      \hspace{8pt}\begin{minipage}[c]{0.5\linewidth}
        {\popdisplay\fontsize{11}{12.5}\selectfont\color{white}\MakeUppercase{Télécharge OnBuch}}\par\vspace{4pt}
        {\raggedright\popsemi\fontsize{8.4}{11}\selectfont\color{white}Gratuit \textbullet\ Google Play\\Tuteur IA, annales, concours, résultats\par}
      \end{minipage}
    \end{tcolorbox}
  \end{minipage}\hfill
  \begin{minipage}[t]{0.487\linewidth}
    \begin{tcolorbox}[enhanced,colback=poppurple,colframe=popink,boxrule=2.2pt,arc=10pt,
        drop shadow=popink,left=11pt,right=11pt,top=10pt,bottom=10pt,height=3.1cm,valign=center]
      \tikz[baseline=-0.7ex]\node[circle,fill=white,draw=popink,line width=1.3pt,minimum size=1.6cm,inner sep=0]{\popdisplay\fontsize{24}{24}\selectfont\color{poppurple}+};%
      \hspace{8pt}\begin{minipage}[c]{0.6\linewidth}
        {\popdisplay\fontsize{11}{12.5}\selectfont\color{white}\MakeUppercase{Passe à OnBuch+}}\par\vspace{4pt}
        {\raggedright\popsemi\fontsize{8.4}{11}\selectfont\color{white}Tous les cours signés, Tuteur IA illimité, fiches PDF — onglet OnBuch+ de l'app\par}
      \end{minipage}
    \end{tcolorbox}
  \end{minipage}
  \vspace{8pt}
  \popcontenu
  \vfill
  \signatureonbuch
  \restoregeometry
  \newpage
}

% Rangée « Ce cours contient » (page de garde)
\newcommand{\poptile}[3]{% #1 couleur #2 gros texte #3 légende
  \begin{tcolorbox}[enhanced,colback=#1,colframe=popink,boxrule=2pt,arc=9pt,shadow={2.5pt}{-2.5pt}{0pt}{popink},
      left=6pt,right=6pt,top=9pt,bottom=9pt,halign=center,valign=center,height=1.75cm,width=\linewidth,nobeforeafter]
    {\popdisplay\fontsize{12.5}{13}\selectfont\color{white}\MakeUppercase{#2}}\par\vspace{3pt}
    {\centering\popsemi\fontsize{7.8}{9.5}\selectfont\color{white}#3\par}
  \end{tcolorbox}}
\newcommand{\popcontenu}{%
  \par\noindent{\popxboldls\fontsize{7.5}{7.5}\selectfont\color{popmuted}CE COURS SIGNÉ CONTIENT}\par\vspace{7pt}
  \noindent\begin{minipage}[t]{0.235\linewidth}\poptile{popblue}{Cours}{complet, illustré, pas à pas}\end{minipage}\hfill
  \begin{minipage}[t]{0.235\linewidth}\poptile{poporange}{Méthodes}{et exercices résolus}\end{minipage}\hfill
  \begin{minipage}[t]{0.235\linewidth}\poptile{poppink}{Exercices}{type examen + corrigés}\end{minipage}\hfill
  \begin{minipage}[t]{0.235\linewidth}\poptile{popgreen}{Fiche bilan}{l'essentiel à retenir}\end{minipage}\par}

% Sommaire
\newtcolorbox{sommaire}{enhanced,colback=white,colframe=popink,boxrule=2.2pt,arc=10pt,
  drop shadow=popink,left=18pt,right=18pt,top=13pt,bottom=13pt,before skip=6pt,after skip=20pt,
  before upper={\poppill{Sommaire}{poppurple}{white}\par\vspace{9pt}\popsemi\fontsize{10.6}{14.5}\selectfont\color{popink}}}

% Signature de fin de cours — poussée tout en bas de la dernière page
\newcommand{\findecours}{%
  \par\vfill\nopagebreak
  \begin{center}\popsquiggle{poporange}\end{center}\vspace{4pt}
  \signatureonbuch
}
\AtBeginDocument{\def\llbracket{\mathopen{[\mkern-3.5mu[}}\def\rrbracket{\mathclose{]\mkern-3.5mu]}}} % intervalles d'entiers (glyphe ⟦ absent de la police)
% Glyphes absents de Fira Math : redéfinis à partir de glyphes disponibles
\AtBeginDocument{%
  \def\setminus{\mathbin{\backslash}}\let\smallsetminus\setminus
  \def\vdots{\mathord{\vbox{\baselineskip3.2pt\lineskiplimit0pt\kern4pt\hbox{.}\hbox{.}\hbox{.}}}}%
  \def\ddots{\mathinner{\mkern1mu\raise6.5pt\hbox{.}\mkern2mu\raise3.6pt\hbox{.}\mkern2mu\raise0.7pt\hbox{.}\mkern1mu}}%
  \def\triangle{\mathord{\scalebox{0.9}{$\symup{\Delta}$}}}%
  \def\nabla{\mathord{\raisebox{\depth}{\scalebox{1}[-1]{$\symup{\Delta}$}}}}%
  \def\checkmark{\popcheck{popdark}}%
  \def\star{\mathbin{\ast}}\def\bigstar{\mathord{\ast}}%
  \def\diamond{\mathbin{\scalebox{0.6}{$\Diamond$}}}%
  \def\bigcup{\mathop{\mathchoice{\raisebox{-0.3em}{\scalebox{1.7}{$\cup$}}}{\scalebox{1.25}{$\cup$}}{\cup}{\cup}}\displaylimits}%
  \def\bigcap{\mathop{\mathchoice{\raisebox{-0.3em}{\scalebox{1.7}{$\cap$}}}{\scalebox{1.25}{$\cap$}}{\cap}{\cap}}\displaylimits}%
  \def\lesseqqgtr{\mathrel{\lessgtr}}\def\bigoplus{\mathop{\oplus}\displaylimits}%
}
\providecommand{\ssi}{si et seulement si} % abréviation fréquente
\AtBeginDocument{\providecommand{\wideparen}{\overparen}} % arc de cercle

%% FIGFIX-BEGIN v3 — correctifs globaux des figures (docs/onbuch-generation/tools/figfix)
\usepackage{adjustbox}\usepackage{etoolbox}
% toute figure TikZ/pgfplots plus large que la ligne est réduite à la largeur disponible
\BeforeBeginEnvironment{tikzpicture}{\begin{adjustbox}{max width=\linewidth}}
\AfterEndEnvironment{tikzpicture}{\end{adjustbox}}
% légendes pgfplots lisibles par-dessus les courbes
\pgfplotsset{every axis/.append style={legend style={fill=white,fill opacity=0.92,text opacity=1,draw=black!15,inner sep=3pt}}}
% étiquettes d'arêtes (syntaxe edge["..."]) avec fond blanc
\tikzset{every edge quotes/.append style={fill=white,inner sep=1.2pt}}
% bulles de cartes mentales : texte à la ligne dans la bulle, bulle un peu plus grande
\tikzset{every concept/.append style={text width=2.0cm,align=center,minimum size=2.3cm,inner sep=2pt}}
%% FIGFIX-END
\begin{document}

\pagegarde{L'exode rural}{Allemand}{1ère A}{Chapitre 5 : Citoyenneté}{ALA23}{2 h}{Lecture et commentaire d'un texte allemand sur l'exode rural (die Landflucht), enrichissement du vocabulaire thématique du village et de la ville, et étude grammaticale de l'expression de la cause et de la conséquence, notamment les subordonnées en weil.
\vspace{4pt}
\begin{multicols}{2}\small
\begin{itemize}
\item À la fin de cette leçon, je sais comprendre globalement puis en détail un texte allemand sur l'exode rural
\item je sais définir et employer le vocabulaire thématique : das Dorf, die Stadt, die Landflucht, die Arbeit, umziehen, die Armut
\item je sais donner l'article et le pluriel des noms du thème et les classer par familles de mots
\item je sais exprimer la cause avec weil, denn et wegen en plaçant correctement le verbe
\item je sais exprimer la conséquence avec deshalb, darum, so dass
\item je sais répondre par écrit et à l'oral à des questions de compréhension en phrases complètes
\end{itemize}\end{multicols}}

\begin{sommaire}
\begin{multicols}{2}
\begin{enumerate}
\item Introduction et découverte du thème
\item Texte support : « Warum verlassen die jungen Leute das Dorf ? »
\item Compréhension du texte
\item Lexique thématique : Dorf, Stadt, Landflucht
\item Grammaire : la cause et la conséquence
\item Méthode du commentaire et commentaire modèle
\item Production guidée
\item Activité d'intégration
\item Méthodes et exercices résolus
\item Exercices
\item Corrigés des exercices
\item Fiche bilan
\end{enumerate}
\end{multicols}
\end{sommaire}

\begin{objectifs}
\begin{itemize}
\item À la fin de cette leçon, je sais comprendre globalement puis en détail un texte allemand sur l'exode rural
\item je sais définir et employer le vocabulaire thématique : das Dorf, die Stadt, die Landflucht, die Arbeit, umziehen, die Armut
\item je sais donner l'article et le pluriel des noms du thème et les classer par familles de mots
\item je sais exprimer la cause avec weil, denn et wegen en plaçant correctement le verbe
\item je sais exprimer la conséquence avec deshalb, darum, so dass
\item je sais répondre par écrit et à l'oral à des questions de compréhension en phrases complètes
\item je sais résumer un texte court en allemand en quelques phrases
\item je sais rédiger un court paragraphe guidé sur les causes et les conséquences de l'exode rural
\end{itemize}
\end{objectifs}

\begin{prerequis}
\begin{itemize}
\item la conjugaison des verbes réguliers et irréguliers au présent (Seconde)
\item la place du verbe en position 2 dans la phrase principale et en fin de subordonnée (début de 1ère)
\item les articles définis der/die/das et le pluriel des noms courants
\item le vocabulaire de base de la famille, du travail et de la ville (Seconde)
\item les questions en W- (Was, Warum, Wohin, Wie)
\end{itemize}
\end{prerequis}

\newpage

\coursec{Introduction et découverte du thème}

\courssub{Situation d'accroche : un phénomène que tu connais déjà}

Ferme les yeux un instant et imagine la scène. Il est six heures du matin à la gare routière de Mvan, à Yaoundé. Un jeune homme de dix-neuf ans serre sa mère dans ses bras, monte dans un car bondé, un sac sur les épaules. Il quitte son village de la région du Centre pour tenter sa chance en ville : trouver un travail, continuer ses études, peut-être devenir chauffeur, commerçant ou fonctionnaire. À Douala comme à Yaoundé, chaque année, des milliers de jeunes quittent ainsi leurs villages pour chercher du travail en ville.

Ce phénomène n'est pas seulement camerounais. Le même mouvement existe en Allemagne, en Autriche et en Suisse : les habitants des campagnes partent vers les grandes villes comme Berlin, Hambourg ou Munich. En allemand, on appelle cela \all{die Landflucht}, littéralement « la fuite de la campagne ».

Voici la problématique de notre leçon : pourquoi les habitants des campagnes partent-ils ? Que trouvent-ils en ville ? Et que devient le village abandonné ? Pour répondre à ces questions, nous allons lire un texte allemand intitulé \all{„Warum verlassen die jungen Leute das Dorf?“} (« Pourquoi les jeunes quittent-ils le village ? »), retenir le vocabulaire du thème et apprendre à exprimer la cause avec \all{weil} et la conséquence.

\begin{aretenir}[Objectifs de la leçon]
À la fin de cette leçon, tu seras capable de :
\begin{itemize}
\item comprendre un texte allemand simple sur l'exode rural (compréhension globale puis détaillée) ;
\item employer correctement le vocabulaire du thème : \all{das Dorf}, \all{die Stadt}, \all{die Landflucht}, \all{die Arbeit}, \all{umziehen}, \all{die Armut} ;
\item exprimer la cause avec la subordonnée en \all{weil} et la conséquence avec \all{deshalb} ou \all{darum} ;
\item répondre à des questions de compréhension et rédiger un court paragraphe sur ce sujet.
\end{itemize}
\end{aretenir}

\courssub{Brainstorming : pourquoi quitte-t-on le village ?}

Avant de lire le moindre texte, mettons en commun ce que nous savons déjà. En classe entière, réponds à la question : \emph{Pourquoi les jeunes quittent-ils le village pour la ville ?} Voici les réponses les plus fréquentes, déjà traduites en allemand — observe bien les articles et les pluriels :

\begin{center}
\begin{tabularx}{\linewidth}{|X|X|X|}
\hline
\rowcolor{popblueL} Français & Deutsch & Remarque \\
\hline
le travail & die Arbeit (Pl. die Arbeiten) & nom féminin ; on dit \all{Arbeit suchen} = chercher du travail \\
\hline
les études & das Studium (Pl. die Studien) & \all{studieren} = faire des études (université) \\
\hline
l'école, la formation & die Ausbildung (Pl. die Ausbildungen) & formation professionnelle \\
\hline
l'hôpital, la santé & das Krankenhaus (Pl. die Krankenhäuser) & les hôpitaux sont en ville \\
\hline
la pauvreté & die Armut (sans pluriel) & nom abstrait, toujours au singulier \\
\hline
l'argent & das Geld (Pl. die Gelder, rare) & \all{Geld verdienen} = gagner de l'argent \\
\hline
la vie moderne & das moderne Leben & cinéma, internet, loisirs \\
\hline
\end{tabularx}
\end{center}

\begin{exemplebox}[Premières phrases de brainstorming]
Voici comment formuler une hypothèse en allemand, avec une phrase simple :
\begin{itemize}
\item \all{Die jungen Leute suchen Arbeit in der Stadt.} « Les jeunes cherchent du travail en ville. »
\item \all{Im Dorf gibt es keine Schule.} « Au village, il n'y a pas d'école. » (structure \all{es gibt} + accusatif)
\item \all{Sie wollen in der Stadt studieren.} « Ils veulent étudier en ville. »
\item \all{Das Leben in der Stadt ist moderner.} « La vie en ville est plus moderne. »
\end{itemize}
\end{exemplebox}

\begin{attention}[Piège : \all{es gibt} et l'accusatif]
Le francophone traduit souvent « il y a » par \all{es ist} ou \all{es sind}. Faux ! En allemand, « il y a » se dit \all{es gibt}, et le nom qui suit est à l'\textbf{accusatif} : \all{Im Dorf gibt es keine Schule.} (et non \all{keine Schule ist}). Retiens : \all{es gibt} + accusatif, au singulier comme au pluriel : \all{In der Stadt gibt es viele Krankenhäuser.}
\end{attention}

\courssub{Le mot-clé de la leçon : \cle{die Landflucht}}

Observons d'abord le mot dans une phrase authentique : \all{Die Landflucht ist ein großes Problem für viele Dörfer.} « L'exode rural est un grand problème pour beaucoup de villages. »

Ce mot est une \textbf{composition} : l'allemand adore coller deux noms pour en créer un nouveau. Décomposons-le :

\begin{center}
\begin{tabular}{|l|l|l|}
\hline
\rowcolor{popgreenL} Élément & Sens & Genre \\
\hline
das Land & la campagne (aussi : le pays) & neutre \\
\hline
die Flucht & la fuite & féminin \\
\hline
die Land + Flucht & la fuite de la campagne = l'exode rural & féminin \\
\hline
\end{tabular}
\end{center}

\begin{definition}[die Landflucht]
\all{Die Landflucht} (sans pluriel usuel) désigne le \textbf{départ massif des habitants de la campagne vers les villes}, à la recherche de travail, d'études ou d'une vie meilleure. En français, on dit « l'exode rural ». Attention au genre : dans un mot composé allemand, c'est le \textbf{dernier élément} qui donne le genre et le pluriel. Ici, \all{die Flucht} est féminin, donc \all{die Landflucht} est féminin.
\end{definition}

\begin{propriete}[Règle d'or des noms composés]
Dans un nom composé allemand \all{A + B}, le sens précis et le genre grammatical sont donnés par le \textbf{dernier} nom (B) ; le premier nom (A) précise ou qualifie. Exemples :
\begin{itemize}
\item \all{das Dorf} + \all{die Schule} = \all{die Dorfschule} « l'école du village » (féminin, comme \all{Schule}) ;
\item \all{die Stadt} + \all{das Leben} = \all{das Stadtleben} « la vie en ville » (neutre, comme \all{Leben}) ;
\item \all{die Arbeit} + \all{der Platz} = \all{der Arbeitsplatz} « le poste de travail » (masculin, comme \all{Platz}).
\end{itemize}
Remarque la petite lettre de liaison parfois ajoutée : \all{Dorf\textbf{s}schule}, \all{Arbeit\textbf{s}platz}.
\end{propriete}

\begin{attention}[Faux ami : \cle{das Land}]
\all{Das Land} ne signifie pas seulement « la campagne ». Il a trois sens : (1) le pays (\all{Deutschland ist ein Land}) ; (2) la campagne, par opposition à la ville (\all{auf dem Land leben} = vivre à la campagne) ; (3) un Land fédéré allemand (\all{Bayern ist ein Bundesland}). Dans \all{Landflucht}, c'est le sens « campagne » qui compte. Et ne confonds pas avec le français « lande » !
\end{attention}

\courssub{Carte mentale : autour du mot Landflucht}

Pour mémoriser durablement un mot, le meilleur professeur te conseille de le placer au centre d'une \textbf{carte mentale} qui relie ses idées voisines : les causes (\all{die Ursachen}), les conséquences (\all{die Folgen}), le village (\all{das Dorf}) et la ville (\all{die Stadt}).

\begin{popfigure}
\begin{tikzpicture}[mindmap, grow cyclic, every node/.style={concept}, concept color=popblue!60, text=white,
level 1/.append style={level distance=3.4cm, sibling angle=90},
level 2/.append style={level distance=2.2cm, sibling angle=45}]
\node {DIE LANDFLUCHT}
child[concept color=poporange!80] { node[align=center]{Ursachen\\ (causes)}
  child { node[align=center] {Armut\\ Arbeitslosigkeit} }
  child { node[align=center] {keine Schulen\\ keine Krankenhäuser} }
}
child[concept color=popgreen!80] { node[align=center]{Folgen\\ (conséquences)}
  child { node[align=center] {leere Dörfer\\ alte Menschen} }
  child { node[align=center] {volle Städte\\ teure Wohnungen} }
}
child[concept color=popgold!85] { node[align=center]{das Dorf\\ (le village)}
  child { node[align=center] {ruhig, klein\\ die Landwirtschaft} }
}
child[concept color=poppurple!80] { node[align=center]{die Stadt\\ (la ville)}
  child { node[align=center] {Arbeit, Schulen\\ das moderne Leben} }
};
\end{tikzpicture}
\legende{Carte mentale du vocabulaire : die Landflucht et ses champs lexicaux}
\end{popfigure}

Recopie cette carte dans ton cahier et enrichis-la au fur et à mesure de la leçon : c'est un excellent outil de révision avant l'évaluation.

\courssub{Comparaison Cameroun / Allemagne : deux pays, un même mouvement}

Regardons maintenant le phénomène des deux côtés du monde. Au Cameroun, les jeunes quittent les villages de l'Ouest, du Nord ou du Centre pour Yaoundé et Douala ; les villes grandissent très vite, et certains villages ne comptent plus que des personnes âgées. En Allemagne, c'est surtout dans les Länder de l'Est (l'ancienne Allemagne de l'Est) que les petits villages se vident : les jeunes partent vers Berlin, Leipzig ou les régions riches de l'Ouest.

\begin{center}
\begin{tabularx}{\linewidth}{|X|X|X|}
\hline
\rowcolor{popblueL} Thème & Au Cameroun & En Allemagne \\
\hline
Destinations principales & Yaoundé, Douala & \all{Berlin, Hamburg, München} \\
\hline
Motifs (Gründe) & Travail, études, santé, vie moderne & \all{Arbeit, Studium, Gesundheit, modernes Leben} \\
\hline
Conséquence pour les villages & Perte des forces vives, vieillissement & \all{Leere Dörfer, Schulschließungen, Überalterung} \\
\hline
Conséquence pour les villes & Croissance rapide, logements précaires & \all{Zuzug, Wohnungsnot, steigende Mieten} \\
\hline
\end{tabularx}
\end{center}

\begin{savaistu}[Les « villages qui meurent » en Allemagne]
En Allemagne, certains petits villages de l'Est perdent chaque année des habitants. Les journalistes allemands parlent de \all{„sterbende Dörfer“}, littéralement « villages mourants » : l'école ferme, puis la boulangerie, puis le petit magasin, et il ne reste que des personnes âgées. Le verbe \all{sterben} (il meurt, il est mort) est un verbe fort : \all{das Dorf stirbt} « le village meurt ». Ce mot dramatique montre combien le sujet est pris au sérieux dans les pays germanophones.
\end{savaistu}

\courssub{Un premier texte de découverte}

Avant le grand texte de la section suivante, voici un court texte de mise en bouche. Lis-le une première fois sans dictionnaire : cherche seulement l'idée générale.

\begin{textebox}[Ein Dorf in Ostdeutschland]
Kleinstein ist ein kleines Dorf in Ostdeutschland. Vor zwanzig Jahren lebten hier 800 Menschen. Heute leben nur noch 300 Menschen im Dorf, und die meisten sind alt. Warum? Die jungen Leute gehen in die Stadt, weil es im Dorf keine Arbeit gibt. Die Schule ist schon geschlossen, weil es zu wenige Kinder gibt. Auch der Bäcker und der kleine Laden sind weg.

Anna ist neunzehn Jahre alt. Nächstes Jahr will sie nach Leipzig umziehen. „Ich liebe mein Dorf“, sagt sie, „aber hier finde ich keine Ausbildung. In der Stadt gibt es Universitäten, Krankenhäuser und viele Arbeitsplätze.“ Ihr Großvater bleibt im Dorf. Er sagt: „Die Stadt ist zu laut und zu teuer. Hier ist mein Zuhause.“

Die Landflucht ist ein Problem für Kleinstein. Der Bürgermeister hat eine Idee: Er will jungen Familien billige Häuser anbieten. Vielleicht kommen dann neue Menschen ins Dorf.
\end{textebox}

\textbf{Glossaire :} \all{vor zwanzig Jahren} « il y a vingt ans » ; \all{nur noch} « seulement encore » ; \all{geschlossen} « fermé » (participe de \all{schließen}) ; \all{der Bäcker, die Bäcker} « le boulanger » ; \all{der Laden, die Läden} « le magasin » ; \all{weg} « parti » ; \all{umziehen (zieht um, zog um, ist umgezogen)} « déménager » ; \all{der Arbeitsplatz, die Arbeitsplätze} « le poste de travail » ; \all{laut} « bruyant » ; \all{der Bürgermeister, die Bürgermeister} « le maire » ; \all{billig} « bon marché » ; \all{anbieten} « proposer, offrir ».

\textbf{Questions rapides (réponds en allemand, phrase complète) :}
\begin{enumerate}
\item \all{Wie viele Menschen leben heute in Kleinstein?}
\item \all{Warum gehen die jungen Leute in die Stadt?}
\item \all{Warum will Anna nach Leipzig umziehen?}
\item \all{Was will der Bürgermeister anbieten?}
\end{enumerate}

\begin{exoresolu}[Première question de compréhension résolue]
Réponds à la question : \all{Warum gehen die jungen Leute in die Stadt?}
\tcblower
Méthode : je repère dans le texte la phrase qui contient la cause : \all{„Die jungen Leute gehen in die Stadt, weil es im Dorf keine Arbeit gibt.“} Je recopie cette phrase ou je la reformule. Réponse attendue : \all{Die jungen Leute gehen in die Stadt, weil es im Dorf keine Arbeit gibt.} « Les jeunes partent en ville parce qu'il n'y a pas de travail au village. » Remarque bien la place du verbe \all{gibt} à la fin de la subordonnée introduite par \all{weil} : c'est la règle de grammaire centrale de cette leçon, que nous étudierons en détail dans la section 5.
\end{exoresolu}

\courssub{Anticiper avant de lire : titre, image, hypothèses}

Un bon lecteur ne lit jamais un texte « à froid ». Il anticipe : il prédit le contenu avant de lire. Cela active tes connaissances et rend la compréhension deux fois plus facile.

\begin{methode}[Anticiper en trois étapes]
\begin{enumerate}
\item \textbf{Lire le titre.} Notre texte s'intitule \all{„Warum verlassen die jungen Leute das Dorf?“}. Le mot interrogatif \all{warum} (« pourquoi ») annonce que le texte donnera des \textbf{causes}. Le verbe \all{verlassen} (quitter) signifie « quitter ».
\item \textbf{Regarder l'image} (ou la carte mentale) : jeunes gens avec des valises, un village presque vide, une grande ville au loin. L'image confirme le thème du départ.
\item \textbf{Formuler deux hypothèses} avant de lire, par exemple : \all{Ich denke, die jungen Leute suchen Arbeit.} « Je pense que les jeunes cherchent du travail. » et \all{Vielleicht gibt es im Dorf keine Schulen.} « Peut-être qu'il n'y a pas d'écoles au village. » Pendant la lecture, tu vérifieras si tes hypothèses étaient justes.
\end{enumerate}
\end{methode}

\begin{exemplebox}[Phrases utiles pour formuler des hypothèses]
\begin{itemize}
\item \all{Ich denke, das Dorf ist sehr arm.} « Je pense que le village est très pauvre. »
\item \all{Vielleicht wollen die jungen Leute in der Stadt studieren.} « Peut-être que les jeunes veulent étudier en ville. »
\item \all{Ich glaube, das Leben in der Stadt ist interessanter.} « Je crois que la vie en ville est plus intéressante. »
\end{itemize}
Remarque : après \all{ich denke} et \all{ich glaube} suivis d'une virgule, la phrase qui suit garde l'ordre normal (verbe en deuxième position) ; avec \all{vielleicht} en tête, le verbe vient juste après : \all{Vielleicht \textbf{gibt} es...}
\end{exemplebox}

\begin{experience}[À toi de jouer : tour de table des hypothèses]
En groupes de quatre, formulez chacun une hypothèse sur le contenu du texte \all{„Warum verlassen die jungen Leute das Dorf?“}, en utilisant \all{ich denke}, \all{ich glaube} ou \all{vielleicht}. Le secrétaire du groupe note les quatre hypothèses au tableau. À la fin de la lecture du texte (section suivante), la classe vérifiera quel groupe avait les meilleures prédictions.
\end{experience}

\begin{aretenir}[Bilan de la section]
Dans cette introduction, tu as découvert le thème de l'exode rural (\all{die Landflucht}), un phénomène présent au Cameroun comme en Allemagne ; tu as appris à décomposer les noms composés (le dernier élément donne le genre) ; tu as construit une carte mentale du vocabulaire ; et tu sais maintenant anticiper un texte en trois étapes : titre, image, hypothèses. Nous sommes prêts pour la lecture du texte support.
\end{aretenir}

\coursec{Texte support : „Warum verlassen die jungen Leute das Dorf?“}

Après avoir découvert le thème de l'exode rural à travers des images et une première discussion en section 1, nous abordons maintenant le cœur de la leçon : un texte authentique qui illustre ce phénomène aussi bien en Allemagne qu'en Afrique. Ce texte servira de base à tout le travail de compréhension, de lexique et de grammaire qui suivra.

\courssub{Découverte du texte support}

\begin{textebox}[Texte original : Warum verlassen die jungen Leute das Dorf?]
In vielen Dörfern Deutschlands und Afrikas gibt es ein großes Problem: die Landflucht. Viele junge Leute verlassen das Dorf, weil es dort keine Arbeit gibt. Sie ziehen in die Stadt um, weil sie hoffen, dort einen Job und ein besseres Leben zu finden. In der Stadt gibt es Fabriken, Büros, Krankenhäuser und Universitäten. Aber das Leben dort ist teuer, und nicht alle finden Arbeit. Manche leben in Armut. Im Dorf bleiben nur die Alten. Die Felder liegen brach, die Schulen schließen. Deshalb müssen die Regierungen etwas tun: Straßen bauen, Schulen und Krankenhäuser im Dorf öffnen und Arbeit auf dem Land schaffen.
\tcblower
\textbf{Glossaire (mots nouveaux) :}
\begin{itemize}
    \item \all{verlassen} (verlässt, verließ, hat verlassen) — \emph{quitter} (un lieu, une personne)
    \item \all{umziehen} (zieht um, zog um, ist umgezogen) — \emph{déménager} ; \all{in die Stadt umziehen} — \emph{déménager en ville}
    \item \all{hoffen} (hofft, hoffte, hat gehofft) — \emph{espérer} ; \all{hoffen, dass…} — \emph{espérer que…} ; \all{hoffen, … zu + infinitif} — \emph{espérer …}
    \item \all{die Fabrik, die Fabriken} — \emph{l'usine}
    \item \all{das Krankenhaus, die Krankenhäuser} — \emph{l'hôpital}
    \item \all{teuer} — \emph{cher, coûteux} (antonyme : \all{billig} — \emph{bon marché})
    \item \all{die Armut} (singulier seulement) — \emph{la pauvreté}
    \item \all{brach liegen} — \emph{être en friche, rester en jachère} (\all{die Felder liegen brach} — \emph{les champs sont en friche})
    \item \all{schaffen} (schafft, schuf, hat geschaffen) — \emph{créer, fonder, mettre en place} (ici : \emph{créer des emplois})
    \item \all{die Landflucht} — \emph{l'exode rural} (litt. « la fuite de la campagne »)
    \item \all{das Dorf, die Dörfer} — \emph{le village}
    \item \all{die Stadt, die Städte} — \emph{la ville}
    \item \all{der Job, die Jobs} — \emph{l'emploi, le boulot} (synonyme : \all{die Stelle, die Stellen})
    \item \all{die Regierung, die Regierungen} — \emph{le gouvernement}
    \item \all{die Straße, die Straßen} — \emph{la route}
    \item \all{die Schule, die Schulen} — \emph{l'école}
    \item \all{das Feld, die Felder} — \emph{le champ}
    \item \all{die Alten} (substantivé, pluriel) — \emph{les personnes âgées, les anciens}
\end{itemize}
\end{textebox}

\courssub{Lecture globale (compréhension générale)}

\begin{methode}[Comment faire une lecture globale efficace]
\begin{enumerate}
    \item \textbf{Lire une première fois sans dictionnaire.} Ne vous arrêtez pas sur chaque mot inconnu ; cherchez le sens global.
    \item \textbf{Repérer les indices visuels :} titre, mots transparents (\all{Problem, Fabriken, Universitäten, Regierungen}), structure en paragraphes.
    \item \textbf{Répondre aux questions « W-Fragen » essentielles :}
    \begin{itemize}
        \item \all{Wer?} (Qui ?) — Qui sont les acteurs principaux ?
        \item \all{Was?} (Quoi ?) — De quoi parle le texte ? Quel est le thème ?
        \item \all{Wo?} (Où ?) — Où se passe l'action ? Quels lieux sont mentionnés ?
        \item \all{Warum?} (Pourquoi ?) — Quelles sont les raisons données ?
    \end{itemize}
    \item \textbf{Résumer en une phrase en français} l'idée principale du texte.
\end{enumerate}
\end{methode}

\begin{exemplebox}[Application de la méthode au texte]
\textbf{Titre :} „Warum verlassen die jungen Leute das Dorf?“ → \emph{Pourquoi les jeunes quittent-ils le village ?}

\textbf{Acteurs (Wer) :} \all{junge Leute} (les jeunes), \all{die Alten} (les anciens), \all{die Regierungen} (les gouvernements).

\textbf{Thème (Was) :} \all{die Landflucht} (l'exode rural) et ses conséquences.

\textbf{Lieux (Wo) :} \all{Dörfer} (villages) en \all{Deutschland} (Allemagne) et \all{Afrika} (Afrique) ; \all{Stadt} (ville) ; \all{auf dem Land} (à la campagne).

\textbf{Raisons (Warum) :} \all{keine Arbeit} (pas de travail), espoir d'un \all{besseres Leben} (meilleure vie) en ville.

\textbf{Résumé :} Le texte explique pourquoi les jeunes quittent les villages pour la ville (manque de travail), les conséquences pour les villages (vieillissement, champs en friche, écoles fermées) et propose des solutions (infrastructures, création d'emplois ruraux).
\end{exemplebox}

\begin{exoresolu}[Questions de compréhension globale]
\textbf{Énoncé :} Répondez en français aux questions suivantes d'après le texte.

1. Quel est le « grand problème » mentionné dans la première phrase ?
2. Qui sont les personnes qui partent ? Où vont-elles ?
3. Donnez deux raisons pour lesquelles elles partent.
4. Citez trois infrastructures que l'on trouve en ville selon le texte.
5. Quelles sont les deux conséquences négatives pour le village mentionnées à la fin du premier paragraphe ?
6. Que propose le texte comme solutions ?

\textbf{Solution détaillée :}
1. Le grand problème est \textbf{l'exode rural} (\all{die Landflucht}).
2. Ce sont \textbf{les jeunes gens} (\all{viele junge Leute}) qui partent. Ils vont \textbf{en ville} (\all{in die Stadt}).
3. Raisons : \textbf{il n'y a pas de travail au village} (\all{keine Arbeit gibt}) et \textbf{ils espèrent trouver un emploi et une vie meilleure en ville} (\all{hoffen, dort einen Job und ein besseres Leben zu finden}).
4. \textbf{Usines, bureaux, hôpitaux, universités} (\all{Fabriken, Büros, Krankenhäuser und Universitäten}).
5. \textbf{Seuls les anciens restent} (\all{nur die Alten bleiben}) ; \textbf{les champs sont en friche} (\all{Felder liegen brach}) et \textbf{les écoles ferment} (\all{Schulen schließen}).
6. Les gouvernements doivent \textbf{construire des routes, ouvrir des écoles et des hôpitaux au village, et créer du travail à la campagne} (\all{Straßen bauen, Schulen und Krankenhäuser im Dorf öffnen und Arbeit auf dem Land schaffen}).
\end{exoresolu}

\courssub{Lecture détaillée (compréhension fine) : causes et conséquences}

Nous allons maintenant analyser la structure logique du texte. Le texte présente une \cle{relation de cause à effet} (Ursache–Wirkung) très claire : le départ des jeunes a des causes (pourquoi ils partent) et des conséquences (ce qui arrive au village).

\begin{center}
\begin{tabularx}{\linewidth}{|X|X|X|}
\hline
\rowcolor{popblueL}
\textbf{Élément du texte (Allemand)} & \textbf{Traduction (Français)} & \textbf{Rôle logique} \\ \hline
\all{weil es dort keine Arbeit gibt} & parce qu'il n'y a pas de travail là-bas & \textbf{Cause 1} du départ \\ \hline
\all{weil sie hoffen, dort einen Job und ein besseres Leben zu finden} & parce qu'ils espèrent y trouver un emploi et une vie meilleure & \textbf{Cause 2} du départ \\ \hline
\all{In der Stadt gibt es Fabriken, Büros, Krankenhäuser und Universitäten} & En ville il y a des usines, bureaux, hôpitaux, universités & \textbf{Facteur d'attraction} (Pull-Faktor) \\ \hline
\all{Aber das Leben dort ist teuer, und nicht alle finden Arbeit} & Mais la vie y est chère et tous ne trouvent pas du travail & \textbf{Nuance / Réalité} \\ \hline
\all{Manche leben in Armut} & Certains vivent dans la pauvreté & \textbf{Conséquence en ville} \\ \hline
\all{Im Dorf bleiben nur die Alten} & Au village ne restent que les anciens & \textbf{Conséquence 1} au village \\ \hline
\all{Die Felder liegen brach} & Les champs sont en friche & \textbf{Conséquence 2} au village \\ \hline
\all{die Schulen schließen} & les écoles ferment & \textbf{Conséquence 3} au village \\ \hline
\all{Deshalb müssen die Regierungen etwas tun} & C'est pourquoi les gouvernements doivent agir & \textbf{Solution politique} (conséquence globale) \\ \hline
\end{tabularx}
\end{center}

\begin{propriete}[La subordonnée causale avec \all{weil} (Ursachensatz)]
\emph{Niveau A1+/A2 — Point de grammaire obligatoire du programme.}
\begin{itemize}
    \item \textbf{Structure :} \all{weil} + sujet + compléments + \textbf{verbe conjugué à la fin}.
    \item \textbf{Fonction :} Exprime la \textbf{cause} (raison, motif). Répond à la question \all{Warum?}.
    \item \textbf{Place dans la phrase :}
    \begin{itemize}
        \item Après la principale : \all{Hauptsatz, weil Nebensatz.} (Pas de virgule obligatoire en nouvelle orthographe, mais recommandée pour la clarté).
        \item Avant la principale : \all{Weil Nebensatz, Hauptsatz.} (La subordonnée occupe la place 1, le verbe de la principale suit immédiatement).
    \end{itemize}
    \item \textbf{Conjonctions synonymes (registre plus soutenu) :} \all{da} (verbe à la fin aussi), \all{denn} (conjonction de coordination, verbe en 2e position, ne prend pas de virgule devant).
\end{itemize}
\end{propriete}

\begin{aretenir}[Vocabulaire logique : exprimer la cause et la conséquence]
\begin{itemize}
    \item \textbf{Cause (Ursache) :} \all{weil} + subordonnée (verbe à la fin), \all{da} + subordonnée, \all{wegen} + génitif (niveau supérieur/B1).
    \item \textbf{Conséquence (Folge) — Adverbes (verbe en 2e position) :} \all{deshalb} / \all{deswegen} / \all{darum} (place 1 ou 3), \all{also} (donc, en conclusion, place 1 ou 3), \all{folglich} (par conséquent, formel, place 1 ou 3).
    \item \textbf{Structure type du texte argumentatif :} \all{Problem → Ursachen (weil...) → Situation actuelle (Aber...) → Folgen (Deshalb...) → Forderung/Lösung}.
\end{itemize}
\end{aretenir}

\courssub{Prononciation : lecture à voix haute}

Entraînez-vous à lire le texte à haute voix. Portez une attention particulière aux mots suivants, sources d'erreurs fréquentes pour les francophones :

\begin{center}
\begin{tabular}{|l|l|l|}
\hline
\rowcolor{poporangeL}
\textbf{Mot allemand} & \textbf{Prononciation approximative (français)} & \textbf{Point d'attention} \\ \hline
\all{Landflucht} & « lant-floocht » & \all{au} = [aʊ] (comme dans \all{Haus}) ; \all{ch} après \all{u} = [x] (guttural, comme \all{Buch}) \\ \hline
\all{verlassen} & « fèr-lass-en » & \all{v} initial = [f] (comme \all{Vater}) ; \all{a} court ; \all{ss} = [s] dur \\ \hline
\all{Armut} & « ar-mout » & \all{u} court [ʊ] (pas [y] comme \all{ü}) ; \all{t} final dur \\ \hline
\all{brach} & « brach » & \all{ch} après \all{a} = [x] (guttural) \\ \hline
\all{schaffen} & « chaf-fen » & \all{sch} = [ʃ] (ch français) ; \all{a} court ; \all{ff} = [f] \\ \hline
\all{Regierungen} & « ré-gué-roung-en » & \all{eu} = [ɔʏ] (comme \all{neu}) ; \all{ng} = [ŋ] ; \all{en} final = [ən] \\ \hline
\end{tabular}
\end{center}

\begin{methode}[Entraînement à la lecture expressive]
\begin{enumerate}
    \item Écoutez le modèle (enseignant ou audio) en suivant le texte des yeux.
    \item Répétez phrase par phrase en imitant l'intonation (montante pour les causes introduites par \all{weil}, descendante pour les affirmations).
    \item Marquez les pauses aux virgules et aux points.
    \item Enregistrez-vous et comparez.
\end{enumerate}
\end{methode}

\courssub{Schémas des causes et conséquences}

\begin{popfigure}
\begin{tikzpicture}[
    node distance=1.2cm and 1.5cm,
    cause/.style={rectangle, rounded corners, draw=poporange, thick, fill=poporangeL, text width=3.5cm, align=center, font=\small},
    consequence/.style={rectangle, rounded corners, draw=popgreen, thick, fill=popgreenL, text width=3.5cm, align=center, font=\small},
    arrow/.style={-Stealth, thick, draw=popdark},
    place/.style={rectangle, rounded corners, draw=popblue, thick, fill=popblueL, text width=2.5cm, align=center, font=\bfseries\small},
    labelstyle/.style={font=\small\bfseries, text=popdark}
]
% Lieux
\node[place] (dorf) {Dorf};
\node[place, right=of dorf] (stadt) {Stadt};

% Flèche principale
\draw[arrow, very thick] (dorf.east) -- node[above, labelstyle, yshift=0.4cm] {Landflucht} (stadt.west);

% Causes au-dessus
\node[cause, above=of dorf, xshift=-1cm] (c1) {keine Arbeit};
\node[cause, above=of stadt, xshift=1cm] (c2) {Hoffnung auf Job \& besseres Leben};
\draw[arrow, draw=poporange] (c1) -- (dorf.north);
\draw[arrow, draw=poporange] (c2) -- (stadt.north);

% Conséquences en-dessous
\node[consequence, below=of dorf, xshift=-1cm] (f1) {nur die Alten bleiben};
\node[consequence, below=of dorf, xshift=1.5cm] (f2) {Felder liegen brach};
\node[consequence, below=of stadt, xshift=-1.5cm] (f3) {Schulen schließen};
\node[consequence, below=of stadt, xshift=1cm] (f4) {teuer, Armut in der Stadt};
\draw[arrow, draw=popgreen] (dorf.south) -- (f1);
\draw[arrow, draw=popgreen] (dorf.south) -- (f2);
\draw[arrow, draw=popgreen] (stadt.south) -- (f3);
\draw[arrow, draw=popgreen] (stadt.south) -- (f4);

% Légendes des zones
\node[labelstyle, above=0.2cm of c1] {Ursachen (Causes)};
\node[labelstyle, above=0.2cm of c2] {Pull-Faktoren};
\node[labelstyle, below=0.2cm of f2] {Folgen für das Dorf};
\node[labelstyle, below=0.2cm of f4] {Realität in der Stadt};
\end{tikzpicture}
\legende{Schéma causal de l'exode rural : causes (orange) au-dessus, conséquences (vert) en dessous.}
\end{popfigure}

\courssub{Point de vigilance : faux amis et pièges lexicaux}

\begin{attention}[« das Land » ≠ « le pays » ; « die Stadt » = la ville]
\begin{itemize}
    \item \textbf{\all{das Land}} (pl. \all{die Länder}) a deux sens principaux :
    \begin{enumerate}
        \item \emph{Le pays, l'État} (ex. : \all{Deutschland ist ein Land} — L'Allemagne est un pays).
        \item \emph{La campagne} (opposé à \all{die Stadt}), surtout avec la préposition \all{auf dem Land} (\emph{à la campagne}) ou \all{aufs Land} (\emph{à la campagne, vers la campagne}).
    \end{enumerate}
    \textbf{Piège :} Dans notre texte, \all{Arbeit auf dem Land schaffen} signifie « créer du travail \emph{à la campagne} », pas « dans le pays ».
    \item \textbf{\all{die Stadt}} (pl. \all{die Städte}) = \emph{la ville} (agglomération urbaine). Ne pas confondre avec \all{der Staat} (l'État).
    \item \textbf{\all{umziehen}} (verbe à particule séparable) : \all{Er zieht nach Yaoundé um.} (Il déménage à Yaoundé.) — La particule \all{um} va à la fin en phrase principale.
    \item \textbf{\all{verlassen}} (verbe transitif, pas de préposition) : \all{Er verlässt das Dorf.} — Ne pas dire *\all{verlässt von dem Dorf}.
    \item \textbf{\all{die Stelle, die Stellen}} = \emph{le poste, l'emploi} (concret). \all{die Arbeit} = \emph{le travail} (activité, notion abstraite ou emploi en général).
\end{itemize}
\end{attention}

\begin{exemplebox}[Exemples traduits : cause (\all{weil}) et conséquence (\all{deshalb})]
\begin{enumerate}
    \item \all{Viele junge Leute verlassen das Dorf, weil es dort keine Arbeit gibt.}
    « Beaucoup de jeunes quittent le village \textbf{parce qu'}il n'y a pas de travail là-bas. »
    \item \all{Sie ziehen in die Stadt um, weil sie einen Job suchen.}
    « Ils déménagent en ville \textbf{parce qu'}ils cherchent un emploi. »
    \item \all{Im Dorf gibt es keine Fabriken. Deshalb ziehen die Jungen weg.}
    « Au village il n'y a pas d'usines. \textbf{C'est pourquoi / Donc} les jeunes partent. »
    \item \all{Das Leben in der Stadt ist teuer. Deswegen leben manche in Armut.}
    « La vie en ville est chère. \textbf{C'est pour cela / Par conséquent} certains vivent dans la pauvreté. »
    \item \all{Die Schulen schließen, weil keine Kinder mehr da sind.}
    « Les écoles ferment \textbf{parce qu'}il n'y a plus d'enfants. »
\end{enumerate}
\end{exemplebox}

\begin{exoresolu}[Exercice résolu : relier causes et conséquences avec \all{weil} / \all{deshalb}]
\textbf{Énoncé :} Reliez les deux phrases en une seule phrase logique. Utilisez \all{weil} (cause) ou \all{deshalb} (conséquence). Conjuguez correctement.

1. Die jungen Leute gehen weg. \_\_\_\_\_\_\_\_ es keine Arbeit gibt. \\
2. Es gibt keine Fabriken im Dorf. \_\_\_\_\_\_\_\_ suchen die Leute Arbeit in der Stadt. \\
3. Das Leben in der Stadt ist teuer. \_\_\_\_\_\_\_\_ haben manche kein Geld. \\
4. Die Schulen schließen. \_\_\_\_\_\_\_\_ keine Kinder mehr da sind. \\
5. Die Regierung baut Straßen. \_\_\_\_\_\_\_\_ sie will die Landflucht stoppen.

\textbf{Solution détaillée :}
1. \all{Die jungen Leute gehen weg, \textbf{weil} es \textbf{keine Arbeit gibt}.} (Subordonnée : verbe \all{gibt} à la fin.)
2. \all{Es gibt keine Fabriken im Dorf. \textbf{Deshalb} \textbf{suchen} die Leute Arbeit in der Stadt.} (Principal : verbe \all{suchen} en 2e position après \all{deshalb}.)
3. \all{Das Leben in der Stadt ist teuer. \textbf{Deshalb} \textbf{haben} manche kein Geld.} (Principal : verbe \all{haben} en 2e position.)
4. \all{Die Schulen schließen, \textbf{weil} keine Kinder mehr \textbf{da sind}.} (Subordonnée : verbe \all{sind} à la fin.)
5. \all{Die Regierung baut Straßen, \textbf{weil} sie die Landflucht \textbf{stoppen will}.} (Subordonnée : verbe \all{will} à la fin, \all{stoppen} avant.)
\end{exoresolu}

\courssub{Production écrite guidée}

\begin{methode}[Rédiger un court paragraphe argumentatif (niveau A2)]
\begin{enumerate}
    \item \textbf{Collecter les idées :} Faites un schéma cause/effet (comme p. 4) pour votre propre contexte (village/quartier d'origine).
    \item \textbf{Structurer :} 1. Constat (Problème) → 2. Causes (\all{weil}) → 3. Conséquences (\all{deshalb}/\all{deswegen}) → 4. Proposition (\all{müssen/sollen} + infinitif).
    \item \textbf{Rédiger en allemand :} Utilisez le vocabulaire de la leçon (\all{Landflucht, umziehen, Arbeit, Armut, schaffen}). Vérifiez la place du verbe (\all{weil} → fin ; \all{deshalb} → 2e place).
    \item \textbf{Relire :} Majuscules aux noms ? Umlaute (ä, ö, ü) ? ß ? Ponctuation ?
\end{enumerate}
\end{methode}

\begin{exercice}[Production écrite : L'exode rural chez nous]
\textbf{Consigne :} Rédigez un paragraphe de 6 à 8 lignes en allemand sur l'exode rural dans votre région (Cameroun : village → Yaoundé/Douala/Bafoussam...).
\begin{itemize}
    \item Utilisez au moins \textbf{deux fois \all{weil}} (cause) et \textbf{une fois \all{deshalb} ou \all{deswegen}} (conséquence).
    \item Employez le vocabulaire : \all{das Dorf, die Stadt, die Landflucht, die Arbeit, umziehen, die Armut, schaffen}.
    \item Modèle de début : \all{In meiner Heimat gibt es auch Landflucht. Viele junge Leute...}
\end{itemize}
\end{exercice}

\begin{corrige}[Exercice : Production écrite]
\textbf{Proposition de correction (exemple) :}
\all{In meiner Heimat gibt es auch Landflucht. Viele junge Leute verlassen das Dorf, weil es dort keine Arbeit gibt. Sie ziehen in die Stadt um, weil sie hoffen, Geld zu verdienen. In der Stadt gibt es viele Fabriken und Büros. Aber das Leben ist teuer. Deswegen leben manche in Armut. Im Dorf bleiben nur die Alten. Die Felder liegen brach. Deshalb muss die Regierung Schulen und Straßen bauen und Arbeit auf dem Land schaffen.}
\end{corrige}

\begin{savaistu}[Contexte comparatif : Cameroun et Allemagne]
Au Cameroun comme en Allemagne, l'exode rural (\all{Landflucht}) est un défi majeur. Au Cameroun, les jeunes quittent les zones rurales (ex. : Nord, Est) pour Yaoundé, Douala ou Bafoussam, attirés par les études, l'informel ou l'espoir d'un salaire. Conséquences : vieillissement des villages, terres non cultivées, pression sur les infrastructures urbaines (logement, transport \all{bendskin}, eau, électricité). En Allemagne, depuis les années 1960, la politique d'aménagement du territoire (\all{Raumordnung}) tente de maintenir l'attractivité des zones rurales par la décentralisation des administrations, le haut débit internet, le soutien à l'agriculture et aux PME. Le texte mentionne explicitement \all{Deutschland und Afrika} pour montrer l'universalité du phénomène.
\end{savaistu}

\coursec{Compréhension du texte}

Après avoir lu et écouté le texte support \all{„Warum verlassen die jungen Leute das Dorf ?“}, nous allons maintenant vérifier notre compréhension. C'est l'étape centrale de toute leçon de texte : d'abord une compréhension \emph{globale} (de quoi parle le texte ?), puis une compréhension \emph{détaillée} (que dit exactement le texte ?), enfin une exploitation personnelle (vrai ou faux, recherche d'informations, résumé). C'est exactement la démarche exigée à l'évaluation de fin de séquence et, plus tard, au Baccalauréat.

\courssub{La méthode pour répondre à une question de compréhension}

\begin{methode}[Comment répondre à une question de compréhension]
Pour répondre correctement à une question sur un texte, suis toujours ces quatre étapes :
\begin{enumerate}
\item \textbf{Lire la question attentivement} et souligner le mot interrogatif : \all{wer} (qui), \all{was} (quoi), \all{wo} (où), \all{wohin} (vers où), \all{warum} (pourquoi), \all{wann} (quand), \all{wie} (comment).
\item \textbf{Repérer dans le texte} la phrase ou le passage qui contient la réponse.
\item \textbf{Formuler une phrase complète} en reprenant les mots de la question : la question \all{Wohin ziehen die jungen Leute ?} appelle une réponse qui commence par \all{Die jungen Leute ziehen ...}.
\item \textbf{Justifier par une courte citation} du texte entre guillemets allemands, introduite par \all{Im Text steht : ...} ou \all{Der Text sagt : ...}.
\end{enumerate}
\end{methode}

\begin{attention}[Le piège de la réponse en un seul mot]
Ne réponds \textbf{jamais} par un seul mot ou par « oui / non ». À l'évaluation comme au Baccalauréat, une réponse en \textbf{phrase complète} est exigée, avec sujet, verbe conjugué à la deuxième position et compléments.

\all{Wohin ziehen die jungen Leute ?} — \emph{Mauvaise réponse :} \all{In die Stadt.}

\emph{Bonne réponse :} \all{Die jungen Leute ziehen in die Stadt, weil sie dort Arbeit suchen.} « Les jeunes s'installent en ville parce qu'ils y cherchent du travail. »
\end{attention}

\courssub{Compréhension globale du texte}

Les questions globales portent sur le sens général du texte. On y répond sans entrer dans les détails.

\begin{exemplebox}[Questions globales et réponses modèles]
\textbf{1.} \all{Worum geht es im Text ?} « De quoi parle le texte ? »

\all{Der Text handelt von der Landflucht : viele junge Leute verlassen das Dorf und ziehen in die Stadt.} « Le texte traite de l'exode rural : beaucoup de jeunes quittent le village et s'installent en ville. »

\textbf{2.} \all{Wer verlässt das Dorf ?} « Qui quitte le village ? »

\all{Die jungen Leute verlassen das Dorf, besonders die Jugendlichen nach der Schule.} « Ce sont les jeunes qui quittent le village, surtout les jeunes après l'école. »

\textbf{3.} \all{Wohin ziehen die jungen Leute ?} « Où vont s'installer les jeunes ? »

\all{Sie ziehen in die großen Städte, zum Beispiel nach Yaoundé oder Douala.} « Ils s'installent dans les grandes villes, par exemple à Yaoundé ou à Douala. »
\end{exemplebox}

Remarque la construction de chaque réponse : on reprend les mots de la question (\all{verlassen das Dorf}, \all{ziehen}) et on forme une phrase complète. Le verbe conjugué reste à la deuxième position, comme dans toute phrase déclarative allemande. Note aussi la préposition : on dit \all{nach Yaoundé} (sans article, pour les villes et la plupart des pays), mais \all{in die Stadt} (avec article, accusatif après \all{in} pour exprimer le mouvement).

\courssub{Compréhension détaillée du texte}

Les questions détaillées exigent de retrouver une information précise dans le texte et de la citer.

\begin{exemplebox}[Questions détaillées et réponses modèles]
\textbf{1.} \all{Warum verlassen die jungen Leute das Dorf ?} « Pourquoi les jeunes quittent-ils le village ? »

\all{Sie verlassen das Dorf, weil es dort keine Arbeit und keine Berufsschulen gibt. Im Text steht : „Auf dem Land finden die Jugendlichen keine Arbeit.“} « Ils quittent le village parce qu'il n'y a là ni travail ni écoles professionnelles. Le texte dit : „À la campagne, les jeunes ne trouvent pas de travail.“ »

\textbf{2.} \all{Was gibt es in der Stadt ?} « Qu'y a-t-il en ville ? »

\all{In der Stadt gibt es Fabriken, Büros, Krankenhäuser, Schulen und Universitäten.} « En ville, il y a des usines, des bureaux, des hôpitaux, des écoles et des universités. »

\textbf{3.} \all{Was passiert im Dorf ?} « Que se passe-t-il au village ? »

\all{Im Dorf bleiben nur die alten Leute. Die Felder liegen brach und die Schulen schließen.} « Au village, il ne reste que les personnes âgées. Les champs restent en friche et les écoles ferment. »

\textbf{4.} \all{Was müssen die Regierungen tun ?} « Que doivent faire les gouvernements ? »

\all{Die Regierungen müssen Straßen, Schulen und Krankenhäuser auf dem Land bauen und Arbeit für die Jugendlichen schaffen.} « Les gouvernements doivent construire des routes, des écoles et des hôpitaux à la campagne et créer du travail pour les jeunes. »
\end{exemplebox}

\begin{exemplebox}[Exemple traduit : la question avec \all{warum}]
\all{Warum ziehen die jungen Leute um ? – Sie ziehen um, weil sie Arbeit suchen.}

« Pourquoi les jeunes déménagent-ils ? – Ils déménagent parce qu'ils cherchent du travail. »

Observe : à la question \all{warum} (pourquoi), on répond avec une subordonnée en \all{weil}, et le verbe conjugué \all{suchen} va à la \textbf{fin} de la subordonnée. Nous étudierons cette règle en détail dans la partie grammaire.
\end{exemplebox}

\courssub{Vrai ou faux justifié}

L'exercice « vrai ou faux » (\all{richtig oder falsch}) est un classique de l'évaluation. La règle d'or : la justification est obligatoire et doit citer le texte.

\begin{exoresolu}[Richtig oder falsch ?]
Énoncé. Lis l'affirmation suivante et dis si elle est vraie (\all{richtig}) ou fausse (\all{falsch}). Justifie ta réponse par une phrase complète et une citation du texte.

\all{„In der Stadt finden alle Arbeit.“} « En ville, tout le monde trouve du travail. »
\tcblower
Solution détaillée. L'affirmation est \textbf{fausse} (\all{falsch}).

Réponse modèle : \all{Das ist falsch. Nicht alle jungen Leute finden in der Stadt Arbeit. Im Text steht : „Manche Jugendliche finden keinen Job und leben in Armut.“} « C'est faux. Tous les jeunes ne trouvent pas de travail en ville. Le texte dit : „Certains jeunes ne trouvent pas d'emploi et vivent dans la pauvreté.“ »

Méthode : on annonce d'abord le jugement (\all{Das ist falsch}), puis on reformule l'information correcte en phrase complète, enfin on cite le texte entre guillemets allemands.
\end{exoresolu}

\begin{attention}[Attention à la négation partielle]
Les francophones traduisent souvent mal \all{nicht alle}. \all{Nicht alle finden Arbeit} signifie « tout le monde ne trouve pas de travail » (certains en trouvent, d'autres non), et non « personne ne trouve de travail ». Pour « personne », l'allemand utilise \all{niemand} ou \all{keiner} : \all{Niemand findet Arbeit.} « Personne ne trouve de travail. » De même, \all{manche} signifie « certains » : \all{Manche Jugendliche finden keinen Job.} « Certains jeunes ne trouvent pas d'emploi. »
\end{attention}

\courssub{Recherche dans le texte : causes et conséquences de la Landflucht}

Le texte est construit autour d'une logique de \textbf{cause} (\all{die Ursache, -n}) et de \textbf{conséquence} (\all{die Folge, -n}). Savoir distinguer les deux est une compétence clé pour le résumé et le commentaire.

\begin{definition}[Cause et conséquence]
La \cle{cause} (\all{die Ursache}) est la raison, le « pourquoi » d'un phénomène : on l'exprime avec \all{weil}, \all{denn}, \all{wegen}. La \cle{conséquence} (\all{die Folge}) est le résultat, le « donc » : on l'exprime avec \all{deshalb}, \all{darum}, \all{so dass}.
\end{definition}

\begin{exemplebox}[Deux causes et deux conséquences relevées dans le texte]
\textbf{Causes (Ursachen) de la Landflucht :}
\begin{itemize}
\item \all{Es gibt keine Arbeit auf dem Land.} « Il n'y a pas de travail à la campagne. »
\item \all{Die Familien leben in Armut und es gibt keine guten Schulen.} « Les familles vivent dans la pauvreté et il n'y a pas de bonnes écoles. »
\end{itemize}
\textbf{Conséquences (Folgen) de la Landflucht :}
\begin{itemize}
\item \all{Nur die alten Leute bleiben im Dorf.} « Seules les personnes âgées restent au village. »
\item \all{Die Felder liegen brach und die Schulen schließen.} « Les champs restent en friche et les écoles ferment. »
\end{itemize}
\end{exemplebox}

\begin{popfigure}
\begin{center}
\begin{tikzpicture}
\node[draw=popblue, fill=popblueL, rounded corners, thick, minimum width=5.6cm, minimum height=1cm, align=center, font=\bfseries] (u) at (0,0) {Ursachen (causes)};
\node[draw=popblue, fill=white, rounded corners, minimum width=5.6cm, align=center] (u1) at (0,-1.1) {keine Arbeit auf dem Land};
\node[draw=popblue, fill=white, rounded corners, minimum width=5.6cm, align=center] (u2) at (0,-2.1) {Armut der Familien};
\node[draw=popblue, fill=white, rounded corners, minimum width=5.6cm, align=center] (u3) at (0,-3.1) {keine Schulen und Krankenhäuser};
\node[draw=poporange, fill=poporangeL, rounded corners, thick, minimum width=5.6cm, minimum height=1cm, align=center, font=\bfseries] (f) at (7.4,0) {Folgen (conséquences)};
\node[draw=poporange, fill=white, rounded corners, minimum width=5.6cm, align=center] (f1) at (7.4,-1.1) {nur Alte bleiben im Dorf};
\node[draw=poporange, fill=white, rounded corners, minimum width=5.6cm, align=center] (f2) at (7.4,-2.1) {die Felder liegen brach};
\node[draw=poporange, fill=white, rounded corners, minimum width=5.6cm, align=center] (f3) at (7.4,-3.1) {die Schulen schließen};
\node[draw=popdark, fill=popcream, rounded corners, thick, minimum width=3.2cm, align=center, font=\bfseries] (m) at (3.7,-4.4) {die Landflucht};
\draw[-{Stealth}, very thick, popblue] (u3.south) .. controls (1.5,-3.9) .. (m.west);
\draw[-{Stealth}, very thick, poporange] (m.east) .. controls (5.9,-3.9) .. (f3.south);
\end{tikzpicture}
\end{center}
\legende{Les causes (à gauche) provoquent l'exode rural, qui entraîne des conséquences (à droite).}
\end{popfigure}

\begin{center}
\begin{tabularx}{\linewidth}{|X|X|X|}
\hline
\rowcolor{popblueL} Français & Deutsch & Remarque \\
\hline
la cause & die Ursache, -n & nom féminin \\
\hline
la conséquence & die Folge, -n & nom féminin \\
\hline
à cause de & wegen + génitif/datif & \all{wegen der Armut} \\
\hline
parce que & weil + verbe final & \all{weil sie Arbeit suchen} \\
\hline
c'est pourquoi / donc & deshalb + verbe en 2\textsuperscript{e} position & \all{deshalb ziehen sie um} \\
\hline
rester en friche & brach liegen & verbe à particule séparable : \all{die Felder liegen brach} \\
\hline
fermer (une école) & schließen, schließt, schloss, hat geschlossen & verbe fort \\
\hline
\end{tabularx}
\end{center}

\courssub{Le résumé guidé en trois phrases}

Le résumé (\all{die Zusammenfassung, -en}) consiste à dire l'essentiel du texte avec ses propres mots, en peu de phrases. Ici, trois connecteurs sont imposés : \all{weil} (cause), \all{deshalb} (conséquence), \all{aber} (opposition). Rappel des positions du verbe :

\begin{itemize}
\item \all{weil} + sujet + ... + \textbf{verbe conjugué à la fin} : \all{..., weil es auf dem Land keine Arbeit gibt.}
\item \all{deshalb} en première position + \textbf{verbe conjugué} + sujet : \all{Deshalb ziehen die jungen Leute in die Stadt.}
\item \all{aber} ne change pas l'ordre des mots : \all{Aber nicht alle finden dort Arbeit.}
\end{itemize}

\begin{exoresolu}[Résumé guidé avec connecteurs imposés]
Énoncé. Résume le texte \all{„Warum verlassen die jungen Leute das Dorf ?“} en trois phrases complètes. Utilise obligatoirement les connecteurs \all{weil}, \all{deshalb} et \all{aber}.
\tcblower
Solution détaillée.

\textbf{Phrase 1 (cause avec \all{weil}) :} \all{Viele junge Leute verlassen das Dorf, weil es auf dem Land keine Arbeit und keine guten Schulen gibt.} « Beaucoup de jeunes quittent le village parce qu'il n'y a à la campagne ni travail ni bonnes écoles. »

\textbf{Phrase 2 (conséquence avec \all{deshalb}) :} \all{Deshalb ziehen sie in die großen Städte und suchen dort einen Job.} « C'est pourquoi ils s'installent dans les grandes villes et y cherchent un emploi. »

\textbf{Phrase 3 (opposition avec \all{aber}) :} \all{Aber nicht alle finden Arbeit, und im Dorf bleiben nur die alten Leute zurück.} « Mais tous ne trouvent pas de travail, et au village il ne reste que les personnes âgées. »

Vérification : trois phrases complètes, les trois connecteurs imposés sont utilisés, le verbe est à la fin après \all{weil}, en deuxième position après \all{deshalb} et après \all{aber}. Remarque : \all{zurückbleiben} est un verbe à particule séparable ; à la fin de la phrase, la particule \all{zurück} se détache : \all{..., bleiben nur die alten Leute zurück.}
\end{exoresolu}

\begin{attention}[Erreurs fréquentes des francophones dans le résumé]
\begin{itemize}
\item Après \all{weil}, le verbe doit aller à la fin : écris \all{weil sie Arbeit suchen} et jamais \all{weil sie suchen Arbeit}.
\item Après \all{deshalb} en tête de phrase, le verbe vient tout de suite après, avant le sujet : \all{Deshalb ziehen sie um.} L'inversion est obligatoire.
\item Ne recopie pas de longues phrases du texte : un résumé reformule l'essentiel avec tes propres mots.
\item N'oublie pas la majuscule à tous les noms : \all{das Dorf}, \all{die Stadt}, \all{die Arbeit}, \all{die Armut}.
\item Attention au verbe \all{umziehen} (déménager), séparable : \all{Sie ziehen in die Stadt um.} La particule \all{um} va à la fin de la proposition principale.
\end{itemize}
\end{attention}

\begin{exercice}[À toi de jouer]
\begin{enumerate}
\item Réponds en phrases complètes : \all{Wer bleibt im Dorf ? Was gibt es in der Stadt ? Was müssen die Regierungen tun ?}
\item Richtig oder falsch ? Justifie : \all{„Die Schulen im Dorf schließen, weil viele Kinder in die Stadt ziehen.“}
\item Relève dans le texte une troisième cause et une troisième conséquence de la Landflucht.
\item Résume le texte en trois phrases avec \all{weil}, \all{deshalb} et \all{aber}, sans regarder le corrigé de l'exercice résolu.
\end{enumerate}
\end{exercice}

\begin{corrige}[Exercice « À toi de jouer »]
\textbf{1.} \all{Im Dorf bleiben nur die alten Leute. In der Stadt gibt es Fabriken, Büros, Krankenhäuser, Schulen und Universitäten. Die Regierungen müssen Straßen, Schulen und Krankenhäuser auf dem Land bauen und Arbeit für die Jugendlichen schaffen.}

\textbf{2.} \all{Das ist richtig. Viele Familien ziehen mit ihren Kindern in die Stadt, deshalb schließen die Schulen im Dorf.} « C'est vrai. Beaucoup de familles s'installent en ville avec leurs enfants, c'est pourquoi les écoles du village ferment. »

\textbf{3.} Troisième cause possible : \all{Es gibt keine Krankenhäuser auf dem Land.} Troisième conséquence possible : \all{Manche Jugendliche finden in der Stadt keinen Job und leben in Armut.}

\textbf{4.} Exemple de résumé : \all{Die jungen Leute verlassen das Dorf, weil sie dort keine Arbeit finden. Deshalb ziehen sie in die großen Städte. Aber nicht alle finden einen Job, und das Dorf wird immer älter.}
\end{corrige}

Cette étape de compréhension te donne maintenant toutes les clés du texte : son thème général, ses informations précises, sa logique de cause et de conséquence. Nous pourrons ainsi, dans la partie suivante, enrichir notre lexique thématique autour de \all{das Dorf}, \all{die Stadt} et \all{die Landflucht}.

\coursec{Lexique thématique : Dorf, Stadt, Landflucht}

Après avoir analysé le texte \emph{« Warum verlassen die jungen Leute das Dorf ? »} et répondu aux questions de compréhension, nous allons maintenant structurer et approfondir le vocabulaire essentiel de ce thème. La maîtrise de ce lexique, avec ses genres, ses pluriels et ses familles de mots, est la clé pour réussir l'expression écrite et orale sur l'exode rural.

\courssub{Observation à partir du texte}

Reprenons quelques phrases du texte support pour observer le vocabulaire en contexte :

\begin{textebox}{Extrait du texte support}
Viele junge Leute \textbf{verlassen} ihr \textbf{Dorf}, weil es dort \textbf{keine Arbeit} gibt. Sie \textbf{ziehen} in die \textbf{Stadt} \textbf{um}, um einen \textbf{Job} zu \textbf{finden}. In der \textbf{Großstadt} gibt es \textbf{Fabriken}, \textbf{Büros} und \textbf{Universitäten}. Aber das Leben in der Stadt ist \textbf{teuer} und \textbf{laut}. Auf dem \textbf{Land} ist es \textbf{ruhig}, aber die \textbf{Armut} ist groß. Die \textbf{Landflucht} ist ein großes \textbf{Problem} für die \textbf{Dörfer}.
\end{textebox}

\noindent \textbf{Glossaire des mots difficiles :}
\begin{itemize}
    \item \all{verlassen (verließ, hat verlassen)} : quitter (verbe fort, irrégulier au Präteritum/Perfekt).
    \item \all{das Dorf, die Dörfer} : le village.
    \item \all{die Arbeit, die Arbeiten} : le travail / l'ouvrage.
    \item \all{umziehen (zog um, ist umgezogen)} : déménager (verbe à particule séparable + fort).
    \item \all{die Stadt, die Städte} : la ville.
    \item \all{der Job, die Jobs} : le boulot, l'emploi (terme familier/quotidien).
    \item \all{die Großstadt, die Großstädte} : la grande ville, la métropole.
    \item \all{die Fabrik, die Fabriken} : l'usine.
    \item \all{das Büro, die Büros} : le bureau.
    \item \all{die Universität, die Universitäten} : l'université.
    \item \all{teuer} : cher (coût de la vie).
    \item \all{laut} : bruyant.
    \item \all{auf dem Land} : à la campagne.
    \item \all{ruhig} : calme, tranquille.
    \item \all{die Armut} : la pauvreté (nom féminin, pas de pluriel courant).
    \item \all{die Landflucht} : l'exode rural (litt. : la fuite de la campagne).
    \item \all{das Problem, die Probleme} : le problème.
\end{itemize}

\courssub{Les quatre champs lexicaux principaux}

Nous organisons le vocabulaire en quatre pôles pour faciliter la mémorisation. Chaque nom est donné avec \textbf{son article} et \textbf{son pluriel} (indispensable en allemand).

\begin{popfigure}
\centering
\begin{tikzpicture}[
    mindmap,
    grow cyclic,
    every node/.style={concept, execute at begin node=\hskip0pt},
    concept color=popblue,
    level 1/.append style={level distance=4.5cm, sibling angle=90},
    level 2/.append style={level distance=3cm, sibling angle=45},
    DAS DORF/.style={concept color=popgreen},
    DIE STADT/.style={concept color=poporange},
    DER UMZUG/.style={concept color=poppurple},
    ARBEIT/.style={concept color=poppink}
]
\node[concept] {LEXIK: LANDFLUCHT}
    child [DAS DORF] { node[concept] {DAS DORF}
        child { node[concept] {das Dorf, die Dörfer} }
        child { node[concept] {das Feld, die Felder} }
        child { node[concept] {der Bauer, die Bauern} }
        child { node[concept] {auf dem Land} }
        child { node[concept] {ruhig} }
    }
    child [DIE STADT] { node[concept] {DIE STADT}
        child { node[concept] {die Stadt, die Städte} }
        child { node[concept] {die Großstadt} }
        child { node[concept] {das Krankenhaus} }
        child { node[concept] {die Fabrik} }
        child { node[concept] {teuer, laut} }
    }
    child [DER UMZUG] { node[concept] {DER UMZUG}
        child { node[concept] {umziehen} }
        child { node[concept] {verlassen} }
        child { node[concept] {der Umzug} }
        child { node[concept] {die Hoffnung} }
        child { node[concept] {die Landflucht} }
    }
    child [ARBEIT] { node[concept] {ARBEIT \& ARMUT}
        child { node[concept] {die Arbeit} }
        child { node[concept] {der Job} }
        child { node[concept] {arbeiten} }
        child { node[concept] {arbeitslos} }
        child { node[concept] {die Armut} }
    };
\end{tikzpicture}
\legende{Carte mentale : Les 4 pôles lexicaux de l'exode rural}
\end{popfigure}

\begin{center}
\begin{tabularx}{\linewidth}{|X|X|X|}
\hline
\rowcolor{popblueL}
\textbf{Champ : DAS DORF (Le village / La campagne)} & \textbf{Champ : DIE STADT (La ville)} & \textbf{Champ : DER UMZUG (Le départ / Déménagement)} \\
\hline
\all{das Dorf, die Dörfer} (village) & \all{die Stadt, die Städte} (ville) & \all{die Landflucht} (exode rural) \\
\all{das Feld, die Felder} (champ) & \all{die Großstadt, die Großstädte} (grande ville) & \all{umziehen, zog um, ist umgezogen} (déménager) \\
\all{der Bauer, die Bauern} (paysan) & \all{das Krankenhaus, die Krankenhäuser} (hôpital) & \all{verlassen, verließ, hat verlassen} (quitter) \\
\all{die Farm, die Farmen} (ferme) & \all{die Fabrik, die Fabriken} (usine) & \all{der Umzug, die Umzüge} (déménagement) \\
\all{auf dem Land} (à la campagne) & \all{das Büro, die Büros} (bureau) & \all{die Hoffnung, die Hoffnungen} (espoir) \\
\all{ruhig} (calme) & \all{die Universität, die Universitäten} (université) & \all{in die Stadt ziehen} (aller en ville) \\
& \all{teuer} (cher) & \all{wegziehen} (partir, s'en aller) \\
& \all{laut} (bruyant) & \\
\hline
\end{tabularx}
\end{center}

\begin{center}
\begin{tabularx}{\linewidth}{|X|X|}
\hline
\rowcolor{popblueL}
\textbf{Champ : ARBEIT \& ARMUT (Travail \& Pauvreté)} & \textbf{Familles de mots (Wortfamilien)} \\
\hline
\all{die Arbeit, die Arbeiten} (travail) & \textbf{arbeiten} (verbe) $\rightarrow$ \all{die Arbeit} (nom) $\rightarrow$ \all{der Arbeiter, die Arbeiter} (ouvrier) $\rightarrow$ \all{arbeitslos} (adj., au chômage) \\
\all{arbeiten} (travailler) & \textbf{wohnen} (verbe) $\rightarrow$ \all{die Wohnung, die Wohnungen} (logement) $\rightarrow$ \all{der Wohnort, die Wohnorte} (lieu de résidence) \\
\all{der Job, die Jobs} (emploi, boulot) & \textbf{fliehen} (verbe, fuir) $\rightarrow$ \all{die Flucht, die Fluchten} (fuite) $\rightarrow$ \all{die Landflucht} (exode rural) \\
\all{arbeitslos} (au chômage) & \textbf{arm} (adj., pauvre) $\rightarrow$ \all{die Armut} (pauvreté) $\rightarrow$ \all{der Arme, die Armen} (le pauvre) \\
\all{die Armut} (pauvreté) & \textbf{reich} (adj., riche) $\rightarrow$ \all{der Reichtum} (richesse) $\rightarrow$ \all{der Reiche, die Reichen} (le riche) \\
\all{arm / reich} (pauvre / riche) & \textbf{bauen} (construire/cultiver) $\rightarrow$ \all{der Bauer} (paysan) $\rightarrow$ \all{die Landwirtschaft} (agriculture) \\
\all{in Armut leben} (vivre dans la pauvreté) & \textbf{städten} (urbaniser) $\rightarrow$ \all{die Stadt} $\rightarrow$ \all{städtisch} (urbain) \\
\all{Arbeit suchen / finden} (chercher/trouver du travail) & \\
\hline
\end{tabularx}
\end{center}

\begin{propriete}[Règle d'or : Apprendre le nom avec article et pluriel]
En allemand, \textbf{le genre (genre grammatical) et le pluriel sont imprévisibles}. On ne peut pas les deviner à coup sûr. \textbf{Apprenez systématiquement :} \cle{Article + Nom + Pluriel}.
\begin{itemize}
    \item \all{das Dorf, die Dörfer} (neutre, pluriel en \textbf{-er} + Umlaut)
    \item \all{die Stadt, die Städte} (féminin, pluriel en \textbf{-e} + Umlaut)
    \item \all{das Krankenhaus, die Krankenhäuser} (neutre, composé $\rightarrow$ genre du dernier élément, pluriel \textbf{-er} + Umlaut)
    \item \all{die Fabrik, die Fabriken} (féminin, pluriel en \textbf{-en})
    \item \all{das Büro, die Büros} (neutre, pluriel en \textbf{-s} car mot d'origine latine/étrangère)
    \item \all{die Universität, die Universitäten} (féminin, pluriel en \textbf{-en})
    \item \all{die Armut} (féminin, \textbf{pas de pluriel} usuel, comme « la pauvreté »)
\end{itemize}
\emph{Astuce :} Les noms en \textbf{-ung}, \textbf{-heit}, \textbf{-keit}, \textbf{-schaft}, \textbf{-ion}, \textbf{-tät}, \textbf{-eur} sont \textbf{toujours féminins} (die Landflucht, die Hoffnung, die Arbeit, die Universität).
\end{propriete}

\courssub{Focus sur les verbes clés : \textit{umziehen} et \textit{verlassen}}

\begin{definition}[umziehen – verbe à particule séparable et verbe fort]
\all{umziehen} signifie « déménager, changer de résidence ».
\begin{itemize}
    \item \textbf{Séparable :} La particule \cle{um-} se détache et va à la fin de la proposition principale.
    \item \textbf{Fort (irrégulier) :} Changement de voyelle radicale au Präteritum (\textbf{a} $\rightarrow$ \textbf{og}) et au Participe II (\textbf{a} $\rightarrow$ \textbf{o}).
\end{itemize}
\textbf{Conjugaison (Präsens / Präteritum / Perfekt) :}
\begin{center}
\begin{tabular}{|l|l|l|l|}
\hline
\rowcolor{popblueL}
Personne & Präsens & Präteritum & Perfekt (sein) \\
\hline
ich & ziehe um & zog um & bin umgezogen \\
du & ziehst um & zogst um & bist umgezogen \\
er/sie/es & zieht um & zog um & ist umgezogen \\
wir & ziehen um & zogen um & sind umgezogen \\
ihr & zieht um & zogt um & seid umgezogen \\
sie/Sie & ziehen um & zogen um & sind umgezogen \\
\hline
\end{tabular}
\end{center}
\emph{Remarque :} Verbe de mouvement $\rightarrow$ auxiliaire \textbf{sein} au Perfekt.
\end{definition}

\begin{exemplebox}[Exemples avec \textit{umziehen}]
\all{Er zieht nächstes Jahr nach Yaoundé um.} = Il déménage à Yaoundé l'an prochain. \\
\all{Wir sind in ein neues Haus umgezogen.} = Nous avons déménagé dans une nouvelle maison. \\
\all{Viele junge Leute ziehen in die Großstädte um.} = Beaucoup de jeunes déménagent dans les grandes villes.
\end{exemplebox}

\begin{definition}[verlassen – verbe fort (irrégulier), non séparable]
\all{verlassen} signifie « quitter (un lieu, une personne) ». Il prend un \textbf{accusatif}.
\begin{itemize}
    \item \textbf{Fort :} \textbf{e} $\rightarrow$ \textbf{ie} (Présens du/er), \textbf{a} $\rightarrow$ \textbf{ie} (Präteritum), \textbf{a} $\rightarrow$ \textbf{assen} (Participe II).
    \item \textbf{Non séparable :} Préfixe \cle{ver-} reste collé.
\end{itemize}
\textbf{Conjugaison (Présens / Präteritum / Perfekt) :}
\begin{center}
\begin{tabular}{|l|l|l|l|}
\hline
\rowcolor{popblueL}
Personne & Präsens & Präteritum & Perfekt (haben) \\
\hline
ich & verlasse & verließ & habe verlassen \\
du & verlässt & verließest & hast verlassen \\
er/sie/es & verlässt & verließ & hat verlassen \\
wir & verlassen & verließen & haben verlassen \\
ihr & verlasst & verließt & habt verlassen \\
sie/Sie & verlassen & verließen & haben verlassen \\
\hline
\end{tabular}
\end{center}
\end{definition}

\begin{exemplebox}[Exemples avec \textit{verlassen}]
\all{Er verlässt das Dorf.} = Il quitte le village. (Accusatif : \textbf{das} Dorf) \\
\all{Sie hat ihren Job verlassen.} = Elle a quitté son emploi. \\
\all{Verlasst bitte den Raum!} = Quittez la salle s'il vous plaît !
\end{exemplebox}

\courssub{Expressions figées et collocations utiles}

Pour parler naturellement de l'exode rural, utilisez ces blocs prêts à l'emploi (chunks) :

\begin{center}
\begin{tabularx}{\linewidth}{|l|X|}
\hline
\rowcolor{popblueL}
\textbf{Expression allemande} & \textbf{Traduction / Contexte} \\
\hline
\all{auf dem Land leben} & Vivre à la campagne (Datif : \textbf{dem} Land) \\
\hline
\all{in die Stadt ziehen / umziehen} & Aller vivre en ville / Déménager en ville (Accusatif : \textbf{in die} Stadt) \\
\hline
\all{Arbeit suchen / finden} & Chercher / Trouver du travail \\
\hline
\all{in Armut leben} & Vivre dans la pauvreté \\
\hline
\all{der Landflucht entgegenwirken} & Lutter contre l'exode rural (verbe + Dativ) \\
\hline
\all{die Arbeitslosigkeit bekämpfen} & Lutter contre le chômage \\
\hline
\all{eine Perspektive haben / sehen} & Avoir / Voir une perspective d'avenir \\
\hline
\all{den Anschluss verpassen} & Rater le coche / Perdre le contact (figé) \\
\hline
\end{tabularx}
\end{center}

\courssub{Texte d'entraînement : « Ein Brief aus dem Dorf »}

Lisez ce texte original (niveau A2) qui réutilise le vocabulaire ci-dessus. Répondez ensuite aux questions.

\begin{textebox}{Ein Brief aus dem Dorf}
Liebe Amina,

vielen Dank für deinen Brief. Hier in Bafia ist alles wie immer. Mein Bruder Kofi hat das Dorf \textbf{verlassen}. Er ist vor zwei Wochen nach Douala \textbf{gezogen}. Dort hat er in einer \textbf{Fabrik} einen \textbf{Job} \textbf{gefunden}. Er \textbf{arbeitet} jetzt als Maschinist. Das \textbf{Gehalt} ist besser als auf dem \textbf{Land}, aber die \textbf{Miete} für seine \textbf{Wohnung} ist sehr \textbf{teuer}.

Das \textbf{Leben in der Großstadt} ist \textbf{laut} und \textbf{stressig}. Kofi vermisst die \textbf{Ruhe} hier. Er sagt: „In der Stadt gibt es viele \textbf{Möglichkeiten}, aber man ist oft \textbf{einsam}.“ Unsere Eltern sind \textbf{traurig}, weil der \textbf{Sohn} weg ist. Die \textbf{Landflucht} ist ein großes \textbf{Thema} bei uns. Viele \textbf{Häuser} im Dorf stehen \textbf{leer}. Die \textbf{Bauern} werden \textbf{älter}, und die \textbf{jungen Leute} wollen \textbf{weg}.

Ich bleibe hier. Ich \textbf{helfe} meinem Vater auf dem \textbf{Feld}. Wir bauen \textbf{Kakao} und \textbf{Mais} an. Vielleicht \textbf{gründe} ich später eine \textbf{Genossenschaft} (Kooperative). Das ist meine \textbf{Hoffnung}.

Schreib mir bald!

Dein Freund
Paul
\end{textebox}

\noindent \textbf{Glossaire complémentaire :}
\begin{itemize}
    \item \all{das Gehalt, die Gehälter} : le salaire.
    \item \all{die Miete} : le loyer.
    \item \all{stressig} : stressant.
    \item \all{vermissen} : manquer à qqn (accusatif : \all{Er vermisst die Ruhe}).
    \item \all{die Möglichkeit, die Möglichkeiten} : la possibilité.
    \item \all{einsam} : seul, solitaire.
    \item \all{leer} : vide.
    \item \all{bauen (baute, hat gebaut)} : cultiver, construire.
    \item \all{der Kakao} : le cacao.
    \item \all{gründen} : fonder, créer (entreprise).
    \item \all{die Genossenschaft, die Genossenschaften} : la coopérative.
\end{itemize}

\begin{exoresolu}[Compréhension et lexique : « Ein Brief aus dem Dorf »]
\textbf{Énoncé :} Répondez en allemand (phrases complètes) aux questions suivantes :
\begin{enumerate}
    \item Wohin ist Kofi gezogen?
    \item Was arbeitet er jetzt?
    \item Warum sind die Eltern traurig?
    \item Was macht Paul im Dorf?
    \item Was ist Pauls Hoffnung für die Zukunft?
\end{enumerate}
\tcblower
\textbf{Solution détaillée :}
\begin{enumerate}
    \item \all{Kofi ist nach Douala gezogen.} (Präteritum de \textit{ziehen}, mouvement $\rightarrow$ sein). \emph{Il a déménagé à Douala.}
    \item \all{Er arbeitet jetzt als Maschinist in einer Fabrik.} (Présens + préposition \textit{als} pour la fonction). \emph{Il travaille maintenant comme machiniste dans une usine.}
    \item \all{Die Eltern sind traurig, weil der Sohn weg ist / das Dorf verlassen hat.} (Subordonnée \textit{weil} + verbe à la fin). \emph{Les parents sont tristes car le fils est parti / a quitté le village.}
    \item \all{Paul hilft seinem Vater auf dem Feld. Er baut Kakao und Mais an.} (Datif \textit{seinem Vater}, séparable \textit{anbauen}). \emph{Paul aide son père au champ. Il cultive du cacao et du maïs.}
    \item \all{Pauls Hoffnung ist, später eine Genossenschaft zu gründen.} (Nominalisation / Infinitif avec \textit{zu}). \emph{L'espoir de Paul est de fonder plus tard une coopérative.}
\end{enumerate}
\end{exoresolu}

\courssub{Exercices d'application}

\begin{exercice}[Mots croisés lexicaux / Complétion]
Complétez les phrases avec le mot correct du champ lexical (forme correcte : genre, cas, nombre, temps verbal).
\begin{enumerate}
    \item Viele junge Leute \_\_\_\_\_\_\_\_\_\_ (verlassen) ihre Heimat, weil sie \_\_\_\_\_\_\_\_\_\_ (Arbeit) suchen.
    \item Mein Onkel ist vor einem Jahr nach Berlin \_\_\_\_\_\_\_\_\_\_ (umziehen). Er \_\_\_\_\_\_\_\_\_\_ (finden) dort einen guten \_\_\_\_\_\_\_\_\_\_ (Job).
    \item Auf dem \_\_\_\_\_\_\_\_\_\_ (Land) ist es \_\_\_\_\_\_\_\_\_\_ (ruhig), aber in der \_\_\_\_\_\_\_\_\_\_ (Stadt) ist es \_\_\_\_\_\_\_\_\_\_ (laut) und \_\_\_\_\_\_\_\_\_\_ (teuer).
    \item Die \_\_\_\_\_\_\_\_\_\_ (Landflucht) führt dazu, dass die \_\_\_\_\_\_\_\_\_\_ (Dorf) \_\_\_\_\_\_\_\_\_\_ (leer) stehen.
    \item Mein Vater ist \_\_\_\_\_\_\_\_\_\_ (Bauer). Er \_\_\_\_\_\_\_\_\_\_ (arbeiten) auf dem \_\_\_\_\_\_\_\_\_\_ (Feld).
    \item Wer \_\_\_\_\_\_\_\_\_\_ (arm) ist, hat oft keine \_\_\_\_\_\_\_\_\_\_ (Perspektive) auf dem Land.
\end{enumerate}
\end{exercice}

\begin{corrige}[Exercice 1]
\begin{enumerate}
    \item \all{verlassen} (Présens 3e pl.) / \all{Arbeit} (Acc. sg. fém.).
    \item \all{gezogen} (Part. II de \textit{ziehen} pour Perfekt avec \textit{ist}) / \all{hat ... gefunden} (Perfekt) / \all{Job} (Acc. sg. masc.).
    \item \all{Land} (Datif sg. n. après \textit{auf dem}) / \all{ruhig} / \all{Stadt} (Datif sg. f. après \textit{in der}) / \all{laut} / \all{teuer}.
    \item \all{Landflucht} (Nom. sg. f.) / \all{Dörfer} (Pluriel de \textit{Dorf} : Acc. pl. ou Nom. pl.) / \all{leer} (adj. attribut).
    \item \all{Bauer} (Nom. sg. masc.) / \all{arbeitet} (Présens 3e sg.) / \all{Feld} (Datif sg. n. après \textit{auf dem}).
    \item \all{arm} (adj. attribut après \textit{ist}) / \all{Perspektive} (Acc. sg. f. après \textit{keine}).
\end{enumerate}
\end{corrige}

\begin{exercice}[Traduction thème : Français $\rightarrow$ Allemand]
Traduisez en allemand en utilisant le vocabulaire de la leçon.
\begin{enumerate}
    \item Le village est calme, mais il n'y a pas de travail. (\textit{il y a} = \all{es gibt} + Accusatif)
    \item Ma sœur déménage en ville la semaine prochaine. (\textit{déménager} = \all{umziehen})
    \item Elle veut trouver un emploi dans un bureau.
    \item La vie dans la grande ville est chère.
    \item L'exode rural est un problème pour l'agriculture.
    \item Les paysans sont pauvres, mais les usines sont riches. (pluriels !)
\end{enumerate}
\end{exercice}

\begin{corrige}[Exercice 2]
\begin{enumerate}
    \item \all{Das Dorf ist ruhig, aber es gibt keine Arbeit.}
    \item \all{Meine Schwester zieht nächste Woche in die Stadt um.} (Séparable : \textit{um} à la fin).
    \item \all{Sie will einen Job in einem Büro finden.} (Acc. masc. \textit{einen Job}, Datif n. \textit{einem Büro} après \textit{in} pour lieu statique du job).
    \item \all{Das Leben in der Großstadt ist teuer.}
    \item \all{Die Landflucht ist ein Problem für die Landwirtschaft.}
    \item \all{Die Bauern sind arm, aber die Fabriken sind reich.} (Pluriels : \textit{die Bauern}, \textit{die Fabriken}).
\end{enumerate}
\end{corrige}

\courssub{Attention : Pièges fréquents pour francophones}

\begin{attention}[Faux amis et erreurs de traduction]
\begin{enumerate}
    \item \textbf{\all{bekommen} $\neq$ devenir.} \all{bekommen} = \textbf{recevoir, obtenir}.
    \begin{itemize}
        \item \all{Er bekommt einen Job.} = Il obtient un emploi. (Correct).
        \item \all{Er wird Arzt.} = Il devient médecin. (Pour « devenir » : \all{werden}).
    \end{itemize}
    \item \textbf{\all{Arbeit finden} vs \all{Arbeit bekommen}.}
    \begin{itemize}
        \item \all{Arbeit suchen / finden} : Chercher / Trouver du travail (action active).
        \item \all{Arbeit bekommen} : Obtenir un emploi (souvent après candidature, insister sur le résultat « recevoir »).
        \item \textbf{Piège :} Ne dites pas \all{Ich bekomme Arbeit} pour « Je trouve du travail » dans le sens général. Préférez \all{Ich finde Arbeit} ou \all{Ich suche Arbeit}.
    \end{itemize}
    \item \textbf{\all{umziehen} vs \all{ziehen}.}
    \begin{itemize}
        \item \all{umziehen} = déménager (changer de domicile). Particule \textbf{um-} à la fin : \all{Er zieht um.}
        \item \all{ziehen} = tirer, traîner ; ou aller/partir (sans changer de domicile officiellement) : \all{Wir ziehen in den Park.} (On va au parc).
    \end{itemize}
    \item \textbf{Majuscules aux noms !} Toujours : \all{das Dorf, die Stadt, die Arbeit, die Armut, die Landflucht}. Jamais en minuscule.
    \item \textbf{Prépositions \textit{in} / \textit{auf} + Accusatif (mouvement) vs Datif (lieu).}
    \begin{itemize}
        \item Mouvement vers ville/pays (frontière) : \all{in die Stadt ziehen} (Acc.), \all{ins Dorf gehen} (Acc. \textit{in das}).
        \item Mouvement vers « Land » (campagne) : \all{aufs Land ziehen} (\textit{auf das} $\rightarrow$ \textbf{aufs}, Acc.).
        \item Lieu (où ?) : \all{in der Stadt} (Datif), \all{auf dem Land} (Datif).
    \end{itemize}
\end{enumerate}
\end{attention}

\begin{savaistu}[« Stadtluft macht frei » – L'air de la ville rend libre]
Ce proverbe médiéval (\textbf{Stadtrecht}) stipulait qu'un serf qui parvenait à vivre « un an et un jour » (\all{Year und Tag}) dans une ville libre (\all{Freie Reichsstadt}) sans être réclamé par son seigneur, devenait libre. L'air de la ville symbolisait l'émancipation juridique et sociale. Aujourd'hui, l'expression est utilisée de façon ironique ou critique pour décrire l'anonymat urbain, mais aussi les opportunités (études, carrière) que la ville offre par rapport aux structures traditionnelles du village. Au Cameroun, l'exode rural vers Yaoundé ou Douala suit souvent cette même logique d'espoir de « liberté » économique.
\end{savaistu}

\courssub{Synthèse pour l'expression écrite (Redemittel)}

Pour rédiger un paragraphe sur l'exode rural, utilisez cette structure :

\begin{methode}[Rédiger un paragraphe sur « Die Landflucht »]
\begin{enumerate}
    \item \textbf{Constat (Ursache) :} \all{Weil es auf dem Land wenig Arbeit gibt, ...} / \all{Da die Armut auf dem Land groß ist, ...}
    \item \textbf{Action (Handlung) :} \all{... verlassen viele junge Leute das Dorf.} / \all{... ziehen sie in die Stadt um.}
    \item \textbf{But (Ziel) :} \all{Sie hoffen, in der Fabrik oder im Büro einen Job zu finden.}
    \item \textbf{Réalité (Folge) :} \all{Aber das Leben in der Großstadt ist teuer und laut.} / \all{Oft werden sie arbeitslos.}
    \item \textbf{Opinion/Solution (Meinung/Lösung) :} \all{Meiner Meinung nach muss man die Landwirtschaft modernisieren, damit die Leute bleiben.}
\end{enumerate}
\end{methode}

\begin{exemplebox}[Modèle de production écrite (niveau A2)]
\all{„Viele junge Leute verlassen das Dorf, weil es dort keine Arbeit gibt. Sie ziehen in die Großstadt um, um Geld zu verdienen. Aber in der Stadt ist das Leben teuer und die Mieten sind hoch. Manche werden arbeitslos und leben in Armut. Die Landflucht schwächt die Dörfer. Ich finde, der Staat muss Fabriken auf dem Land bauen, damit die Jugend eine Perspektive hat.“}
\end{exemplebox}

\coursec{Grammaire : la cause et la conséquence}

Après avoir exploré le lexique de l'exode rural (\all{das Dorf, die Stadt, die Landflucht, die Armut, umziehen, die Arbeit}) dans la section précédente, nous disposons maintenant du vocabulaire nécessaire pour comprendre \textbf{comment} la langue allemande articule les raisons (causes) et les résultats (conséquences) de ce phénomène. Le texte \emph{„Warum verlassen die jungen Leute das Dorf ?“} que nous venons d'étudier regorge de ces structures logiques. Les maîtriser, c'est passer de la compréhension passive à l'expression active et correcte — une compétence exigée à l'examen, aussi bien en compréhension qu'en production écrite.

\courssub{Observation dans le texte support}

Reprenons quelques phrases extraites du texte \emph{„Warum verlassen die jungen Leute das Dorf ?“} (section 2). Identifiez les mots en \textbf{gras} qui relient deux idées : une cause et une conséquence.

\begin{textebox}{Extraits du texte „Warum verlassen die jungen Leute das Dorf ?“}
\textbf{1.} Viele junge Leute verlassen das Dorf, \textbf{weil} es dort keine Arbeitsplätze gibt.\\
\textbf{2.} \textbf{Weil} die Landwirtschaft wenig Ertrag bringt, \textbf{ziehen} viele Familien in die Stadt.\\
\textbf{3.} Auf dem Land gibt es wenige Schulen, \textbf{deshalb} müssen die Kinder lange Wege gehen.\\
\textbf{4.} \textbf{Wegen} der Armut können die Eltern die Schulgebühren nicht zahlen.\\
\textbf{5.} Die Situation ist hoffnungslos, \textbf{so dass} die Jugend keine Zukunft sieht.\\
\textbf{6.} In der Stadt gibt es Fabriken, \textbf{denn} die Industrie braucht Arbeitskräfte.
\end{textebox}

\noindent
\textbf{Analyse rapide :}
\begin{itemize}
    \item Phrases 1 et 2 : \cle{weil} introduit la \textbf{cause} (la raison). Le verbe conjugué de la proposition introduite par \all{weil} se trouve à la \textbf{fin} (\all{gibt, bringt}). Dans la phrase 2, \all{weil} est en tête de phrase $\rightarrow$ inversion dans la principale (\all{ziehen viele Familien}).
    \item Phrase 3 : \cle{deshalb} introduit la \textbf{conséquence} (le résultat). Le verbe reste en \textbf{position 2} (\all{müssen}).
    \item Phrase 4 : \cle{wegen} + nom (\all{der Armut}) exprime la cause sans verbe (cause nominale). Cas : génitif (formel) ou datif (courant).
    \item Phrase 5 : \cle{so dass} introduit une conséquence (souvent avec une nuance d'intensité ou de résultat). Verbe à la fin (\all{sieht}).
    \item Phrase 6 : \cle{denn} relie deux phrases principales (coordination). Verbe en position 2 (\all{braucht}). Pas de verbe à la fin.
\end{itemize}

\courssub{La subordonnée de cause avec \textit{weil} (verbe à la fin)}

C'est la structure la plus fréquente pour dire « parce que ». La proposition subordonnée (\all{Nebensatz}) commence par \all{weil} et \textbf{le verbe conjugué est rejeté à la toute fin de cette proposition}.

\begin{propriete}[Règle fondamentale : \textit{weil} + verbe à la fin]
Dans une phrase introduite par \textbf{\textit{weil}}, le verbe conjugué occupe la \textbf{dernière place} de la subordonnée.
\begin{center}
\begin{tabular}{|l|l|l|}
\hline
\rowcolor{popblueL} \textbf{Hauptsatz (Principale)} & \textbf{Konjunktion} & \textbf{Nebensatz (Subordonnée weil)} \\ \hline
Sie verlassen das Dorf & , weil & es dort keine Arbeit \textbf{gibt}. \\ \hline
\end{tabular}
\end{center}
\all{Sie verlassen das Dorf, weil es dort keine Arbeit gibt.} « Ils quittent le village parce qu'il n'y a pas de travail là-bas. »
\end{propriete}

\begin{exemplebox}[Exemples variés avec \textit{weil}]
\begin{enumerate}
    \item \all{Die Bauern verkaufen ihre Ernte, weil die Preise niedrig sind.} (Les paysans vendent leur récolte parce que les prix sont bas.)
    \item \all{Wir lernen Deutsch, weil wir in Deutschland studieren wollen.} (Nous apprenons l'allemand parce que nous voulons étudier en Allemagne.)
    \item \all{Mein Bruder bleibt zu Hause, weil er krank ist.} (Mon frère reste à la maison parce qu'il est malade.)
\end{enumerate}
\end{exemplebox}

\courssub{\textit{Weil} en tête de phrase : l'inversion du sujet}

Lorsque la subordonnée en \all{weil} est placée \textbf{au début de la phrase} (pour insister sur la cause), la phrase principale qui suit commence par le \textbf{verbe conjugué} (inversion verbe-sujet). C'est la règle générale de l'allemand : \textbf{la subordonnée initiale occupe la position 1 de la phrase entière}, le verbe de la principale vient donc en position 2, juste après la virgule.

\begin{propriete}[Inversion après un \textit{Weil}-Satz initial]
Structure : \textbf{[Weil-Satz ,] + Verbe + Sujet + ...}
\begin{center}
\begin{tabular}{|l|l|l|l|}
\hline
\rowcolor{popblueL} \textbf{Position 1 (Weil-Satz)} & \textbf{Position 2 (Verbe)} & \textbf{Sujet} & \textbf{Reste} \\ \hline
\textit{Weil es keine Arbeit gibt,} & \textbf{verlassen} & \textbf{sie} & \textit{das Dorf.} \\ \hline
\end{tabular}
\end{center}
\all{Weil es keine Arbeit gibt, verlassen sie das Dorf.} « Parce qu'il n'y a pas de travail, ils quittent le village. »
\end{propriete}

\begin{popfigure}
\begin{tikzpicture}[
    node distance=0.8cm and 0.9cm,
    box/.style={draw=popblue, thick, rounded corners, fill=popblueL, minimum height=1cm, align=center, font=\small},
    arrow/.style={-Stealth, thick, popdark}
]
% Structure 1: Hauptsatz + weil + Nebensatz
\node[box, minimum width=3.5cm, align=center] (hs){Hauptsatz\\ \textbf{Sie verlassen} das Dorf};
\node[box, right=of hs, minimum width=0.8cm] (komma) {,};
\node[box, right=of komma, minimum width=1.4cm, fill=poporangeL, draw=poporange] (konj) {\textit{weil}};
\node[box, right=of konj, minimum width=4cm, fill=popgreenL, draw=popgreen, align=center] (ns){Nebensatz\\ es dort keine Arbeit \textbf{gibt}};

\draw[arrow] (hs) -- (komma);
\draw[arrow] (komma) -- (konj);
\draw[arrow] (konj) -- (ns);

\draw[popgreen, very thick, decorate, decoration={snake, amplitude=1pt, segment length=5pt}] (ns.south) -- ++(0,-0.6) node[below, font=\tiny, popgreen] {Verbe conjugué à la FIN};

% Structure 2: Weil-Satz + Hauptsatz (Inversion)
\node[box, minimum width=4cm, fill=popgreenL, draw=popgreen, below=2.6cm of hs.south west, anchor=north west, align=center] (ws){Weil-Satz (Pos 1)\\ \textit{Weil} es keine Arbeit \textbf{gibt}};
\node[box, right=of ws, minimum width=0.8cm] (komma2) {,};
\node[box, right=of komma2, minimum width=2.6cm, fill=poppinkL, draw=poppink] (v2) {Verbe \textbf{verlassen} (Pos 2)};
\node[box, right=of v2, minimum width=1.8cm] (subj) {Sujet \textbf{sie}};
\node[box, right=of subj, minimum width=2.2cm] (rest) {Rest \textit{das Dorf}};

\draw[arrow] (ws) -- (komma2);
\draw[arrow] (komma2) -- (v2);
\draw[arrow] (v2) -- (subj);
\draw[arrow] (subj) -- (rest);

\draw[poppink, very thick, decorate, decoration={snake, amplitude=1pt, segment length=5pt}] (v2.south) -- ++(0,-0.6) node[below, font=\tiny, poppink] {Verbe en POSITION 2};
\end{tikzpicture}
\legende{La place du verbe : en haut, structure classique (Hauptsatz + weil-Satz, verbe à la fin) ; en bas, structure avec inversion (Weil-Satz en tête, verbe de la principale en position 2).}
\end{popfigure}

\courssub{L'alternative coordonnée : \textit{denn} (verbe en position 2)}

\all{Denn} signifie aussi « car / parce que », mais c'est une \textbf{conjonction de coordination} (comme \all{und, aber, oder}). Elle relie deux \textbf{phrases principales}. Le verbe reste donc en \textbf{position 2} dans la seconde phrase. \all{Denn} ne déclenche jamais l'inversion ni le renvoi du verbe à la fin.

\begin{propriete}[\textit{denn} = coordination (verbe en position 2)]
\all{Denn} relie deux propositions principales. L'ordre des mots est celui d'une phrase affirmative normale : \textbf{Sujet + Verbe (pos 2) + ...}
\begin{center}
\begin{tabular}{|l|l|l|l|}
\hline
\rowcolor{popblueL} \textbf{Hauptsatz 1} & \textbf{,} & \textbf{denn} & \textbf{Hauptsatz 2 (Sujet + Verbe pos 2)} \\ \hline
Sie gehen in die Stadt & , & \textit{denn} & dort \textbf{gibt} es Arbeit. \\ \hline
\end{tabular}
\end{center}
\all{Sie gehen in die Stadt, denn dort gibt es Arbeit.} « Ils vont en ville, car il y a du travail là-bas. »
\end{propriete}

\begin{attention}[Piège majeur : \textit{weil} vs \textit{denn}]
Ne confondez pas la place du verbe !
\begin{itemize}
    \item \textbf{\textit{weil}} $\rightarrow$ subordonnée $\rightarrow$ \textbf{verbe À LA FIN} : \all{..., weil es dort Arbeit \textbf{gibt}.}
    \item \textbf{\textit{denn}} $\rightarrow$ coordination $\rightarrow$ \textbf{verbe en POSITION 2} : \all{..., denn dort \textbf{gibt} es Arbeit.}
\end{itemize}
\textbf{Erreur fréquente (calque du français) :} \all{*..., weil es gibt Arbeit.} FAUX ! $\rightarrow$ Correct : \all{..., weil es Arbeit \textbf{gibt}.}
\end{attention}

\courssub{La cause nominale : \textit{wegen} + Génitif / Datif}

Quand la cause est exprimée par un \textbf{nom} (et non par une phrase avec verbe), on utilise la préposition \textbf{\textit{wegen}} (+ génitif à l'écrit formel, + datif à l'oral courant). En français : « à cause de », « en raison de ».

\begin{propriete}[\textit{wegen} + cas (Génitif / Datif)]
\begin{itemize}
    \item \textbf{Écrit / formel :} \all{wegen} + \textbf{Génitif} (\all{des / der / des / der} + nom).
    \item \textbf{Parlé / courant :} \all{wegen} + \textbf{Datif} (\all{dem / der / dem / den} + nom).
\end{itemize}
\begin{center}
\begin{tabular}{|l|l|l|}
\hline
\rowcolor{popblueL} \textbf{Contexte} & \textbf{Allemand} & \textbf{Français} \\ \hline
Génitif (écrit) & \all{Wegen \textbf{der Armut} verlassen viele das Dorf.} & À cause de la pauvreté, beaucoup quittent le village. \\ \hline
Datif (parlé) & \all{Wegen \textbf{dem Regen} bleibe ich zu Hause.} & À cause de la pluie, je reste à la maison. \\ \hline
Génitif (masc. sg.) & \all{Wegen \textbf{des Regens} bleibe ich zu Hause.} & À cause de la pluie, je reste à la maison. \\ \hline
Génitif (fém. sg.) & \all{Wegen \textbf{der schlechten Ernte} haben die Bauern kein Geld.} & À cause de la mauvaise récolte, les paysans n'ont pas d'argent. \\ \hline
\end{tabular}
\end{center}
\end{propriete}

\begin{aretenir}[Génitif ou Datif après \textit{wegen} ?]
\begin{itemize}
    \item Noms masculins et neutres au singulier : le génitif se voit grâce à la marque \textbf{-s / -es} : \all{wegen des Regens, wegen des Vaters}.
    \item Noms féminins : génitif et datif sont identiques (\all{wegen der Armut} = les deux cas possibles).
    \item Pluriel : \all{wegen der Eltern} (génitif) ou \all{wegen den Eltern} (datif, courant).
    \item \textbf{Conseil examen :} à l'écrit (copie, rédaction), utilisez le \textbf{génitif} : c'est la norme scolaire attendue.
\end{itemize}
\end{aretenir}

\courssub{La conséquence : \textit{deshalb}, \textit{darum}, \textit{also} (verbe en position 2)}

Pour exprimer le résultat, la conséquence (« c'est pourquoi », « donc », « par conséquent »), on utilise des \textbf{adverbes de liaison} : \textbf{\textit{deshalb}}, \textbf{\textit{darum}}, \textbf{\textit{also}}. Placés en tête de phrase, ils occupent la \textbf{position 1} $\rightarrow$ le verbe conjugué suit immédiatement en \textbf{position 2}, puis vient le sujet (inversion).

\begin{propriete}[Conséquence : \textit{deshalb / darum / also} + verbe en position 2]
Ces mots ne sont pas des conjonctions de subordination : la phrase reste une principale, avec inversion si l'adverbe est en tête.
\begin{center}
\begin{tabular}{|l|l|l|l|}
\hline
\rowcolor{popblueL} \textbf{Phrase 1 (cause)} & \textbf{Liaison (Pos 1)} & \textbf{Verbe (Pos 2)} & \textbf{Sujet + Reste} \\ \hline
Es gibt keine Arbeit, & \textit{deshalb} & \textbf{ziehen} & \textit{sie um.} \\ \hline
Die Schulen sind weit, & \textit{darum} & \textbf{bleiben} & \textit{die Kinder zu Hause.} \\ \hline
Ich bin müde, & \textit{also} & \textbf{gehe} & \textit{ich jetzt.} \\ \hline
\end{tabular}
\end{center}
\all{Es gibt keine Arbeit, deshalb ziehen sie um.} « Il n'y a pas de travail, c'est pourquoi ils déménagent. »
\end{propriete}

\begin{exemplebox}[Nuances entre \textit{deshalb}, \textit{darum}, \textit{also}]
\begin{itemize}
    \item \textbf{\textit{deshalb}} : neutre, standard, écrit et oral. « C'est pourquoi / par conséquent. »
    \item \textbf{\textit{darum}} : un peu plus oral, insiste sur le lien logique direct. « Pour cette raison. »
    \item \textbf{\textit{also}} : très oral, marque une conclusion immédiate, souvent en début de phrase. « Donc / alors. » \all{Also fahren wir jetzt.} « Alors nous partons maintenant. »
\end{itemize}
\end{exemplebox}

\courssub{La conséquence avec \textit{so dass} (verbe à la fin)}

\all{So dass} introduit une \textbf{subordonnée de conséquence} (consécutive). Elle exprime souvent un résultat intense, une conséquence inévitable ou un degré élevé (« si bien que », « de sorte que »). Comme dans toute subordonnée, \textbf{le verbe conjugué va à la fin}.

\begin{propriete}[\textit{so dass} + verbe à la fin]
Structure : \textbf{Hauptsatz , so dass + ... + Verbe (fin)}.\\
\all{Die Armut ist groß, so dass viele weggehen.} « La pauvreté est grande, si bien que beaucoup partent. »\\
\all{Es regnet stark, so dass die Straße überflutet ist.} « Il pleut fort, de sorte que la rue est inondée. »
\end{propriete}

\begin{attention}[\textit{so dass} ou \textit{dass} ?]
Ne confondez pas :
\begin{itemize}
    \item \textbf{\textit{so dass}} = conséquence (si bien que) : \all{Er ist müde, so dass er schläft.} « Il est fatigué, si bien qu'il dort. »
    \item \textbf{\textit{dass}} = complétive (que) : \all{Er sagt, dass er müde ist.} « Il dit qu'il est fatigué. »
\end{itemize}
Dans les deux cas, le verbe conjugué va à la fin ; c'est le \textbf{sens} qui change.
\end{attention}

\courssub{Tableau récapitulatif comparatif : Français / Allemand}

Ce tableau synthétise les correspondances et les pièges de place du verbe pour un francophone.

\begin{center}
\begin{tabularx}{\linewidth}{|X|X|X|X|}
\hline
\rowcolor{popblueL} \textbf{Fonction} & \textbf{Français} & \textbf{Allemand (connecteur)} & \textbf{Place du verbe all.} \\ \hline
Cause (avec verbe) & parce que / comme / puisque & \textbf{\textit{weil}} & \textbf{FIN} (subordonnée) \\ \hline
Cause (avec verbe, coordination) & car & \textbf{\textit{denn}} & \textbf{Pos 2} (principale) \\ \hline
Cause (nom seul) & à cause de / en raison de & \textbf{\textit{wegen}} + Gén./Dat. & pas de verbe (groupe nominal) \\ \hline
Conséquence (adverbe) & c'est pourquoi / donc / par conséquent & \textbf{\textit{deshalb / darum / also}} & \textbf{Pos 2} (principale, inversion) \\ \hline
Conséquence (subordonnée) & si bien que / de sorte que & \textbf{\textit{so dass}} & \textbf{FIN} (subordonnée) \\ \hline
\end{tabularx}
\end{center}

\courssub{Méthode : choisir le bon connecteur (cause / conséquence)}

\begin{methode}[Comment choisir le bon mot-outil ?]
\begin{enumerate}
    \item \textbf{Analysez la nature de la cause :}
    \begin{itemize}
        \item C'est une \textbf{phrase complète avec un verbe conjugué} $\rightarrow$ \all{weil} (standard, verbe à la fin) ou \all{denn} (coordination, verbe en position 2).
        \item C'est un \textbf{nom seul} (la pauvreté, la pluie, le manque d'écoles) $\rightarrow$ \textbf{\textit{wegen}} + génitif (écrit) ou datif (oral).
    \end{itemize}
    \item \textbf{Analysez la nature de la conséquence :}
    \begin{itemize}
        \item Vous voulez une \textbf{phrase principale} (verbe en position 2, inversion) $\rightarrow$ \all{deshalb / darum / also}.
        \item Vous voulez une \textbf{subordonnée} (verbe à la fin, pour insister sur l'intensité du résultat) $\rightarrow$ \all{so dass}.
    \end{itemize}
    \item \textbf{Vérifiez la place du verbe :} appliquez la règle stricte du connecteur choisi (fin pour \all{weil} / \all{so dass} ; position 2 pour \all{denn} / \all{deshalb} / \all{also}).
    \item \textbf{Relisez votre phrase à voix haute} et repérez le verbe conjugué : est-il à la place attendue ?
\end{enumerate}
\end{methode}

\courssub{Exercice résolu : transformation et choix du connecteur}

\begin{exoresolu}[Transformer les phrases en utilisant le connecteur indiqué]
\textbf{Consigne :} reliez les deux phrases en utilisant le mot entre parenthèses. Attention à l'ordre des mots !

\textbf{1.} Es gibt keine Busse. Die Schüler laufen zur Schule. (\textit{weil})\\
\textbf{2.} Die Ernte war schlecht. Die Bauern haben kein Geld. (\textit{deshalb})\\
\textbf{3.} Der Boden ist trocken. Es hat nicht geregnet. (\textit{denn})\\
\textbf{4.} Die Armut. Viele Familien ziehen weg. (\textit{wegen})\\
\textbf{5.} Die Wege sind weit. Die Kinder kommen zu spät. (\textit{so dass})

\tcblower
\textbf{Solution détaillée :}

\textbf{1.} \all{Die Schüler laufen zur Schule, weil es keine Busse gibt.}\\
\textit{Analyse :} cause avec verbe $\rightarrow$ \all{weil} $\rightarrow$ verbe \textbf{gibt} à la fin. (Autre ordre possible : \all{Weil es keine Busse gibt, laufen die Schüler zur Schule.} avec inversion dans la principale.)

\textbf{2.} \all{Die Ernte war schlecht, deshalb haben die Bauern kein Geld.}\\
\textit{Analyse :} conséquence $\rightarrow$ \all{deshalb} en position 1 $\rightarrow$ verbe \textbf{haben} en position 2 $\rightarrow$ sujet \all{die Bauern} (inversion).

\textbf{3.} \all{Der Boden ist trocken, denn es hat nicht geregnet.}\\
\textit{Analyse :} cause avec verbe, coordination $\rightarrow$ \all{denn} $\rightarrow$ ordre normal d'une principale : sujet \all{es} + verbe \textbf{hat} en position 2. Pas de verbe à la fin !

\textbf{4.} \all{Wegen der Armut ziehen viele Familien weg.}\\
\textit{Analyse :} cause nominale $\rightarrow$ \all{wegen} + génitif (\all{der Armut} ; pour un nom féminin, génitif et datif sont identiques). Pas de verbe conjugué après \all{wegen}. La principale suit avec inversion : verbe \textbf{ziehen} en position 2. Attention au verbe à particule : \all{wegziehen} $\rightarrow$ \all{ziehen ... weg}.

\textbf{5.} \all{Die Wege sind weit, so dass die Kinder zu spät kommen.}\\
\textit{Analyse :} conséquence subordonnée $\rightarrow$ \all{so dass} $\rightarrow$ verbe \textbf{kommen} à la fin.
\end{exoresolu}

\courssub{Attention : les erreurs classiques des francophones}

\begin{attention}[Top 3 des erreurs à ne plus commettre]
\begin{enumerate}
    \item \textbf{Le calque « parce que » $\rightarrow$ *\textit{weil} + verbe en position 2* :}\\
    \all{*Er geht weg, weil er hat keine Arbeit.} \textcolor{poppink}{\textbf{FAUX}}\\
    \all{Er geht weg, weil er keine Arbeit \textbf{hat}.} \textcolor{popgreen}{\textbf{RICHTIG}} (verbe \all{hat} à la fin).
    \item \textbf{La confusion \textit{deshalb} / \textit{weil} (place du verbe) :}\\
    \all{*Es regnet, deshalb er \textbf{bleibt} zu Hause.} \textcolor{poppink}{\textbf{FAUX}} (verbe en 3\up{e} position)\\
    \all{Es regnet, \textbf{deshalb bleibt} er zu Hause.} \textcolor{popgreen}{\textbf{RICHTIG}} (\all{deshalb} = position 1, \all{bleibt} = position 2).
    \item \textbf{\textit{Wegen} + accusatif (calque de « à cause de » + nom) :}\\
    \all{*Wegen \textbf{den} Regen...} \textcolor{poppink}{\textbf{FAUX}} (accusatif)\\
    \all{Wegen \textbf{des} Regens...} \textcolor{popgreen}{\textbf{RICHTIG}} (génitif, écrit) ou \all{Wegen \textbf{dem} Regen...} \textcolor{poporange}{\textbf{DATIF}} (parlé).
\end{enumerate}
\end{attention}

\courssub{Entraînement : texte d'application « Die Landflucht in Kamerun »}

Le texte suivant réutilise le vocabulaire de la section 4 et les structures grammaticales de cette section. Lisez-le, repérez les connecteurs logiques (\all{weil, deshalb, wegen, so dass, denn, also, darum}) et répondez aux questions.

\begin{textebox}{Die Landflucht in Kamerun — ein aktuelles Problem}
In Kamerun verlassen viele junge Menschen ihre Dörfer auf dem Land. \textbf{Wegen} der Armut und der schlechten Infrastruktur sehen sie keine Zukunft in der Landwirtschaft. \textbf{Weil} es in den Dörfern keine Fabriken und keine großen Unternehmen gibt, finden sie dort keine Arbeit. \textbf{Denn} die Regierung investiert wenig in die ländlichen Gebiete. Die Straßen sind oft schlecht, \textbf{so dass} der Transport der Ernte in die Stadt teuer und schwierig ist. \textbf{Deshalb} entscheiden sich die Jugendlichen für den Umzug nach Yaoundé oder Douala. \textbf{Also} hoffen sie auf ein besseres Einkommen in der Stadt. Aber oft ist die Enttäuschung groß: \textbf{weil} sie keine Ausbildung haben, \textbf{arbeiten} sie nur als Gelegenheitsarbeiter oder Moto-Taxi-Fahrer (\textit{Bendskin-Fahrer}). \textbf{Darum} bleibt die Armut oft bestehen, nur an einem anderen Ort.
\end{textebox}

\begin{definition}[Glossaire du texte]
\begin{itemize}
    \item \textbf{die Landflucht} (nur Singular) : l'exode rural.
    \item \textbf{die Infrastruktur} (die \textit{-en}) : les infrastructures (routes, électricité, eau).
    \item \textbf{investieren} (investiert, hat investiert) : investir.
    \item \textbf{der Transport} (die \textit{-e}) : le transport.
    \item \textbf{sich entscheiden für + Akk.} (entscheidet sich, entschied sich, hat sich entschieden) : décider de, opter pour.
    \item \textbf{der Umzug} (die \textit{Umzüge}) : le déménagement.
    \item \textbf{die Enttäuschung} (die \textit{-en}) : la déception.
    \item \textbf{die Ausbildung} (die \textit{-en}) : la formation, l'apprentissage.
    \item \textbf{der Gelegenheitsarbeiter} (die \textit{-}) : le travailleur occasionnel, journalier.
    \item \textbf{bestehen bleiben} (bleibt bestehen, blieb bestehen, ist bestehen geblieben) : persister, rester.
\end{itemize}
\end{definition}

\begin{exercice}[Questions de compréhension et de grammaire sur le texte]
\begin{enumerate}
    \item \textbf{Repérage :} surlignez dans le texte tous les connecteurs de cause (\all{weil, denn, wegen}) et de conséquence (\all{deshalb, so dass, also, darum}). Comptez-les.
    \item \textbf{Transformation (Weil $\rightarrow$ Deshalb) :} reprenez la phrase \all{„Weil es in den Dörfern keine Fabriken und keine großen Unternehmen gibt, finden sie dort keine Arbeit.“} et réécrivez-la en deux phrases principales reliées par \all{deshalb}.
    \item \textbf{Transformation (Wegen $\rightarrow$ Weil) :} reprenez le début du texte (\all{„Wegen der Armut ...“}) et transformez la cause nominale en subordonnée avec \all{weil} (verbe : \all{sein}).
    \item \textbf{Production courte :} écrivez 3 phrases en allemand pour expliquer pourquoi \textbf{vous} (ou un ami) pourriez quitter votre village. Utilisez obligatoirement : une phrase avec \all{weil}, une avec \all{deshalb}, une avec \all{wegen}.
\end{enumerate}
\end{exercice}

\begin{corrige}[Exercice « Die Landflucht in Kamerun »]
\begin{enumerate}
    \item \textbf{Connecteurs trouvés (8 au total) :}
    \begin{itemize}
        \item Cause : \all{Wegen} (1), \all{weil} (2 occurrences), \all{denn} (1). Total cause : 4.
        \item Conséquence : \all{so dass} (1), \all{deshalb} (1), \all{also} (1), \all{darum} (1). Total conséquence : 4.
    \end{itemize}
    \item \textbf{Weil $\rightarrow$ Deshalb :} \all{Es gibt in den Dörfern keine Fabriken und keine großen Unternehmen, deshalb finden die jungen Leute dort keine Arbeit.} (Inversion : \all{deshalb} position 1, \all{finden} position 2.)
    \item \textbf{Wegen $\rightarrow$ Weil :} \all{Weil die Armut groß ist, sehen sie keine Zukunft in der Landwirtschaft.} (Verbe \all{ist} à la fin de la subordonnée ; inversion dans la principale : \all{sehen sie}.)
    \item \textbf{Production (exemples de réponses) :}
    \begin{itemize}
        \item \all{Ich würde das Dorf verlassen, weil es hier keine weiterführende Schule gibt.}
        \item \all{Meine Eltern sind krank, deshalb muss ich in der Stadt Geld verdienen.}
        \item \all{Wegen der schlechten Straßen kann man die Ernte nicht verkaufen.}
    \end{itemize}
\end{enumerate}
\end{corrige}

\coursec{Méthode du commentaire et commentaire modèle}

Dans la section précédente, nous avons appris à exprimer la \cle{cause} et la \cle{conséquence} avec \all{weil}, \all{deshalb} ou \all{darum}. Ces outils grammaticaux ne servent pas seulement à construire des phrases isolées : ils sont la colonne vertébrale du \cle{commentaire de texte}, l'exercice roi de l'épreuve d'allemand au Bac. En effet, commenter le texte \all{„Warum verlassen die jungen Leute das Dorf ?“}, c'est précisément identifier les causes de l'exode rural, expliquer ses conséquences et donner son avis argumenté. Cette section vous donne la méthode complète, un plan modèle et toutes les phrases témoins pour réussir.

\courssub{Les quatre étapes fixes du commentaire}

\begin{methode}[Commenter un texte au Bac : les 4 étapes]
Un commentaire de texte en allemand suit toujours le même plan, en quatre étapes fixes. Ne l'improvisez jamais :
\begin{enumerate}
\item \textbf{Présenter le texte} (\all{Einleitung}) : indiquer le type de texte (article de journal, lettre, dialogue), le thème et, si elle est donnée, la source. Une ou deux phrases suffisent.
\item \textbf{Résumer les idées principales} (\all{Zusammenfassung}) : reformuler avec ses propres mots le contenu essentiel, sans recopier le texte et sans donner son avis.
\item \textbf{Analyser le texte} (\all{Analyse / Hauptteil}) : dégager les causes et les conséquences présentées par l'auteur, relever le lexique du thème et les procédés (questions, chiffres, exemples concrets).
\item \textbf{Donner son avis argumenté} (\all{eigene Meinung / Schluss}) : dire ce que l'on pense du problème et proposer éventuellement des solutions, en justifiant avec \all{weil} ou \all{denn}.
\end{enumerate}
\end{methode}

\begin{popfigure}
\begin{center}
\begin{tikzpicture}[
node distance=0.55cm,
etape/.style={rectangle, rounded corners=4pt, draw=popblue, fill=popblueL, thick, minimum width=7.2cm, minimum height=0.95cm, align=center, font=\small},
fleche/.style={-{Stealth[length=3mm]}, thick, poporange}
]
\node[etape] (e1) {\textbf{1. Einleitung} : Text sortieren, Thema und Quelle nennen};
\node[etape, below=of e1] (e2) {\textbf{2. Zusammenfassung} : die Hauptideen mit eigenen Worten};
\node[etape, below=of e2] (e3) {\textbf{3. Analyse} : Ursachen, Folgen, Wortschatz, Mittel};
\node[etape, below=of e3] (e4) {\textbf{4. Eigene Meinung} : Stellung nehmen und begründen};
\draw[fleche] (e1) -- (e2);
\draw[fleche] (e2) -- (e3);
\draw[fleche] (e3) -- (e4);
\end{tikzpicture}
\end{center}
\legende{Organigramme du commentaire de texte : quatre étapes, dans cet ordre.}
\end{popfigure}

\begin{attention}[Le résumé n'est pas une copie]
Deux erreurs coûtent cher au Bac : recopier des phrases entières du texte dans le résumé, et donner son avis trop tôt. Le résumé doit être \emph{objectif} et rédigé \emph{avec vos mots} ; l'opinion personnelle n'apparaît qu'à la quatrième étape, annoncée par \all{meiner Meinung nach} ou \all{ich denke, dass ...}.
\end{attention}

\courssub{Le plan modèle appliqué à notre texte}

Appliquons la méthode au texte de la leçon, \all{„Warum verlassen die jungen Leute das Dorf ?“}. Le plan se décline ainsi :

\begin{center}
\begin{tabularx}{\linewidth}{|l|X|X|}
\hline
\rowcolor{popblueL} \textbf{Partie} & \textbf{Contenu attendu} & \textbf{Phrase témoin} \\
\hline
\all{Einleitung} & Type de texte + thème & \all{Der Text handelt von der Landflucht in Kamerun.} \\
\hline
\all{Zusammenfassung} & Idées principales & \all{Viele junge Leute verlassen das Dorf und ziehen in die Stadt.} \\
\hline
\all{Hauptteil : Ursachen} & Causes de l'exode & \all{Der Autor nennt zwei Ursachen: die Armut und die fehlende Arbeit.} \\
\hline
\all{Hauptteil : Folgen} & Conséquences & \all{Deshalb werden die Dörfer immer älter und die Städte immer voller.} \\
\hline
\all{Hauptteil : Lösungen} & Solutions proposées & \all{Der Autor schlägt vor, Schulen und Krankenhäuser auf dem Land zu bauen.} \\
\hline
\all{Schluss : Meinung} & Avis argumenté & \all{Meiner Meinung nach ist die Landflucht ein Problem, weil das Land seine Jugend verliert.} \\
\hline
\end{tabularx}
\end{center}

\courssub{Le commentaire modèle}

Voici un commentaire complet, rédigé au niveau A2, que vous pouvez prendre comme modèle. Observez d'abord les connecteurs en gras : ils guident le correcteur à travers les quatre étapes.

\begin{textebox}[Kommentar : „Warum verlassen die jungen Leute das Dorf ?“ — Modell]
Der Text ist ein Zeitungsartikel. Er handelt von der Landflucht: Viele junge Leute verlassen das Dorf und ziehen in die Stadt, zum Beispiel nach Yaoundé oder Douala.

\textbf{Zuerst} fasst der Autor das Problem zusammen: Auf dem Land gibt es wenig Arbeit, und die Familien sind oft arm. \textbf{Dann} nennt er zwei Ursachen für die Landflucht. Erstens gibt es in den Dörfern keine Fabriken und keine Büros, deshalb finden die jungen Leute keine Arbeit. Zweitens fehlen gute Schulen und Krankenhäuser. \textbf{Außerdem} zeigt der Text die Folgen: Die Dörfer werden immer älter, weil nur die Großeltern bleiben, und die Städte werden immer größer und teurer.

Der Autor schlägt auch Lösungen vor: Die Regierung soll Straßen bauen und Arbeit auf dem Land schaffen, damit die Jugend im Dorf bleiben kann.

\textbf{Meiner Meinung nach} ist die Landflucht ein großes Problem, weil das Land seine Jugend verliert. Ich denke, dass die Regierung mehr Schulen und Krankenhäuser in den Dörfern bauen muss. Ich finde es auch wichtig, dass die jungen Leute stolz auf ihr Dorf sind. \textbf{Zum Schluss}: Die Stadt ist nicht immer das Paradies — das Dorf hat eine Zukunft, wenn man ihm hilft.
\end{textebox}

\textbf{Glossaire :} die Landflucht, — (l'exode rural) ; verlassen (verlässt, verließ, hat verlassen) = quitter ; die Fabrik, -en (l'usine) ; fehlen = manquer ; das Krankenhaus, die Krankenhäuser (l'hôpital) ; die Regierung, -en (le gouvernement) ; vorschlagen (séparable) = proposer ; stolz (fier) ; die Zukunft, — (l'avenir).

\begin{exemplebox}[Phrases témoins à réutiliser]
\begin{itemize}
\item \all{Der Text handelt von der Landflucht.} « Le texte traite de l'exode rural. »
\item \all{Der Autor nennt zwei Ursachen: die Armut und die fehlende Arbeit.} « L'auteur cite deux causes : la pauvreté et le manque de travail. »
\item \all{Meiner Meinung nach müssen die Regierungen Arbeit auf dem Land schaffen.} « À mon avis, les gouvernements doivent créer du travail à la campagne. »
\item \all{Ich finde es wichtig, dass die Jugend im Dorf bleiben kann.} « Je trouve important que les jeunes puissent rester au village. »
\end{itemize}
\end{exemplebox}

\courssub{Les connecteurs et l'expression de l'opinion}

Le commentaire doit être \emph{articulé}. Voici la boîte à outils des connecteurs de cette leçon :

\begin{center}
\begin{tabularx}{\linewidth}{|X|X|X|}
\hline
\rowcolor{popblueL} \textbf{Deutsch} & \textbf{Français} & \textbf{Fonction} \\
\hline
\all{zuerst} & d'abord & première étape \\
\hline
\all{dann} & ensuite & enchaînement \\
\hline
\all{außerdem} & de plus, en outre & ajout d'un argument \\
\hline
\all{deshalb / darum} & c'est pourquoi & conséquence \\
\hline
\all{meiner Meinung nach} & à mon avis & introduction de l'opinion \\
\hline
\all{zum Schluss} & pour conclure & conclusion \\
\hline
\end{tabularx}
\end{center}

\begin{propriete}[Exprimer son opinion : \all{ich denke, dass ...} / \all{ich finde es wichtig, dass ...}]
Pour donner son avis, on utilise deux constructions principales :
\begin{itemize}
\item \all{Ich denke, dass} + subordonnée (verbe conjugué \textbf{à la fin}) : \all{Ich denke, dass die Landflucht ein Problem ist.} « Je pense que l'exode rural est un problème. »
\item \all{Ich finde es wichtig, dass} + subordonnée : \all{Ich finde es wichtig, dass die Regierung hilft.} « Je trouve important que le gouvernement aide. »
\item Variante sans \all{dass} : \all{Meiner Meinung nach} + phrase principale (verbe en 2\up{e} position) : \all{Meiner Meinung nach ist das Dorf schön.}
\end{itemize}
\end{propriete}

\begin{attention}[Le verbe va aussi à la fin après \all{dass}]
Les francophones placent souvent le verbe juste après le sujet, comme en français. C'est faux : dans \all{Ich denke, dass die Landflucht ein Problem ist}, le verbe \all{ist} va \textbf{à la fin} de la subordonnée, exactement comme après \all{weil}. Comparez : « je pense que l'exode rural \emph{est} un problème » — \all{Ich denke, dass die Landflucht ein Problem \emph{ist}}. Ne dites jamais \all{dass die Landflucht ist ein Problem}.
\end{attention}

\begin{exoresolu}[Transformer un avis simple en phrase de commentaire]
Énoncé : voici une opinion simple : \all{Die Landflucht ist ein Problem.} Transformez-la en deux phrases de commentaire, l'une avec \all{ich denke, dass}, l'autre avec \all{meiner Meinung nach}.
\tcblower
Solution détaillée :
\begin{enumerate}
\item Avec \all{ich denke, dass}, on crée une subordonnée : le verbe \all{ist} passe à la fin : \all{Ich denke, dass die Landflucht ein Problem ist.}
\item Avec \all{meiner Meinung nach}, la phrase reste une principale : le verbe reste en deuxième position : \all{Meiner Meinung nach ist die Landflucht ein Problem.}
\end{enumerate}
On peut enrichir avec une cause : \all{Ich denke, dass die Landflucht ein Problem ist, weil die Dörfer ihre Jugend verlieren.} « Je pense que l'exode rural est un problème parce que les villages perdent leur jeunesse. »
\end{exoresolu}

\begin{exercice}[À vous de commenter]
Rédigez le début d'un commentaire (4 à 6 phrases) sur le texte de la leçon. Consignes : présentez le texte, résumez le problème en une phrase, citez une cause avec \all{weil} et donnez votre avis avec \all{ich denke, dass}. Utilisez au moins deux connecteurs du tableau.
\end{exercice}

\begin{corrige}[Exercice : À vous de commenter]
Proposition de corrigé : \all{Der Text handelt von der Landflucht in Kamerun. Viele junge Leute verlassen das Dorf, weil es dort keine Arbeit gibt. Deshalb ziehen sie in die Stadt. Ich denke, dass die Landflucht ein großes Problem ist, weil die Dörfer ohne Jugend sterben. Meiner Meinung nach muss die Regierung Arbeit auf dem Land schaffen.} « Le texte traite de l'exode rural au Cameroun. Beaucoup de jeunes quittent le village parce qu'il n'y a pas de travail là-bas. C'est pourquoi ils s'installent en ville. Je pense que l'exode rural est un grand problème parce que les villages meurent sans jeunesse. À mon avis, le gouvernement doit créer du travail à la campagne. »
\end{corrige}

\begin{aretenir}[L'essentiel du commentaire]
Quatre étapes fixes : présenter, résumer, analyser, donner son avis. Connecteurs obligatoires : \all{zuerst}, \all{dann}, \all{außerdem}, \all{deshalb}, \all{meiner Meinung nach}, \all{zum Schluss}. Après \all{dass} et \all{weil}, le verbe conjugué va à la fin. L'avis personnel vient en dernier, jamais dans le résumé.
\end{aretenir}

\begin{savaistu}[Comment le commentaire est-il noté au Bac ?]
Au baccalauréat, le commentaire de texte est évalué selon trois critères : la \emph{compréhension} (avez-vous saisi les idées du texte ?), la \emph{correction grammaticale} (place du verbe, cas, conjugaison) et la \emph{richesse du vocabulaire}. Astuce de correcteur : réutiliser le lexique de la leçon — \all{die Landflucht}, \all{die Armut}, \all{umziehen}, \all{die Arbeit schaffen} — rapporte directement des points, car cela prouve que vous maîtrisez le thème. Un commentaire simple mais correct et bien articulé vaut toujours mieux qu'une phrase compliquée pleine de fautes.
\end{savaistu}

\coursec{Production guidée}

Après avoir appris à lire, comprendre et commenter le texte \all{Warum verlassen die jungen Leute das Dorf ?}, tu es maintenant capable de \emph{produire} toi-même : d'abord un court texte écrit de 8 à 10 lignes, puis un mini-exposé oral d'une minute. C'est l'aboutissement de la leçon : réutiliser le vocabulaire du thème (\all{das Dorf}, \all{die Stadt}, \all{umziehen}, \all{die Arbeit}, \all{die Armut}) et les structures grammaticales étudiées (les subordonnées en \all{weil}, les connecteurs de conséquence comme \all{deshalb}) dans un énoncé personnel, organisé et correct.

\courssub{Production écrite guidée : « L'exode rural dans ton pays »}

\begin{methode}[Comment réussir ta production écrite en 5 étapes]
\begin{enumerate}
\item \textbf{Respecte le plan imposé.} Ton texte doit contenir quatre parties, dans l'ordre : 1) la situation (2 phrases), 2) deux causes (2 phrases avec \all{weil}), 3) deux conséquences (dont 1 phrase avec \all{deshalb}), 4) une solution et ton avis personnel (2-3 phrases).
\item \textbf{Utilise au moins 5 mots du lexique obligatoire} : \all{das Dorf}, \all{die Stadt}, \all{umziehen}, \all{die Arbeit}, \all{die Armut}. Souligne-les dans ton brouillon pour vérifier.
\item \textbf{Place correctement le verbe.} Dans la phrase principale, le verbe conjugué est en position 2 ; après \all{weil}, il va à la fin ; après \all{deshalb}, le verbe reste en position 2 et le sujet passe après le verbe (inversion).
\item \textbf{Rédige d'abord au brouillon}, phrase par phrase, en suivant la grille d'aide ci-dessous. Phrases courtes et simples : mieux vaut 10 phrases correctes que 3 phrases compliquées et fausses.
\item \textbf{Relis et corrige} avec la liste de contrôle (voir la boîte \emph{Attention} plus bas) avant de recopier au propre.
\end{enumerate}
\end{methode}

\begin{popfigure}
\begin{tikzpicture}[
  box/.style={draw=popblue, fill=popblueL, rounded corners=3pt, text width=3.1cm, align=center, font=\small, minimum height=2.6cm},
  head/.style={font=\small\bfseries, text=popdark},
  arr/.style={-{Stealth[length=2.5mm]}, thick, poporange}
]
\node[head] at (0,2.0) {1. Situation};
\node[box] at (0,0) {In meinem Land verlassen viele junge Leute das Dorf.};
\node[head] at (3.9,2.0) {2. Ursachen};
\node[box] at (3.9,0) {\all{weil} + verbe final : \emph{weil es auf dem Land keine Arbeit gibt.}};
\node[head] at (7.8,2.0) {3. Folgen};
\node[box] at (7.8,0) {\all{deshalb} + inversion : \emph{Deshalb bleiben nur die Alten im Dorf.}};
\node[head] at (11.7,2.0) {4. Lösung \& Meinung};
\node[box] at (11.7,0) {\all{Ich denke, dass... / Meiner Meinung nach...}};
\draw[arr] (1.7,0) -- (2.2,0);
\draw[arr] (5.6,0) -- (6.1,0);
\draw[arr] (9.5,0) -- (10.0,0);
\end{tikzpicture}
\legende{Grille d'aide : les quatre colonnes du plan imposé avec les connecteurs obligatoires.}
\end{popfigure}

\begin{exemplebox}[Phrases modèles pour chaque partie du plan]
\begin{itemize}
\item \textbf{Situation :} \all{Viele junge Leute verlassen das Dorf und ziehen in die Stadt.} « Beaucoup de jeunes quittent le village et s'installent en ville. »
\item \textbf{Cause 1 :} \all{Viele junge Leute verlassen das Dorf, weil sie Arbeit suchen.} « Beaucoup de jeunes quittent le village parce qu'ils cherchent du travail. »
\item \textbf{Cause 2 :} \all{Sie ziehen in die Stadt, weil es dort Schulen und Krankenhäuser gibt.} « Ils s'installent en ville parce qu'il y a là des écoles et des hôpitaux. »
\item \textbf{Conséquence :} \all{Deshalb bleiben nur die Alten im Dorf.} « C'est pourquoi seuls les vieux restent au village. »
\item \textbf{Solution et avis :} \all{Ich denke, dass der Staat Fabriken auf dem Land bauen muss. Meiner Meinung nach ist das Leben im Dorf auch schön.} « Je pense que l'État doit construire des usines à la campagne. À mon avis, la vie au village est belle aussi. »
\end{itemize}
\end{exemplebox}

\begin{textebox}[Production modèle : „Die Landflucht in meinem Land“]
In Kamerun verlassen viele junge Leute das Dorf und ziehen in die Stadt, zum Beispiel nach Yaoundé oder Douala. Sie verlassen das Dorf, weil es auf dem Land wenig Arbeit gibt. Sie ziehen auch in die Stadt, weil sie dort mehr Geld verdienen und besser leben wollen. Die Armut auf dem Land ist ein großes Problem.

Deshalb bleiben oft nur die alten Leute und die Kinder im Dorf. Die Felder liegen brach, und die Dörfer werden immer kleiner.

Ich denke, dass der Staat Schulen, Krankenhäuser und Fabriken auf dem Land bauen muss. Meiner Meinung nach ist das Leben im Dorf ruhig und gesund, aber die jungen Leute brauchen Arbeit und Hoffnung.
\end{textebox}

\textbf{Glossaire du texte modèle :} \all{die Landflucht} « l'exode rural » ; \all{verdienen} « gagner (de l'argent) » ; \all{die Armut} « la pauvreté » ; \all{brach liegen} « être en friche » ; \all{ruhig} « calme » ; \all{gesund} « sain » ; \all{die Hoffnung, -en} « l'espoir ».

\textbf{Questions sur le modèle :} 1) Woher kommen die jungen Leute, die nach Yaoundé ziehen ? 2) Nenne zwei Gründe für die Landflucht. 3) Was ist die Lösung des Autors ?

\begin{exoresolu}[Rédiger les deux premières parties du plan]
Énoncé : rédige la situation et les deux causes de l'exode rural au Cameroun, avec deux subordonnées en \all{weil} et les mots \all{das Dorf}, \all{die Stadt}, \all{die Arbeit}.
\tcblower
\textbf{Solution détaillée.} Situation : \all{In meinem Land verlassen viele junge Leute das Dorf und ziehen in die Stadt.} Cause 1 : \all{Sie verlassen das Dorf, weil sie auf dem Land keine Arbeit finden.} (verbe \all{finden} en position finale après \all{weil}). Cause 2 : \all{Sie ziehen in die Stadt, weil das Leben dort moderner ist.} (verbe \all{ist} en finale). Vérification : les trois mots obligatoires sont employés, les deux verbes des subordonnées sont bien à la fin.
\end{exoresolu}

\begin{attention}[Relis ton texte : la liste de contrôle]
Avant de rendre ta copie, vérifie point par point :
\begin{itemize}
\item \textbf{Verbe en position 2} dans chaque phrase principale : \all{Viele Leute \emph{ziehen} in die Stadt.} (et non : \all{Viele Leute in die Stadt ziehen.})
\item \textbf{Verbe à la fin} après \all{weil} et \all{dass} : \all{..., weil sie Arbeit \emph{suchen}.} (et non : \all{weil sie suchen Arbeit}).
\item \textbf{Inversion après deshalb} : \all{Deshalb \emph{bleiben} nur die Alten.} (le sujet \all{die Alten} passe après le verbe ; jamais \all{Deshalb die Alten bleiben...}).
\item \textbf{Article et pluriel corrects} : \all{das Dorf, die Dörfer} ; \all{die Stadt, die Städte} ; \all{die Arbeit, -en}.
\item \textbf{Majuscule à tous les noms} : \all{das Dorf, die Armut, die Arbeit}.
\item \textbf{Verbe à particule} : \all{umziehen} se sépare en position 2 : \all{Sie \emph{ziehen} in die Stadt \emph{um}.}
\end{itemize}
\end{attention}

\courssub{Production orale : mini-exposé d'une minute}

\begin{experience}[Mini-exposé : „Warum ziehen die jungen Leute in die Stadt ?“]
Par groupes de deux, préparez un exposé oral d'environ une minute (6 à 8 phrases) en suivant ce schéma :
\begin{enumerate}
\item \textbf{Introduction :} \all{Ich spreche heute über die Landflucht.} « Je parle aujourd'hui de l'exode rural. »
\item \textbf{Deux causes} avec \all{weil} : \all{Die jungen Leute ziehen in die Stadt, weil...}
\item \textbf{Une conséquence} avec \all{deshalb} : \all{Deshalb...}
\item \textbf{Conclusion personnelle :} \all{Meiner Meinung nach...}
\end{enumerate}
L'un expose, l'autre écoute et coche sur une fiche : les 5 mots du lexique sont-ils utilisés ? Les verbes sont-ils bien placés ? Puis on échange les rôles. Parle lentement, articule, et regarde ton public : un exposé n'est pas une lecture !
\end{experience}

\courssub{Correction collective et auto-évaluation}

\begin{methode}[La correction collective au tableau]
\begin{enumerate}
\item Deux ou trois élèves écrivent une phrase de leur production au tableau.
\item La classe cherche les erreurs ensemble : place du verbe après \all{weil} (finale) et après \all{deshalb} (position 2 avec inversion), articles, pluriels, majuscules.
\item Le professeur corrige en expliquant \emph{pourquoi} : chaque erreur est rattachée à la règle de grammaire correspondante.
\item Chaque élève corrige ensuite sa propre copie en rouge et note la règle à revoir.
\end{enumerate}
\end{methode}

\begin{center}
\begin{tabularx}{\linewidth}{|X|l|}
\hline
\rowcolor{popblueL} \textbf{Objectif de la leçon} & \textbf{Oui / Non} \\
\hline
Je comprends un texte sur l'exode rural. & \\
\hline
J'emploie le vocabulaire : \all{das Dorf, die Stadt, umziehen, die Arbeit, die Armut}. & \\
\hline
Je construis une subordonnée avec \all{weil} (verbe à la fin). & \\
\hline
J'exprime une conséquence avec \all{deshalb} (inversion). & \\
\hline
Je rédige un texte de 8 à 10 lignes avec un plan en 4 parties. & \\
\hline
Je présente un mini-exposé oral d'une minute. & \\
\hline
\end{tabularx}
\end{center}

Coche chaque case honnêtement. Si une case reste décochée, revois la section correspondante de la leçon et refais l'exercice associé : c'est ainsi qu'on progresse durablement en allemand.

\begin{exercice}[À faire à la maison]
Rédige la version définitive de ta production « L'exode rural dans ton pays » (8 à 10 lignes) en appliquant les corrections faites en classe. Souligne les 5 mots du lexique, encadre les verbes des subordonnées en \all{weil} et entoure le verbe en position 2 après \all{deshalb}.
\end{exercice}

\newpage

\coursec{Activité d'intégration}

\courssub{Objectifs de l'activité}

À la fin de cette activité d'intégration, tu seras capable de :
\begin{itemize}
\item mobiliser le vocabulaire thématique de l'exode rural (\all{das Dorf}, \all{die Stadt}, \all{die Landflucht}, \all{die Arbeit}, \all{die Armut}, \all{umziehen}) dans une situation de communication réelle ;
\item exprimer la cause avec une subordonnée en \all{weil} et la conséquence avec \all{deshalb} / \all{darum} ;
\item défendre oralement un point de vue dans un mini-débat, avec des arguments clairs et des phrases correctes ;
\item rédiger un court article de journal (environ 10 lignes) avec une introduction, des arguments et une conclusion personnelle ;
\item travailler en groupe, écouter les autres et réagir à leurs arguments.
\end{itemize}

\courssub{Situation}

Tu es élève en classe de Première A au lycée bilingue de Bafoussam. Le journal de ton école, \all{„Die Brücke“}, prépare un numéro spécial sur un sujet qui touche toutes les familles du Cameroun : le départ des jeunes des villages vers les grandes villes comme Douala et Yaoundé. Le rédacteur en chef du journal lance un appel aux élèves germanistes : qui écrira le meilleur article sur ce thème ?

Pour préparer les articles, la classe organise d'abord un débat. La moitié de la classe défend le village, l'autre moitié défend la ville. Voici le document qui lance la réflexion, affiché au tableau par ton professeur.

\begin{textebox}[Affiche du journal de l'école]
\textbf{Débat du mois : Soll man im Dorf bleiben oder in die Stadt ziehen?}\par
Immer mehr junge Leute verlassen die Dörfer in Kamerun. Sie suchen in der Stadt Arbeit, eine bessere Schule und ein modernes Leben. Aber das Leben in der Stadt ist teuer und oft schwer.\par
Was meinst du? Schreib einen Artikel für unsere Schülerzeitung „Die Brücke“!\par
Dein Artikel muss haben: mindestens 5 Wörter aus unserem Wortschatz, 2 Sätze mit \emph{weil}, 1 Satz mit \emph{deshalb} und eine persönliche Meinung am Ende.
\end{textebox}

\begin{definition}[Glossaire de l'affiche]
\begin{itemize}
\item \all{immer mehr} : de plus en plus
\item \all{verlassen (verlässt, verließ, hat verlassen)} : quitter
\item \all{suchen} : chercher
\item \all{teuer} : cher (coûteux)
\item \all{schwer} : difficile, lourd
\item \all{die Schülerzeitung, -en} : le journal des élèves
\item \all{die Meinung, -en} : l'opinion, l'avis
\item \all{der Wortschatz, „-e} : le vocabulaire
\end{itemize}
\end{definition}

\courssub{Consignes}

\textbf{Étape 1 — Compréhension du document (5 min, individuel).} Lis l'affiche du journal et réponds par écrit en allemand, avec des phrases complètes :
\begin{enumerate}
\item \all{Was machen viele junge Leute in Kamerun?}
\item \all{Was suchen sie in der Stadt?} (cite deux choses)
\item \all{Welches Problem gibt es in der Stadt?}
\end{enumerate}

\textbf{Étape 2 — Préparation des arguments en groupe (15 min).} La classe est divisée en deux équipes :
\begin{itemize}
\item \textbf{Groupe A} : \all{Wir bleiben im Dorf!} (les partisans du village)
\item \textbf{Groupe B} : \all{Wir ziehen in die Stadt!} (les partisans de la ville)
\end{itemize}
Chaque groupe prépare \textbf{4 arguments}. Chaque argument doit contenir une subordonnée en \all{weil} ; au moins deux arguments doivent être complétés par une phrase de conséquence avec \all{deshalb}. Exemple de modèle : \all{Man soll im Dorf bleiben, weil das Leben dort ruhig und billig ist.} « On doit rester au village parce que la vie y est calme et bon marché. »

\textbf{Étape 3 — Le mini-débat oral (10 min).} Un porte-parole de chaque groupe présente les arguments à tour de rôle. Règles du débat : on parle allemand, on écoute sans couper la parole, et chaque groupe doit réagir à au moins un argument de l'autre groupe avec \all{Ich bin nicht einverstanden, weil …} (« Je ne suis pas d'accord parce que… ») ou \all{Das stimmt, aber …} (« C'est vrai, mais… »).

\textbf{Étape 4 — Rédaction de l'article (20 min, individuel).} Chaque élève rédige un article d'environ 10 lignes intitulé \all{„Die Landflucht in Kamerun“}. L'article doit contenir :
\begin{enumerate}
\item une introduction qui présente le problème (1 ou 2 phrases) ;
\item au moins \textbf{5 mots du lexique} de la leçon (souligne-les) ;
\item \textbf{2 subordonnées en \all{weil}} (attention au verbe à la fin !) ;
\item \textbf{1 phrase avec \all{deshalb}} (attention à l'inversion !) ;
\item une conclusion avec ton avis personnel : \all{Ich meine, dass …} / \all{Meiner Meinung nach …}
\end{enumerate}

\textbf{Étape 5 — Correction et affichage (10 min).} Échange ton article avec ton voisin : il vérifie la place du verbe dans les subordonnées et la présence des 5 mots du lexique. Les meilleurs articles sont lus à voix haute puis affichés au tableau d'allemand.

\courssub{Ressources à mobiliser}

\begin{aretenir}[Boîte à outils de l'article]
\textbf{Lexique :} \all{das Dorf, „-er} (le village) ; \all{die Stadt, „-e} (la ville) ; \all{die Landflucht} (l'exode rural) ; \all{die Arbeit, -en} (le travail) ; \all{die Armut} (la pauvreté) ; \all{umziehen (zieht um, zog um, ist umgezogen)} (déménager) ; \all{verlassen} (quitter) ; \all{die Familie, -n} (la famille) ; \all{der Bauer, -n} (le paysan) ; \all{das Feld, -er} (le champ).\par
\textbf{Cause :} phrase principale + \all{weil} + sujet + … + verbe conjugué \textbf{à la fin}.\par
\all{Viele Jugendliche gehen in die Stadt, weil es im Dorf keine Arbeit gibt.}\par
\textbf{Conséquence :} \all{deshalb} / \all{darum} + verbe conjugué en \textbf{position 2} + sujet.\par
\all{Es gibt im Dorf keine Arbeit, deshalb ziehen viele Jugendliche in die Stadt.}
\end{aretenir}

\begin{propriete}[Rappel de grammaire : cause et conséquence]
\begin{center}
\begin{tabularx}{\linewidth}{|X|X|X|}\hline
\rowcolor{popblueL} Outil & Place du verbe & Exemple \\ \hline
\all{weil} (parce que) & verbe conjugué à la fin de la subordonnée & \all{Er bleibt, weil er das Dorf liebt.} \\ \hline
\all{denn} (car) & verbe en position 2 (phrase coordonnée) & \all{Er bleibt, denn er liebt das Dorf.} \\ \hline
\all{deshalb} / \all{darum} (c'est pourquoi) & inversion : verbe en position 2, puis le sujet & \all{Er liebt das Dorf, deshalb bleibt er.} \\ \hline
\end{tabularx}
\end{center}
Attention : \all{weil} introduit la \textbf{cause} (la raison), \all{deshalb} introduit la \textbf{conséquence} (le résultat). Ne les confonds pas !
\end{propriete}

\begin{popfigure}
\begin{tikzpicture}[node distance=0.4cm]
\node[draw=popblue, fill=popblueL, rounded corners, align=center, font=\small] (haupt) {Phrase principale\\ \all{Viele Jugendliche gehen in die Stadt,}};
\node[draw=poporange, fill=poporangeL, rounded corners, align=center, font=\small, below=of haupt] (neb) {Subordonnée en \all{weil}\\ \all{weil es im Dorf keine Arbeit} \textbf{\all{gibt.}}};
\node[draw=popgreen, fill=popgreenL, rounded corners, align=center, font=\small, below=of neb] (folge) {Conséquence\\ \all{Deshalb} \textbf{\all{ziehen}} \all{sie in die Stadt.}};
\draw[-{Stealth}, thick, popdark] (haupt) -- (neb);
\draw[-{Stealth}, thick, popdark] (neb) -- (folge);
\end{tikzpicture}
\legende{De la cause à la conséquence : après \all{weil}, le verbe va à la fin ; après \all{deshalb}, le verbe reste en position 2.}
\end{popfigure}

\begin{attention}[Pièges à éviter]
\begin{itemize}
\item Après \all{weil}, le verbe conjugué va \textbf{à la fin} : on écrit \all{…, weil das Leben dort billig \textbf{ist}.} et non \all{weil das Leben ist billig}.
\item Après \all{deshalb}, il y a \textbf{inversion} : le verbe reste en position 2, le sujet le suit : \all{Deshalb \textbf{ziehen} viele Leute in die Stadt.}
\item \all{die Stadt} prend l'accusatif après \all{in} quand il y a mouvement, et le datif quand il y a lieu : \all{Sie zieht in die Stadt. Sie wohnt jetzt in der Stadt.}
\item N'oublie pas la majuscule à tous les noms : \all{das Dorf}, \all{die Arbeit}, \all{die Armut}.
\item Faux ami : \all{bekommen} ne veut pas dire « devenir » mais « recevoir, obtenir » : \all{Sie bekommt eine Arbeit} = « elle obtient un travail ». Pour « devenir », utilise \all{werden} : \all{Er wird Bauer.} « Il devient paysan. »
\item \all{umziehen} est un verbe à particule séparable : \all{Wir \textbf{ziehen} nächste Woche \textbf{um}.} « Nous déménageons la semaine prochaine. »
\end{itemize}
\end{attention}

\courssub{Proposition de production et corrigé commenté}

\begin{exoresolu}[Article modèle : „Die Landflucht in Kamerun“]
Rédige un article d'environ 10 lignes pour le journal de l'école, avec 5 mots du lexique, 2 subordonnées en \all{weil}, 1 phrase avec \all{deshalb} et une conclusion personnelle.
\tcblower
\textbf{Article modèle :}\par
\all{\textbf{Die Landflucht in Kamerun}}\par
\all{In Kamerun verlassen immer mehr junge Leute das Dorf und ziehen in die Stadt. Warum? Viele Jugendliche gehen nach Douala oder Yaoundé, weil es im Dorf wenig Arbeit gibt. Auch die Schulen und Krankenhäuser sind in der Stadt besser. Die Armut auf dem Land ist groß, deshalb suchen die jungen Leute ihr Glück in der Stadt. Aber das Leben dort ist nicht einfach, weil alles sehr teuer ist: die Miete, das Essen und der Transport. Manche finden keine Arbeit und leben in schlechten Vierteln. Ich meine, dass das Dorf auch Vorteile hat: Das Leben ist dort ruhig, billig und gesund. Meiner Meinung nach muss der Staat den Dörfern helfen, damit die Jugendlichen nicht alle weggehen.}\par
\textbf{Traduction :} « Au Cameroun, de plus en plus de jeunes quittent le village et déménagent en ville. Pourquoi ? Beaucoup de jeunes vont à Douala ou Yaoundé parce qu'il y a peu de travail au village. Les écoles et les hôpitaux sont aussi meilleurs en ville. La pauvreté à la campagne est grande, c'est pourquoi les jeunes cherchent leur chance en ville. Mais la vie là-bas n'est pas facile, parce que tout est très cher : le loyer, la nourriture et le transport. Certains ne trouvent pas de travail et vivent dans des quartiers pauvres. Je pense que le village a aussi des avantages : la vie y est calme, bon marché et saine. À mon avis, l'État doit aider les villages, pour que les jeunes ne partent pas tous. »\par
\textbf{Glossaire complémentaire :} \all{das Krankenhaus, „-er} (l'hôpital) ; \all{die Miete, -n} (le loyer) ; \all{das Viertel, -} (le quartier) ; \all{der Vorteil, -e} (l'avantage) ; \all{gesund} (sain, en bonne santé) ; \all{weggehen (geht weg, ging weg, ist weggegangen)} (partir, s'en aller).\par
\textbf{Commentaire des choix de langue :}
\begin{enumerate}
\item \textbf{Lexique (plus de 5 mots) :} \all{das Dorf}, \all{die Stadt}, \all{die Landflucht} (dans le titre), \all{die Arbeit}, \all{die Armut}, \all{verlassen}, \all{in die Stadt ziehen}.
\item \textbf{Deux subordonnées en \all{weil}} avec le verbe conjugué à la fin : \all{…, weil es im Dorf wenig Arbeit \textbf{gibt}.} et \all{…, weil alles sehr teuer \textbf{ist}.}
\item \textbf{Une conséquence avec \all{deshalb}} et inversion correcte : \all{…, deshalb \textbf{suchen} die jungen Leute ihr Glück in der Stadt.}
\item \textbf{Conclusion personnelle} avec \all{Ich meine, dass …} (verbe à la fin après \all{dass}) et \all{Meiner Meinung nach} + verbe en position 2 : \all{Meiner Meinung nach \textbf{muss} der Staat …}
\item \textbf{Structure} : introduction (le problème) → causes du départ → conséquences et revers de la ville → avis personnel. C'est le plan attendu d'un court article d'opinion.
\end{enumerate}
\end{exoresolu}

\begin{exercice}[À toi de jouer : cause ou conséquence ?]
Complète les phrases avec \all{weil} ou \all{deshalb}, puis traduis-les en français :
\begin{enumerate}
\item \all{Viele Bauern verkaufen ihre Felder, … sie sehr arm sind.}
\item \all{Im Dorf gibt es kein Krankenhaus, … gehen kranke Leute in die Stadt.}
\item \all{Mein Bruder wohnt jetzt in Douala, … er dort eine Arbeit gefunden hat.}
\item \all{Das Leben in der Stadt ist teuer, … spart er jeden Monat Geld.}
\end{enumerate}
\end{exercice}

\begin{corrige}[Exercice : cause ou conséquence ?]
\begin{enumerate}
\item \all{Viele Bauern verkaufen ihre Felder, \textbf{weil} sie sehr arm sind.} « Beaucoup de paysans vendent leurs champs parce qu'ils sont très pauvres. » (cause : verbe \all{sind} à la fin)
\item \all{Im Dorf gibt es kein Krankenhaus, \textbf{deshalb} gehen kranke Leute in die Stadt.} « Il n'y a pas d'hôpital au village, c'est pourquoi les malades vont en ville. » (conséquence : inversion \all{gehen kranke Leute})
\item \all{Mein Bruder wohnt jetzt in Douala, \textbf{weil} er dort eine Arbeit gefunden hat.} « Mon frère habite maintenant à Douala parce qu'il y a trouvé un travail. » (cause : auxiliaire \all{hat} à la fin, après le participe \all{gefunden})
\item \all{Das Leben in der Stadt ist teuer, \textbf{deshalb} spart er jeden Monat Geld.} « La vie en ville est chère, c'est pourquoi il économise de l'argent chaque mois. » (conséquence : inversion \all{spart er})
\end{enumerate}
\end{corrige}

\begin{experience}[Variante : la table ronde]
Si le temps le permet, organise une \all{Runde Diskussion} (table ronde) : quatre élèves jouent les rôles d'un jeune du village, de sa mère, d'un ami déjà installé à Douala et du chef du village. Chacun défend sa position en 2 ou 3 phrases avec \all{weil} et \all{deshalb}. La classe vote ensuite : \all{Wer hat die besten Argumente?} (« Qui a les meilleurs arguments ? »)
\end{experience}

\begin{savaistu}[Le Cameroun et l'Allemagne : deux histoires de Landflucht]
L'exode rural n'est pas seulement un phénomène camerounais. En Allemagne aussi, beaucoup de jeunes quittent les petits villages, surtout dans les régions de l'Est, pour étudier ou travailler dans les grandes villes comme Berlin, Hamburg ou München. Le mot \all{die Landflucht} est composé de \all{das Land} (la campagne) et \all{die Flucht} (la fuite) : littéralement « la fuite de la campagne ». Pour retenir les jeunes, les communes allemandes misent sur le \all{Breitband} (le haut débit internet), les transports et les écoles : des pistes de réflexion pour ton article !
\end{savaistu}

\courssub{Bilan de l'activité}

\begin{aretenir}[Ce que j'ai appris à faire]
À l'issue de cette activité, je sais :
\begin{itemize}
\item comprendre et exploiter un document authentique (une affiche de journal scolaire) sur l'exode rural ;
\item employer à l'oral et à l'écrit le vocabulaire du thème : \all{das Dorf}, \all{die Stadt}, \all{die Landflucht}, \all{die Arbeit}, \all{die Armut}, \all{umziehen}, \all{verlassen} ;
\item construire une subordonnée de cause en \all{weil} avec le verbe conjugué à la fin ;
\item exprimer la conséquence avec \all{deshalb} / \all{darum} suivi de l'inversion verbe-sujet ;
\item structurer un court article d'opinion : introduction, arguments, avis personnel ;
\item défendre un point de vue et réagir aux arguments des autres dans un débat respectueux.
\end{itemize}
\textbf{Auto-évaluation :} coche ce que tu as réussi dans ton article : 5 mots du lexique soulignés ? 2 \all{weil} avec le verbe à la fin ? 1 \all{deshalb} avec inversion ? une conclusion avec ton avis ? Si oui, ton article est prêt pour \all{„Die Brücke“} !
\end{aretenir}

\newpage

\coursec{Méthodes et exercices résolus}

\courssub{Méthode 1 : Comprendre un texte sur l'exode rural (global puis détaillé)}

\begin{methode}[Lire et comprendre un texte en deux étapes]
\begin{enumerate}
\item \textbf{Lecture globale} : lis le texte une première fois sans dictionnaire. Repère le thème (ici : \all{die Landflucht}, l'exode rural), les personnages (\all{die jungen Leute}), les lieux (\all{das Dorf} / \all{die Stadt}) et la question posée dans le titre.
\item \textbf{Repérage des mots-clés} : souligne les mots du champ lexical de la campagne (\all{das Dorf, der Bauernhof, die Arbeit auf dem Land}) et de la ville (\all{die Stadt, die Arbeit, die Schule, das Krankenhaus}).
\item \textbf{Lecture détaillée} : pour chaque question, localise la ligne du texte qui contient la réponse. Cherche les connecteurs de cause (\all{weil, denn, da}) et de conséquence (\all{deshalb, darum, deswegen}) : ils signalent les raisons du départ et ses effets.
\item \textbf{Reformulation} : réponds par une phrase complète en allemand, en reprenant les mots de la question. Ne recopie jamais tout le texte.
\item \textbf{Vérification} : contrôle la place du verbe (2\up{e} position dans la phrase principale, dernière position après \all{weil}).
\end{enumerate}
\end{methode}

\begin{exoresolu}[Questions de compréhension sur le texte]
\textbf{Énoncé.} Lis le texte « Warum verlassen die jungen Leute das Dorf ? » et réponds en allemand par des phrases complètes :
\begin{enumerate}
\item Warum verlassen viele junge Leute das Dorf?
\item Was gibt es in der Stadt, aber nicht auf dem Land?
\item Was machen die jungen Leute in der Stadt?
\end{enumerate}
\tcblower
\textbf{Solution détaillée.}
\begin{enumerate}
\item La question porte sur la \emph{cause}. On cherche dans le texte une subordonnée en \all{weil} : « Weil es auf dem Land wenig Arbeit gibt, ziehen viele junge Leute in die Stadt. » Réponse rédigée : \all{Viele junge Leute verlassen das Dorf, weil es auf dem Land wenig Arbeit gibt.} « Beaucoup de jeunes quittent le village parce qu'il y a peu de travail à la campagne. » Justification : après \all{weil}, le verbe conjugué \all{gibt} va à la fin de la subordonnée.
\item Question sur le \emph{contraste} ville/campagne. Le texte oppose \all{die Stadt} (écoles, hôpitaux, cinémas) et \all{das Land}. Réponse : \all{In der Stadt gibt es Schulen, Krankenhäuser und mehr Arbeit, aber auf dem Land gibt es das nicht.} « En ville, il y a des écoles, des hôpitaux et plus de travail, mais pas à la campagne. »
\item Question sur les \emph{activités}. Réponse : \all{In der Stadt arbeiten die jungen Leute in Fabriken oder Büros, und sie besuchen eine Schule oder eine Universität.} « En ville, les jeunes travaillent dans des usines ou des bureaux et fréquentent une école ou une université. »
\end{enumerate}
\textbf{Conclusion.} Une bonne réponse de compréhension reprend les mots de la question, cite l'idée du texte et respecte l'ordre des mots allemand.
\end{exoresolu}

\courssub{Méthode 2 : Employer le vocabulaire du thème (Dorf, Stadt, Landflucht)}

\begin{methode}[Construire et utiliser le champ lexical de l'exode rural]
\begin{enumerate}
\item \textbf{Apprends chaque nom avec son article et son pluriel} : \all{das Dorf, die Dörfer} ; \all{die Stadt, die Städte} ; \all{die Landflucht} (sans pluriel) ; \all{die Arbeit} (au sens de « travail », le pluriel \all{die Arbeiten} est rare et signifie plutôt « les travaux, les devoirs ») ; \all{die Armut} (sans pluriel).
\item \textbf{Retiens les couples d'opposition} : \all{das Dorf} $\leftrightarrow$ \all{die Stadt} ; \all{das Land} (la campagne) $\leftrightarrow$ \all{die Stadt} ; \all{bleiben} $\leftrightarrow$ \all{weggehen / verlassen / umziehen}.
\item \textbf{Mémorise les verbes à particule} : \all{umziehen} (déménager) se sépare : \all{Er zieht in die Stadt um.} Au Perfekt : \all{Er ist in die Stadt umgezogen.} (auxiliaire \all{sein}, car il y a déplacement).
\item \textbf{Utilise les familles de mots} : \all{fliehen} $\rightarrow$ \all{die Flucht} $\rightarrow$ \all{die Landflucht} ; \all{arm} $\rightarrow$ \all{die Armut} ; \all{arbeiten} $\rightarrow$ \all{die Arbeit, der Arbeiter, die Arbeitslosigkeit}.
\item \textbf{Choisis la bonne préposition} : \all{auf dem Land} (à la campagne), \all{in der Stadt} (en ville), \all{in die Stadt ziehen} (accusatif, mouvement), \all{vom Dorf weggehen}.
\end{enumerate}
\end{methode}

\begin{center}
\begin{tabularx}{\linewidth}{|X|X|X|}\hline
\rowcolor{popblueL} Français & Deutsch & Remarque \\ \hline
le village & das Dorf, die Dörfer & neutre, Umlaut au pluriel \\ \hline
la ville & die Stadt, die Städte & féminin, Umlaut au pluriel \\ \hline
l'exode rural & die Landflucht & sans pluriel ; genre du dernier élément (\all{die Flucht}) \\ \hline
le travail & die Arbeit & féminin ; pluriel rare \\ \hline
la pauvreté & die Armut & féminin, sans pluriel \\ \hline
déménager & umziehen (zieht um, ist umgezogen) & particule séparable, auxiliaire \all{sein} \\ \hline
quitter & verlassen (verlässt, hat verlassen) & particule inséparable \all{ver-} \\ \hline
partir & weggehen (geht weg, ist weggegangen) & particule séparable \\ \hline
rester & bleiben (bleibt, ist geblieben) & auxiliaire \all{sein} \\ \hline
\end{tabularx}
\end{center}

\begin{exoresolu}[Lexique : compléter et employer]
\textbf{Énoncé.}
\begin{enumerate}
\item Complète avec le bon article : \ldots{} Dorf, \ldots{} Stadt, \ldots{} Landflucht, \ldots{} Arbeit, \ldots{} Armut.
\item Complète les phrases avec : \all{umgezogen, Landflucht, Armut, Arbeit, Dörfer}.\\
a) Viele \ldots{} in Kamerun sind fast leer.\\
b) Die \ldots{} ist ein großes Problem für das Land.\\
c) Mein Bruder ist letztes Jahr nach Douala \ldots{}\\
d) Auf dem Land gibt es wenig \ldots{}\\
e) Viele Familien leben in großer \ldots{}
\end{enumerate}
\tcblower
\textbf{Solution détaillée.}
\begin{enumerate}
\item \all{das Dorf} (neutre, pluriel \all{die Dörfer}), \all{die Stadt} (féminin, \all{die Städte}), \all{die Landflucht} (féminin, nom composé : le genre est celui du dernier élément \all{die Flucht}), \all{die Arbeit} (féminin), \all{die Armut} (féminin).
\item a) \all{Dörfer} : pluriel de \all{das Dorf} avec Umlaut. b) \all{Landflucht} : nom composé, sujet de la phrase. c) \all{umgezogen} : Perfekt de \all{umziehen} avec \all{sein} (déménagement = mouvement) ; la particule \all{um-} reste collée au participe : \all{um-ge-zogen}. d) \all{Arbeit} : « peu de travail », sans article après \all{wenig}. e) \all{Armut} : \all{in großer Armut leben} = vivre dans une grande pauvreté (datif féminin : \all{großer}).
\end{enumerate}
\textbf{Conclusion.} Pour le lexique, l'examen exige l'article, le pluriel et la bonne forme du verbe à particule : trois réflexes à vérifier systématiquement.
\end{exoresolu}

\courssub{Méthode 3 : Exprimer la cause avec \all{weil}}

\begin{methode}[Construire une subordonnée causale en \all{weil}]
\begin{enumerate}
\item \textbf{Repère la relation logique} : la cause répond à la question « pourquoi ? » (\all{warum?}).
\item \textbf{Place le verbe à la fin} : après \all{weil}, le verbe conjugué va en \emph{dernière position} de la subordonnée : \all{Er bleibt, weil er hier Arbeit findet.} (et non *\all{weil er findet Arbeit}).
\item \textbf{Choisis la place de la subordonnée} : après la principale (\all{Sie gehen weg, weil es keine Arbeit gibt.}) ou avant — alors la subordonnée occupe la position 1 et le verbe de la principale vient juste après la virgule : \all{Weil es keine Arbeit gibt, gehen viele junge Leute weg.}
\item \textbf{Variantes} : \all{denn} (coordonnant, verbe en position 2 : \all{Er geht weg, denn er findet keine Arbeit.}) ; \all{da} (cause présentée comme connue, souvent en tête de phrase : \all{Da es keine Arbeit gibt, geht er weg.}).
\item \textbf{Au Perfekt dans la subordonnée} : le participe précède l'auxiliaire à la fin : \all{..., weil er keine Arbeit gefunden hat.}
\end{enumerate}
\end{methode}

\begin{center}
\begin{tabularx}{\linewidth}{|X|X|X|}\hline
\rowcolor{popblueL} Connecteur & Place du verbe & Exemple \\ \hline
\all{weil} (parce que) & verbe conjugué à la fin & \all{Er bleibt, weil er Arbeit findet.} \\ \hline
\all{da} (puisque, comme) & verbe conjugué à la fin & \all{Da er keine Arbeit findet, geht er weg.} \\ \hline
\all{denn} (car) & verbe en position 2 & \all{Er geht weg, denn er findet keine Arbeit.} \\ \hline
\end{tabularx}
\end{center}

\begin{exoresolu}[Grammaire : subordonnées en \all{weil}]
\textbf{Énoncé.}
\begin{enumerate}
\item Relie les phrases avec \all{weil} : a) Die jungen Leute verlassen das Dorf. / Es gibt auf dem Land wenig Arbeit. b) Aminatou zieht nach Yaoundé um. / Sie möchte dort studieren.
\item Mets la subordonnée en tête de phrase : \all{Peter bleibt im Dorf, weil er der Familie hilft.}
\item Traduis en allemand : « Beaucoup de jeunes quittent le village parce qu'ils ne trouvent pas de travail. »
\end{enumerate}
\tcblower
\textbf{Solution détaillée.}
\begin{enumerate}
\item a) \all{Die jungen Leute verlassen das Dorf, weil es auf dem Land wenig Arbeit gibt.} Le verbe \all{gibt} passe à la fin ; virgule obligatoire avant \all{weil}. b) \all{Aminatou zieht nach Yaoundé um, weil sie dort studieren möchte.} Attention : verbe à particule \all{umziehen} séparé dans la principale (\all{zieht \ldots{} um}) ; dans la subordonnée, le verbe modal \all{möchte} va à la fin, après l'infinitif \all{studieren}.
\item \all{Weil er der Familie hilft, bleibt Peter im Dorf.} Deux pièges : d'abord, \all{helfen} (aider) se construit avec le \emph{datif} : \all{er hilft der Familie} (et non *\all{die Familie}) ; ensuite, la subordonnée compte comme la position 1, donc le verbe \all{bleibt} reste en position 2 et le sujet \all{Peter} passe après le verbe (inversion).
\item \all{Viele junge Leute verlassen das Dorf, weil sie keine Arbeit finden.} « Ne pas trouver de travail » = \all{keine Arbeit finden} (négation avec \all{kein-} devant un nom sans article) ; \all{finden} à la fin de la subordonnée.
\end{enumerate}
\textbf{Conclusion.} La règle d'or : \all{weil} + sujet + compléments + \textbf{verbe à la fin}. Toute faute de place du verbe est sanctionnée.
\end{exoresolu}

\courssub{Méthode 4 : Exprimer la conséquence (\all{deshalb, darum, deswegen})}

\begin{methode}[Construire une phrase de conséquence]
\begin{enumerate}
\item \textbf{Repère la relation logique} : la conséquence répond à « que se passe-t-il donc ? ». En allemand : \all{deshalb}, \all{darum}, \all{deswegen} (= c'est pourquoi, donc).
\item \textbf{Applique l'inversion} : ces connecteurs occupent la position 1 ; le verbe reste en position 2, le sujet passe après le verbe : \all{Es gibt wenig Arbeit auf dem Land. Deshalb ziehen die jungen Leute in die Stadt.}
\item \textbf{Alternative sans inversion} : place le connecteur après le verbe : \all{Die jungen Leute ziehen deshalb in die Stadt.}
\item \textbf{Ne confonds pas cause et conséquence} : \all{weil} introduit la cause (verbe final), \all{deshalb} introduit le résultat (inversion).
\item \textbf{Combine les deux} pour un texte riche : \all{Weil es wenig Arbeit gibt, verlassen viele das Dorf. Deshalb werden die Dörfer immer leerer.}
\end{enumerate}
\end{methode}

\begin{exoresolu}[Thème : cause et conséquence]
\textbf{Énoncé.} Traduis en allemand :
\begin{enumerate}
\item Il n'y a pas de lycée dans le village. C'est pourquoi les élèves vont en ville.
\item Parce que la pauvreté est grande à la campagne, beaucoup de familles déménagent à Douala.
\item Mon frère ne trouve pas de travail, c'est pourquoi il part.
\end{enumerate}
\tcblower
\textbf{Solution détaillée.}
\begin{enumerate}
\item \all{Im Dorf gibt es kein Gymnasium. Deshalb gehen die Schüler in die Stadt.} « Lycée » = \all{das Gymnasium} (attention, faux ami : pas « gymnase ») ; négation \all{kein} ; inversion après \all{deshalb} : verbe \all{gehen} avant le sujet \all{die Schüler}.
\item \all{Weil die Armut auf dem Land groß ist, ziehen viele Familien nach Douala um.} Subordonnée en tête : le verbe de la principale \all{ziehen} suit immédiatement la virgule ; \all{umziehen} se sépare : \all{ziehen \ldots{} um} ; « à la campagne » = \all{auf dem Land} (et non *\all{im Land}).
\item \all{Mein Bruder findet keine Arbeit, deshalb geht er weg.} « Partir » = \all{weggehen}, verbe à particule : \all{geht \ldots{} weg}. On peut aussi écrire deux phrases : \all{Mein Bruder findet keine Arbeit. Deshalb geht er weg.}
\end{enumerate}
\textbf{Conclusion.} Cause = \all{weil} + verbe final ; conséquence = \all{deshalb} + inversion. Ces deux mécanismes sont au cœur de l'épreuve de thème.
\end{exoresolu}

\courssub{Méthode 5 : Résumer et produire un court paragraphe}

\begin{methode}[Rédiger un résumé ou un court texte sur l'exode rural]
\begin{enumerate}
\item \textbf{Plan en trois temps} : situation (la vie au village) $\rightarrow$ causes du départ (\all{weil}) $\rightarrow$ conséquences (\all{deshalb}) pour le village et la ville.
\item \textbf{Utilise un vocabulaire simple et sûr} : n'écris que ce que tu sais orthographier ; préfère une phrase correcte à une phrase ambitieuse fautive.
\item \textbf{Varie les connecteurs} : \all{zuerst, dann, weil, denn, deshalb, aber, auch}.
\item \textbf{Respecte la structure de phrase} : verbe en position 2 dans la principale, verbe final après \all{weil}, particule séparable renvoyée en fin de phrase.
\item \textbf{Relis-toi} : majuscule à tous les noms, Umlaute (\all{Dörfer, Städte}), accord de l'adjectif, auxiliaire \all{sein} avec \all{umziehen} au Perfekt.
\end{enumerate}
\end{methode}

\begin{exoresolu}[Production écrite guidée]
\textbf{Énoncé.} Rédige un court paragraphe (5 à 6 phrases) en allemand : « Pourquoi les jeunes de ton village partent-ils en ville ? Quelles sont les conséquences ? »
\tcblower
\textbf{Solution détaillée (modèle de production).}\\
\all{In meinem Dorf gibt es wenig Arbeit und keine große Schule. Weil die jungen Leute keine Arbeit finden, ziehen viele nach Yaoundé oder Douala um. Dort arbeiten sie in Büros oder Fabriken, und sie verdienen Geld. Deshalb bleiben im Dorf nur die alten Leute und die Kinder. Die Landflucht ist ein großes Problem, denn das Dorf braucht die jungen Leute. Ich möchte später im Dorf bleiben und hier arbeiten.}\\
« Dans mon village, il y a peu de travail et pas de grande école. Parce que les jeunes ne trouvent pas de travail, beaucoup déménagent à Yaoundé ou Douala. Là-bas, ils travaillent dans des bureaux ou des usines et gagnent de l'argent. C'est pourquoi il ne reste au village que les personnes âgées et les enfants. L'exode rural est un grand problème, car le village a besoin des jeunes. Moi, je voudrais plus tard rester au village et travailler ici. »\\
\textbf{Points de grammaire vérifiés} : verbe final après \all{weil} (\all{finden}) ; inversion après la subordonnée initiale (\all{ziehen viele}) ; particule \all{um} en fin de phrase ; \all{deshalb} + inversion (\all{bleiben}) ; \all{denn} + verbe en position 2 ; majuscules à tous les noms.
\end{exoresolu}

\begin{attention}[Les 5 erreurs qui coûtent des points à l'examen]
\begin{enumerate}
\item \textbf{Verbe mal placé après \all{weil}} : *\all{weil er findet keine Arbeit} est faux. Le verbe conjugué va à la fin : \all{weil er keine Arbeit findet.}
\item \textbf{Oubli de l'inversion après \all{deshalb}} : *\all{Deshalb die jungen Leute ziehen weg} est faux. Écris : \all{Deshalb ziehen die jungen Leute weg.} (verbe en position 2).
\item \textbf{Particule séparable oubliée} : *\all{Er zieht in die Stadt} ne suffit pas si tu veux dire « déménager » : \all{Er zieht in die Stadt um.} Au Perfekt : \all{Er ist umgezogen} (auxiliaire \all{sein}, pas \all{haben}).
\item \textbf{Majuscules et Umlaute} : tous les noms prennent une majuscule (\all{das Dorf, die Stadt, die Arbeit}) ; les pluriels ont souvent un Umlaut : \all{die Dörfer, die Städte}. *\all{die stadte} est une faute.
\item \textbf{Faux amis et calques du français} : \all{das Gymnasium} = le lycée (pas le gymnase) ; « à la campagne » se dit \all{auf dem Land} (pas *\all{im Land} pour ce sens) ; « chercher du travail » = \all{Arbeit suchen}, « trouver » = \all{finden} ; enfin, \all{helfen} (aider) se construit avec le datif : \all{Er hilft der Familie} (et non *\all{die Familie}). Ne traduis jamais mot à mot.
\end{enumerate}
\end{attention}

\newpage

\coursec{Exercices}

\courssub{Je vérifie mes connaissances}

\begin{exercice}[QCM : le texte et le thème]
Entoure la bonne réponse (a, b ou c).
\begin{enumerate}
\item \all{Die Landflucht} bedeutet :
\begin{itemize}
\item[a)] le départ des jeunes vers la ville ;
\item[b)] la construction de nouvelles écoles ;
\item[c)] les vacances à la campagne.
\end{itemize}
\item \all{Viele junge Leute verlassen das Dorf, \ldots{} es keine Arbeit gibt.}
\begin{itemize}
\item[a)] deshalb \quad b) weil \quad c) aber
\end{itemize}
\item Dans une subordonnée avec \all{weil}, le verbe conjugué se trouve :
\begin{itemize}
\item[a)] en deuxième position ;
\item[b)] à la fin de la phrase ;
\item[c)] juste après \all{weil}.
\end{itemize}
\item Le pluriel de \all{das Dorf} est :
\begin{itemize}
\item[a)] die Dorfs \quad b) die Dörfer \quad c) die Dorfen
\end{itemize}
\end{enumerate}
\end{exercice}

\begin{exercice}[Vrai ou faux ?]
Lis le texte de la leçon \all{„Warum verlassen die jungen Leute das Dorf ?“} et réponds par \all{richtig} ou \all{falsch}. Justifie chaque réponse par une citation du texte.
\begin{enumerate}
\item \all{Im Dorf gibt es viele Arbeitsplätze für junge Leute.}
\item \all{Die Stadt bietet bessere Schulen und Krankenhäuser.}
\item \all{Die alten Leute bleiben oft allein im Dorf.}
\item \all{Das Leben in der Stadt ist immer einfach.}
\item \all{Die Regierung will das Leben auf dem Land verbessern.}
\end{enumerate}
\end{exercice}

\begin{exercice}[Vocabulaire : le tableau des mots]
Complète le tableau avec l'article, le pluriel et la traduction française de chaque nom.

\begin{center}
\begin{tabularx}{\linewidth}{|l|l|l|X|}
\hline
\rowcolor{popblueL} Article & Nom & Pluriel & Traduction \\
\hline
 & Dorf &  &  \\
\hline
 & Stadt &  &  \\
\hline
 & Landflucht &  &  \\
\hline
 & Arbeit &  &  \\
\hline
 & Armut &  &  \\
\hline
 & Bauer &  &  \\
\hline
 & Fabrik &  &  \\
\hline
 & Schule &  &  \\
\hline
 & Zukunft &  &  \\
\hline
 & Leben &  &  \\
\hline
\end{tabularx}
\end{center}
\end{exercice}

\begin{exercice}[Grammaire : à compléter]
Complète les phrases avec \all{deshalb}, \all{darum}, \all{so dass} ou \all{denn} selon la structure de la phrase. Attention à la place du verbe !
\begin{enumerate}
\item Es gibt keine Arbeit im Dorf, \ldots{} ziehen viele junge Leute in die Stadt.
\item Die Schulen in der Stadt sind besser, \ldots{} wollen die Familien dort wohnen.
\item Viele Bauern sind arm, \ldots{} sie können ihre Kinder nicht zur Schule schicken.
\item Der Weg zur Stadt ist lang, \ldots{} die jungen Leute müde ankommen.
\item Im Dorf gibt es keinen Arzt, \ldots{} müssen die Kranken weit fahren.
\end{enumerate}
\end{exercice}

\courssub{Je m'entraîne}

\begin{exercice}[Phrases à traduire]
Traduis les phrases suivantes en allemand. Fais attention à la place du verbe et à la majuscule des noms.
\begin{enumerate}
\item Beaucoup de jeunes quittent le village parce qu'il n'y a pas de travail.
\item La ville offre plus d'écoles et d'hôpitaux, c'est pourquoi les familles déménagent.
\item Les paysans sont pauvres, donc ils ne peuvent pas acheter de machines.
\item Mes grands-parents restent au village, mais mon frère habite à Douala.
\item L'exode rural est un grand problème dans notre pays.
\end{enumerate}
\end{exercice}

\begin{exercice}[Relier avec weil]
Relie les deux parties de phrases avec \all{weil} en mettant le verbe conjugué à la fin.
\begin{enumerate}
\item Viele Leute ziehen um. / Es gibt keine Arbeit im Dorf.
\item Die jungen Leute gehen in die Stadt. / Sie suchen eine bessere Zukunft.
\item Die Kinder bleiben zu Hause. / Die Schule ist zu weit weg.
\item Die Familie verkauft das Land. / Sie braucht Geld für das Leben in der Stadt.
\item Die alten Leute sind traurig. / Ihre Kinder kommen selten zurück.
\end{enumerate}
\end{exercice}

\begin{exercice}[Transformer : weil / wegen]
Transforme les phrases : remplace \all{weil} par \all{wegen} et inversement. Attention : \all{wegen} + génitif (langue soutenue) ou datif (langue courante) !
\begin{enumerate}
\item \all{Viele Leute verlassen das Dorf, weil sie arm sind.} $\rightarrow$ \all{wegen \ldots}
\item \all{Wegen des Regens können die Bauern nicht arbeiten.} $\rightarrow$ \all{weil \ldots}
\item \all{Die Familie zieht um, weil der Vater eine neue Arbeit hat.} $\rightarrow$ \all{wegen \ldots}
\item \all{Wegen der schlechten Straßen kommen die Busse nicht ins Dorf.} $\rightarrow$ \all{weil \ldots}
\end{enumerate}
\end{exercice}

\begin{exercice}[Texte à trous]
Complète le texte avec les mots suivants : \all{Arbeit -- Landflucht -- deshalb -- weil -- Stadt -- Dorf -- Zukunft -- Armut}.

\begin{textebox}[Die Landflucht in Kamerun]
In vielen Regionen Kameruns verlassen die jungen Leute das \ldots{}, \ldots{} es dort keine \ldots{} gibt. Die \ldots{} ist ein großes Problem für das Land. In der \ldots{} suchen die jungen Leute eine bessere \ldots{}. Aber das Leben dort ist teuer, und viele Familien leben in \ldots{}. \ldots{} will die Regierung das Leben auf dem Land verbessern.
\end{textebox}
\end{exercice}

\courssub{Je me prépare au Bac}

\begin{exercice}[Compréhension de texte type Bac]
Lis le texte suivant, puis réponds aux questions \textbf{en allemand, par des phrases complètes}.

\begin{textebox}[„Die Stadt ruft -- aber zu welchem Preis?“]
Jedes Jahr verlassen tausende junge Kameruner ihre Dörfer und ziehen nach Yaoundé oder Douala, \all{weil} es auf dem Land wenig Arbeit gibt. Die Gründe sind bekannt: auf dem Land gibt es wenig Arbeit, die Schulen sind weit weg, und die Krankenhäuser sind oft schlecht. In der Stadt hoffen die jungen Leute auf einen Job, eine gute Ausbildung und ein modernes Leben.

Aber die Realität ist oft anders. Viele finden keine feste Arbeit, \all{wenn} sie keine Ausbildung haben, und müssen als Händler oder Motorradfahrer -- die berühmten „Bendskins“ -- Geld verdienen. Die Mieten in der Stadt sind hoch, und viele Familien wohnen in kleinen Häusern ohne Wasser und Strom.

Im Dorf bleiben die Alten allein zurück. Niemand hilft ihnen mehr bei der Arbeit auf den Feldern, und die Ernte wird kleiner. Manche Dörfer sind fast leer.

Experten sagen: Man muss das Landleben attraktiver machen -- mit Straßen, Schulen, Strom und Arbeit für die Jugend. Dann kommen vielleicht auch junge Leute zurück aufs Land.
\end{textebox}

\textbf{Questions de compréhension :}
\begin{enumerate}
\item Warum verlassen die jungen Leute ihre Dörfer? (Nenne zwei Gründe.)
\item Was machen viele junge Leute in der Stadt, wenn sie keine feste Arbeit finden?
\item Welche Probleme gibt es im Dorf, wenn die Jungen weg sind?
\item Was schlagen die Experten vor?
\end{enumerate}

\textbf{Questions de langue :}
\begin{enumerate}
\item Trouve dans le texte deux subordonnées avec \all{weil} ou \all{wenn} et recopie-les.
\item Mets à l'infinitif : \all{verlassen -- hoffen -- bleiben -- wird}.
\item Donne le contraire de : \all{klein -- hoch -- alt}.
\end{enumerate}
\end{exercice}

\begin{exercice}[Thème et version]
\textbf{A. Version.} Traduis en français le passage suivant :

\begin{textebox}[À traduire]
Die Landflucht ist ein großes Problem in vielen afrikanischen Ländern. Weil es auf dem Land keine Arbeit gibt, ziehen die jungen Leute in die Städte. Deshalb bleiben die alten Leute allein, und die Dörfer werden leer.
\end{textebox}

\textbf{B. Thème.} Traduis en allemand :
\begin{enumerate}
\item Beaucoup de jeunes gens quittent le village parce qu'ils cherchent du travail.
\item Il n'y a pas d'hôpital dans notre village, c'est pourquoi les malades vont en ville.
\item La pauvreté est la cause de l'exode rural.
\item Si le gouvernement construit des écoles et des routes, les jeunes resteront au village.
\end{enumerate}
\end{exercice}

\begin{exercice}[Rédaction type Bac]
\textbf{Sujet :} \all{Die Landflucht in deinem Land : Ursachen und Folgen.}

Rédige un court paragraphe de \textbf{8 à 10 lignes} en allemand. Tu présenteras :
\begin{itemize}
\item deux ou trois causes de l'exode rural dans ton pays ;
\item deux ou trois conséquences pour les villages et pour les villes ;
\item ta proposition personnelle pour améliorer la situation.
\end{itemize}

\textbf{Contraintes :} utilise obligatoirement les connecteurs \all{weil}, \all{aber}, \all{deshalb} et \all{zum Schluss}. Soigne la place du verbe, la majuscule des noms et les articles.
\end{exercice}



\newpage

\coursec{Corrigés des exercices}

\begin{corrige}[Exercice 1 — QCM : le texte et le thème]
\begin{enumerate}
\item \textbf{a)} \all{Die Landflucht} = l'exode rural, le départ des jeunes de la campagne vers la ville (mot composé : \all{das Land} + \all{die Flucht}, littéralement « la fuite de la campagne »).
\item \textbf{b)} \all{weil} introduit une cause : \all{Viele junge Leute verlassen das Dorf, weil es keine Arbeit gibt.} Attention : \all{deshalb} et \all{aber} exigeraient le verbe en 2\cle{e} position (\all{deshalb gibt es ...}), pas à la fin.
\item \textbf{b)} Règle fondamentale : dans toute subordonnée introduite par \all{weil}, \all{dass}, \all{wenn}, \all{als}, \all{ob}..., le verbe conjugué est rejeté \textbf{à la fin} de la phrase.
\item \textbf{b)} \all{das Dorf} $\rightarrow$ \all{die Dörfer} : pluriel en \all{-er} avec Umlaut (\all{o} $\rightarrow$ \all{ö}), comme \all{das Haus} $\rightarrow$ \all{die Häuser}.
\end{enumerate}
\textbf{Points de vigilance :} ne jamais écrire \all{die Dorfs} (pluriel anglais !) ; mémoriser les pluriels avec Umlaut dès le niveau A1 : \all{das Dorf -- die Dörfer}, \all{die Stadt -- die Städte}, \all{das Haus -- die Häuser}.
\end{corrige}

\begin{corrige}[Exercice 2 — Vrai ou faux ?]
\begin{enumerate}
\item \all{falsch} — Justification : \all{„Es gibt kaum Arbeit für junge Leute.“} Le texte dit le contraire : il n'y a presque pas de travail au village.
\item \all{richtig} — Justification : \all{„In der Stadt gibt es bessere Schulen und Krankenhäuser.“}
\item \all{richtig} — Justification : \all{„Die alten Leute bleiben allein zurück.“}
\item \all{falsch} — Justification : \all{„Das Leben in der Stadt ist teuer und schwer.“} Le mot \all{immer} (« toujours ») rend l'affirmation fausse : la vie en ville est difficile, mais le texte ne dit pas qu'elle est \emph{toujours} facile ou agréable.
\item \all{richtig} — Justification : \all{„Die Regierung plant, das Leben auf dem Land zu verbessern.“}
\end{enumerate}
\textbf{Méthode :} pour justifier, recopier toujours une courte citation exacte du texte entre guillemets allemands „…“ ; ne pas paraphraser. Commencer la réponse par \all{richtig} ou \all{falsch}, puis donner la citation.
\end{corrige}

\begin{corrige}[Exercice 3 — Vocabulaire : le tableau des mots]
\begin{center}
\begin{tabularx}{\linewidth}{|l|l|l|X|}
\hline
\rowcolor{popblueL} Article & Nom & Pluriel & Traduction \\
\hline
\all{das} & \all{Dorf} & \all{die Dörfer} & le village \\
\hline
\all{die} & \all{Stadt} & \all{die Städte} & la ville \\
\hline
\all{die} & \all{Landflucht} & \all{--} (sans pluriel) & l'exode rural \\
\hline
\all{die} & \all{Arbeit} & \all{die Arbeiten} & le travail \\
\hline
\all{die} & \all{Armut} & \all{--} (sans pluriel) & la pauvreté \\
\hline
\all{der} & \all{Bauer} & \all{die Bauern} & le paysan, l'agriculteur \\
\hline
\all{die} & \all{Fabrik} & \all{die Fabriken} & l'usine \\
\hline
\all{die} & \all{Schule} & \all{die Schulen} & l'école \\
\hline
\all{die} & \all{Zukunft} & \all{--} (singulier usuel) & l'avenir \\
\hline
\all{das} & \all{Leben} & \all{die Leben} & la vie \\
\hline
\end{tabularx}
\end{center}
\textbf{Points de vigilance :}
\begin{itemize}
\item \all{Landflucht}, \all{Armut} et \all{Zukunft} sont des noms abstraits utilisés presque toujours au singulier (\emph{Singulariatantum}).
\item \all{der Bauer} suit la \cle{déclinaison faible} : \all{den Bauern, dem Bauern, des Bauern} ; pluriel \all{die Bauern}.
\item Apprendre chaque nom \textbf{avec son article et son pluriel} : c'est la clé des cas et des accords.
\end{itemize}
\end{corrige}

\begin{corrige}[Exercice 4 — Grammaire : à compléter]
\begin{enumerate}
\item \all{Es gibt keine Arbeit im Dorf, deshalb ziehen viele junge Leute in die Stadt.} (\all{deshalb}/\all{darum} = adverbe conjonctif $\rightarrow$ verbe \all{ziehen} en 2\cle{e} position, avant le sujet.)
\item \all{Die Schulen in der Stadt sind besser, deshalb wollen die Familien dort wohnen.} (Même structure : inversion sujet-verbe après \all{deshalb}.)
\item \all{Viele Bauern sind arm, denn sie können ihre Kinder nicht zur Schule schicken.} (\all{denn} = conjonction de coordination $\rightarrow$ verbe en 2\cle{e} position, ordre normal.)
\item \all{Der Weg zur Stadt ist lang, so dass die jungen Leute müde ankommen.} (\all{so dass} = subordonnée de conséquence $\rightarrow$ verbe à la fin ; le verbe à particule \all{an|kommen} se recolle : \all{ankommen}.)
\item \all{Im Dorf gibt es keinen Arzt, deshalb müssen die Kranken weit fahren.} (Adverbe conjonctif $\rightarrow$ inversion ; \all{keinen Arzt} : accusatif masculin de la négation.)
\end{enumerate}
\textbf{Récapitulatif des trois structures :}
\begin{center}
\begin{tabularx}{\linewidth}{|l|X|X|}
\hline
\rowcolor{popblueL} Connecteur & Type & Place du verbe \\
\hline
\all{denn} & coordination & 2\cle{e} position (ordre normal) \\
\hline
\all{deshalb}, \all{darum} & adverbe conjonctif & 2\cle{e} position (inversion) \\
\hline
\all{weil}, \all{so dass} & subordination & à la fin \\
\hline
\end{tabularx}
\end{center}
\end{corrige}

\begin{corrige}[Exercice 5 — Phrases à traduire]
\begin{enumerate}
\item \all{Viele junge Leute verlassen das Dorf, weil es dort keine Arbeit gibt.} (« parce que » = \all{weil} $\rightarrow$ verbe \all{gibt} à la fin.)
\item \all{Die Stadt bietet mehr Schulen und Krankenhäuser, deshalb ziehen die Familien um.} (« déménager » = \all{um|ziehen}, particule \all{um} à la fin de la principale.)
\item \all{Die Bauern sind arm, deshalb können sie keine Maschinen kaufen.} (Négation d'un nom pluriel indéfini : \all{keine Maschinen}.)
\item \all{Meine Großeltern bleiben im Dorf, aber mein Bruder wohnt in Douala.} (\all{im} = \all{in dem} ; \all{aber} ne change pas l'ordre des mots.)
\item \all{Die Landflucht ist ein großes Problem in unserem Land.} (\all{in} + datif pour le lieu : \all{unserem Land} ; adjectif neutre après \all{ein} : \all{großes}.)
\end{enumerate}
\textbf{Points de vigilance :} majuscule à tous les noms (\all{Dorf, Schulen, Krankenhäuser, Maschinen, Problem}) ; ne pas oublier le \all{ß} de \all{Großeltern, großes} ; « le problème » est neutre : \all{das Problem}.
\end{corrige}

\begin{corrige}[Exercice 6 — Relier avec weil]
\begin{enumerate}
\item \all{Viele Leute ziehen um, weil es im Dorf keine Arbeit gibt.}
\item \all{Die jungen Leute gehen in die Stadt, weil sie eine bessere Zukunft suchen.}
\item \all{Die Kinder bleiben zu Hause, weil die Schule zu weit weg ist.}
\item \all{Die Familie verkauft das Land, weil sie Geld für das Leben in der Stadt braucht.}
\item \all{Die alten Leute sind traurig, weil ihre Kinder selten zurückkommen.}
\end{enumerate}
\textbf{Points de vigilance :}
\begin{itemize}
\item Virgule obligatoire devant \all{weil}.
\item Verbes à particule séparable : dans la subordonnée, la particule se \textbf{recolle} au verbe à la fin : \all{weg ist}, \all{zurückkommen}, \all{umzieht}.
\item Le pronom \all{sie} reprend le sujet de la principale : vérifier l'accord du verbe (\all{sie suchen} au pluriel ; \all{sie braucht} pour \all{die Familie}, nom singulier).
\end{itemize}
\end{corrige}

\begin{corrige}[Exercice 7 — Transformer : weil / wegen]
\begin{enumerate}
\item \all{Viele Leute verlassen das Dorf, weil sie arm sind.} $\rightarrow$ \all{Wegen ihrer Armut verlassen viele Leute das Dorf.} (Le pronom \all{sie} devient le possessif \all{ihrer} au génitif féminin ; \all{wegen} en tête $\rightarrow$ verbe \all{verlassen} en 2\cle{e} position.)
\item \all{Wegen des Regens können die Bauern nicht arbeiten.} $\rightarrow$ \all{Weil es regnet, können die Bauern nicht arbeiten.} (Le nom \all{der Regen} redevient une proposition avec le verbe impersonnel \all{regnen} ; subordonnée en tête $\rightarrow$ verbe de la principale \all{können} juste après la virgule.)
\item \all{Die Familie zieht um, weil der Vater eine neue Arbeit hat.} $\rightarrow$ \all{Wegen der neuen Arbeit des Vaters zieht die Familie um.} (Double génitif : \all{der neuen Arbeit} + \all{des Vaters}.)
\item \all{Wegen der schlechten Straßen kommen die Busse nicht ins Dorf.} $\rightarrow$ \all{Weil die Straßen schlecht sind, kommen die Busse nicht ins Dorf.} (L'adjectif épithète \all{schlechten} devient attribut : \all{schlecht sind}, verbe à la fin.)
\end{enumerate}
\textbf{Méthode :} \all{weil} + phrase (sujet + verbe conjugué à la fin) ; \all{wegen} + nom (génitif soutenu : \all{wegen des Regens} ; datif courant à l'oral : \all{wegen dem Regen}). Transformer revient à passer du nom au verbe, ou inversement.
\end{corrige}

\begin{corrige}[Exercice 8 — Texte à trous]
\textbf{Texte complété :}\par
\all{In vielen Regionen Kameruns verlassen die jungen Leute das Dorf, weil es dort keine Arbeit gibt. Die Landflucht ist ein großes Problem für das Land. In der Stadt suchen die jungen Leute eine bessere Zukunft. Aber das Leben dort ist teuer, und viele Familien leben in Armut. Deshalb will die Regierung das Leben auf dem Land verbessern.}\par
« Dans de nombreuses régions du Cameroun, les jeunes quittent le village parce qu'il n'y a pas de travail là-bas. L'exode rural est un grand problème pour le pays. En ville, les jeunes cherchent un avenir meilleur. Mais la vie là-bas est chère, et beaucoup de familles vivent dans la pauvreté. C'est pourquoi le gouvernement veut améliorer la vie à la campagne. »

\textbf{Justifications :}
\begin{itemize}
\item \all{weil} : subordonnée de cause, verbe \all{gibt} à la fin.
\item \all{keine Arbeit} : négation d'un nom féminin indéfini $\rightarrow$ \all{keine}.
\item \all{in der Stadt} : datif féminin (lieu, pas de mouvement).
\item \all{in Armut} : expression figée sans article.
\item \all{Deshalb} en tête de phrase $\rightarrow$ inversion : \all{will die Regierung} (verbe avant sujet).
\end{itemize}
\end{corrige}

\begin{corrige}[Exercice 9 — Compréhension de texte type Bac]
\textbf{Barème indicatif :} compréhension 8 pts (2 pts par question) ; questions de langue 6 pts (2 pts par question). Toute réponse en phrase complète et correcte ; une citation exacte vaut la justification.

\textbf{Questions de compréhension :}
\begin{enumerate}
\item \all{Die jungen Leute verlassen ihre Dörfer, weil es auf dem Land wenig Arbeit gibt und die Schulen weit weg sind.} (On accepte aussi : \all{die Krankenhäuser sind oft schlecht}.)
\item \all{Wenn sie keine feste Arbeit finden, arbeiten sie als Händler oder Motorradfahrer (Bendskins), um Geld zu verdienen.}
\item \all{Die alten Leute bleiben allein zurück, niemand hilft ihnen bei der Arbeit auf den Feldern, die Ernte wird kleiner und manche Dörfer sind fast leer.}
\item \all{Die Experten schlagen vor, das Landleben attraktiver zu machen — mit Straßen, Schulen, Strom und Arbeit für die Jugend.}
\end{enumerate}

\textbf{Questions de langue :}
\begin{enumerate}
\item Deux subordonnées du texte : \all{„..., weil es auf dem Land wenig Arbeit gibt.“} et \all{„..., wenn sie keine Ausbildung haben, ...“}. (Vérifier : verbe conjugué \all{gibt} / \all{haben} bien à la fin.)
\item Infinitifs : \all{verlassen} $\rightarrow$ \all{verlassen} ; \all{hoffen} $\rightarrow$ \all{hoffen} ; \all{bleiben} $\rightarrow$ \all{bleiben} ; \all{wird} $\rightarrow$ \all{werden} (3\cle{e} pers. sing. du présent de \all{werden}).
\item Contraires : \all{klein} $\leftrightarrow$ \all{groß} ; \all{hoch} $\leftrightarrow$ \all{niedrig} (ou \all{tief}) ; \all{alt} $\leftrightarrow$ \all{jung} (pour les personnes) ou \all{neu} (pour les choses).
\end{enumerate}

\textbf{Points de vigilance :} répondre par des \textbf{phrases complètes} (sujet + verbe en 2\cle{e} position), pas par des mots isolés ; reprendre le vocabulaire du texte sans recopier de longs passages ; attention aux majuscules des noms dans les citations : \all{Arbeit, Schulen, Ausbildung}.
\end{corrige}

\begin{corrige}[Exercice 10 — Thème et version]
\textbf{Barème indicatif :} version 5 pts ; thème 5 pts (1 pt par phrase : vocabulaire 0,5 + grammaire/ordre des mots 0,5).

\textbf{A. Version :}

« L'exode rural est un grand problème dans de nombreux pays africains. Parce qu'il n'y a pas de travail à la campagne, les jeunes partent s'installer dans les villes. C'est pourquoi les personnes âgées restent seules, et les villages se vident. »

Correspondances à noter : \all{die Landflucht} = l'exode rural ; \all{auf dem Land} = à la campagne ; \all{in die Städte ziehen} = partir s'installer en ville ; \all{leer werden} = se vider.

\textbf{B. Thème :}
\begin{enumerate}
\item \all{Viele junge Leute verlassen das Dorf, weil sie Arbeit suchen.} (\all{weil} $\rightarrow$ verbe \all{suchen} à la fin.)
\item \all{Es gibt kein Krankenhaus in unserem Dorf, deshalb gehen die Kranken in die Stadt.} (\all{kein} neutre ; \all{in unserem Dorf} datif ; \all{die Kranken} = « les malades », adjectif substantivé avec majuscule.)
\item \all{Die Armut ist die Ursache der Landflucht.} (\all{die Ursache} + génitif : \all{der Landflucht}.)
\item \all{Wenn die Regierung Schulen und Straßen baut, werden die jungen Leute im Dorf bleiben.} (Subordonnée \all{wenn} en tête $\rightarrow$ verbe \all{baut} à la fin ; principale au \cle{Futur I} : \all{werden ... bleiben}, infinitif à la fin.)
\end{enumerate}

\textbf{Points de vigilance :} ne pas traduire « c'est pourquoi » par \all{das ist warum} ! La bonne solution est \all{deshalb} (ou \all{darum}) avec inversion. « Resteront » annonce un futur : en allemand, le présent est souvent accepté (\all{bleiben die jungen Leute im Dorf}), mais le \all{Futur I} est plus sûr.
\end{corrige}

\begin{corrige}[Exercice 11 — Rédaction type Bac]
\textbf{Barème indicatif (sur 10) :} respect du sujet et du plan (causes / conséquences / proposition) : 3 pts ; emploi des quatre connecteurs imposés (\all{weil}, \all{aber}, \all{deshalb}, \all{zum Schluss}) : 2 pts ; correction grammaticale (place du verbe, cas, articles, majuscules) : 3 pts ; richesse du vocabulaire du thème : 2 pts.

\textbf{Proposition de rédaction modèle :}\par
\all{In Kamerun ist die Landflucht ein großes Problem. Viele junge Leute verlassen ihre Dörfer, weil es dort keine Arbeit, keine guten Schulen und keine Krankenhäuser gibt. Auch die Straßen sind oft schlecht. Deshalb ziehen sie in die großen Städte wie Yaoundé oder Douala und hoffen auf ein besseres Leben. Aber die Realität ist schwer : die Mieten sind hoch, und viele junge Leute arbeiten als Bendskin-Fahrer oder Händler. Für die Dörfer ist die Landflucht auch negativ : die alten Leute bleiben allein, und niemand arbeitet mehr auf den Feldern. Zum Schluss meine ich : der Staat muss Schulen, Straßen und Fabriken auf dem Land bauen. Dann bleiben die jungen Leute im Dorf, und die Landflucht wird kleiner.}\par
« Au Cameroun, l'exode rural est un grand problème. Beaucoup de jeunes quittent leurs villages parce qu'il n'y a là-bas ni travail, ni bonnes écoles, ni hôpitaux. Les routes aussi sont souvent mauvaises. C'est pourquoi ils partent dans les grandes villes comme Yaoundé ou Douala et espèrent une vie meilleure. Mais la réalité est dure : les loyers sont élevés, et beaucoup de jeunes travaillent comme conducteurs de bendskin ou commerçants. Pour les villages, l'exode rural est aussi négatif : les personnes âgées restent seules, et personne ne travaille plus dans les champs. Pour conclure, je pense que l'État doit construire des écoles, des routes et des usines à la campagne. Alors les jeunes resteront au village, et l'exode rural diminuera. »

\textbf{Points de vigilance :}
\begin{itemize}
\item Vérifier chaque subordonnée : après \all{weil}, verbe à la fin (\all{... gibt}).
\item Après \all{deshalb} et \all{zum Schluss} en tête de phrase : inversion verbe-sujet (\all{ziehen sie}, \all{meine ich}).
\item Majuscules : \all{Landflucht, Dörfer, Arbeit, Schulen, Krankenhäuser, Straßen, Realität, Mieten, Staat}.
\item Éviter les faux amis : « réaliser » n'est pas \all{realisieren} ici ; « le problème » = \all{das Problem} (neutre !) ; « rester » = \all{bleiben}, jamais \all{resten}.
\end{itemize}
\end{corrige}

\newpage

\coursec{Fiche bilan}

\begin{popfigure}
\begin{tikzpicture}[mindmap, grow cyclic, every node/.style={concept, align=center, font=\small},
  root concept/.append style={concept color=popblue, font=\bfseries\small, minimum size=2.6cm},
  level 1 concept/.append style={sibling angle=60, level distance=3.6cm, minimum size=2.2cm}]
\node[root concept, align=center]{Die Landflucht\\ \emph{l'exode rural}}
  child[concept color=poporange]{ node[align=center]{Le texte\\ „Warum verlassen\\ die jungen Leute\\ das Dorf?“} }
  child[concept color=popgreen]{ node[align=center]{Lexique\\ das Dorf, die Stadt\\ die Arbeit, die Armut\\ umziehen} }
  child[concept color=popgold]{ node[align=center]{La cause\\ weil + verbe final\\ wegen + nom} }
  child[concept color=poppurple]{ node[align=center]{La conséquence\\ deshalb, darum\\ so dass} }
  child[concept color=poppink]{ node[align=center]{Subordonnée\\ en \emph{weil}\\ verbe conjugué\\ à la fin} }
  child[concept color=popblue!70]{ node[align=center]{Méthode\\ lire, repérer\\ résumer, produire} };
\end{tikzpicture}
\legende{Carte mentale de la leçon : l'exode rural (die Landflucht)}
\end{popfigure}

\begin{bilan}
\textbf{1. Lexique clé à connaître par cœur}
\begin{itemize}
  \item \all{das Dorf, die Dörfer} : le village ; \all{die Stadt, die Städte} : la ville
  \item \all{die Landflucht} (n. f., sans pluriel courant) : l'exode rural
  \item \all{die Arbeit, -en} : le travail ; \all{der Arbeitsplatz, die Arbeitsplätze} : l'emploi ; \all{arbeitslos} : au chômage
  \item \all{die Armut} : la pauvreté ; \all{arm} : pauvre ; \all{reich} : riche
  \item \all{umziehen (zieht um, zog um, ist umgezogen)} : déménager ; \all{verlassen (verlässt, verließ, hat verlassen)} : quitter
  \item \all{der Bauer, -n} : le paysan ; \all{das Land} : la campagne ; \all{auf dem Land leben} : vivre à la campagne
  \item \all{die Schule, -n} : l'école ; \all{das Krankenhaus, die Krankenhäuser} : l'hôpital ; \all{die Möglichkeit, -en} : la possibilité
\end{itemize}

\textbf{2. Grammaire : la cause et la conséquence}

\begin{center}
\begin{tabularx}{\linewidth}{|l|X|X|}
\hline
\rowcolor{popblueL} Outil & Règle & Exemple \\
\hline
\all{weil} (parce que) & subordonnée : verbe conjugué à la \textbf{fin} & \all{Er zieht um, weil er keine Arbeit findet.} \\
\hline
\all{denn} (car) & coordination : verbe en 2\textsuperscript{e} position & \all{Er zieht um, denn er findet keine Arbeit.} \\
\hline
\all{wegen} + génitif/datif & devant un \textbf{nom}, pas de verbe & \all{Wegen der Armut verlassen viele das Dorf.} \\
\hline
\all{deshalb / darum} (c'est pourquoi) & conséquence : verbe en 2\textsuperscript{e} position, le sujet après le verbe & \all{Es gibt keine Arbeit. Deshalb ziehen die Jungen weg.} \\
\hline
\all{so dass} (si bien que) & subordonnée de conséquence, verbe à la fin & \all{Die Armut ist groß, so dass viele wegziehen.} \\
\hline
\end{tabularx}
\end{center}

\textbf{3. Règles d'or}
\begin{itemize}
  \item Après \all{weil} : sujet + … + \textbf{verbe conjugué en dernière position}. Jamais : \emph{weil er findet keine Arbeit}.
  \item Si la subordonnée en \all{weil} ouvre la phrase, le verbe de la principale suit immédiatement la virgule : \all{Weil es keine Arbeit gibt, ziehen viele junge Leute in die Stadt.}
  \item \all{wegen} + nom (souvent datif à l'oral : \all{wegen dem Wetter}) ; \all{weil} + phrase verbale.
  \item Majuscule à tous les noms : \all{das Dorf, die Armut, die Landflucht}.
\end{itemize}

\textbf{4. Méthode express : comprendre et résumer un texte}
\begin{enumerate}
  \item Lire le titre et la source : anticiper le thème (\all{Landflucht} = exode rural).
  \item Première lecture : qui ? où ? quel problème ?
  \item Deuxième lecture : souligner les causes (\all{weil, wegen}) et les conséquences (\all{deshalb}).
  \item Répondre aux questions avec des phrases complètes, en reprenant les mots du texte.
  \item Résumer en 3 à 5 phrases : problème, causes, conséquences, solutions possibles.
\end{enumerate}
\end{bilan}

\begin{aretenir}[Checklist avant le Bac]
\begin{itemize}
  \item $\square$ Je connais l'article et le pluriel de \all{das Dorf, die Stadt, die Arbeit, die Schule, der Arbeitsplatz}.
  \item $\square$ Je sais définir \all{die Landflucht} en une phrase simple.
  \item $\square$ Je conjugue sans erreur \all{umziehen} et \all{verlassen} au Präsens et au Perfekt (auxiliaire \all{sein} pour \all{umziehen}).
  \item $\square$ Je place le verbe conjugué à la fin dans une subordonnée en \all{weil}.
  \item $\square$ Je distingue \all{weil} (+ phrase) et \all{wegen} (+ nom).
  \item $\square$ J'emploie \all{deshalb} avec le verbe en deuxième position et le sujet après le verbe.
  \item $\square$ Je sais construire une phrase avec \all{so dass} pour exprimer la conséquence.
  \item $\square$ Je mets une majuscule à tous les noms allemands et j'utilise les guillemets „…“.
  \item $\square$ Je sais citer au moins deux causes de l'exode rural (chômage, manque d'écoles et d'hôpitaux) et deux conséquences (vieillissement des villages, surpopulation des villes).
  \item $\square$ Je sais rédiger un court paragraphe de 5 phrases sur l'exode rural au Cameroun avec au moins une subordonnée en \all{weil}.
\end{itemize}
\end{aretenir}

\findecours

\end{document}
